\documentclass[11pt]{article}

\usepackage[T1]{fontenc}
\usepackage[utf8]{inputenc}
\usepackage[english]{babel}
\usepackage[margin=1in]{geometry}
\usepackage{lmodern}
\usepackage{microtype}
\usepackage{mathtools}
\usepackage{amsmath,amssymb,amsthm}
\usepackage{array}
\usepackage{booktabs}
\usepackage{tabularx}
\usepackage{float}
\usepackage{graphicx}
\usepackage{pdflscape}
\usepackage{caption}
\usepackage{enumitem}
\usepackage{fancyhdr}
\usepackage[numbers]{natbib}
\usepackage{doi}
\usepackage{xurl}
\usepackage{hyperref}
\usepackage{cleveref}
\usepackage{orcidlink}
\usepackage{titling}



\newcommand{\FOne}{\textsc{Foundations~I}}
\newcommand{\FTwo}{\textsc{Foundations~II}}
\newcommand{\FThree}{\textsc{Foundations~III}}
\newcommand{\PaperAdequacy}{\textsc{Paper~1 (Adequacy Residuals)}}
\newcommand{\PaperSDTC}{\textsc{Paper~2 (SDTC)}}
\newcommand{\PaperAOR}{\textsc{Paper~3 (AOR)}}
\newcommand{\PaperNeedles}{\textsc{Paper~4 (Needle Killer)}}
\newcommand{\PaperNS}{\textsc{Paper~5 (NS)}}
\newcommand{\PaperRH}{\textsc{Paper~6 (RH)}}
\newcommand{\PaperPvNP}{\textsc{Paper~7 (PvNP)}}
\newcommand{\PaperCSL}{\textsc{Paper~8 (CSL)}}



\newcommand{\Xequiv}{$X$-equivalent}
\newcommand{\nonXequiv}{non-$X$-equivalent}
\newcommand{\CRETE}{\textsc{CRE-TE}}
\newcommand{\CREBF}{\textsc{CRE-BF}}
\newcommand{\CRECTMT}{\textsc{CRE-CTMT}}
\newcommand{\NativeClosure}{native-closure outcome}
\newcommand{\XstrengthObl}{$X$-strength obligation}

\newcommand{\thmCarrierDichotomy}{Theorem~\ref{thm:carrier-dichotomy} (Carrier Dichotomy)}
\newcommand{\corBridgeImpossibility}{Corollary~\ref{cor:bridge-impossibility} (Bridge Impossibility)}
\newcommand{\thmCoverage}{Conditional schema~\ref{thm:coverage} (Coverage)}
\newcommand{\thmTFF}{Theorem~\ref{thm:type-finite-foreclosure} (Type-Finite Foreclosure)}
\newcommand{\thmAFvFour}{Conditional schema~\ref{thm:attack-foreclosure-v4} (Attack Foreclosure v4)}
\newcommand{\thmCTMT}{Theorem~\ref{thm:ctmt-recursion} (CTMT Recursion)}

\newcommand{\fim}{framework-internal move}
\newcommand{\carrierPivot}{carrier pivot}
\newcommand{\bridgeImport}{bridge import}
\newcommand{\cascadeInternal}{cascade-internal computation}
\newcommand{\TCD}{typed-context discipline}
\newcommand{\CTMTTower}[1]{CTMT tower of depth~$#1$}
\newcommand{\TypedBridgeCarrier}{carrier-side typed bridge}

\newcommand{\VDTSO}{V-Differential trace-state-only column}
\newcommand{\VDWCE}{V-Differential witness-content-exposing column}

\newcommand{\SevenStage}{seven-stage recognition-mode landing template}
\newcommand{\StageOne}{Stage~1 (closure construction)}
\newcommand{\StageTwo}{Stage~2 (translation theorem)}
\newcommand{\StageThree}{Stage~3 (applicability audit)}
\newcommand{\StageFour}{Stage~4 (smuggle audit)}
\newcommand{\StageFive}{Stage~5 (anti-tautology hardening)}
\newcommand{\StageSix}{Stage~6 (three-option derivation sweep)}
\newcommand{\StageSeven}{Stage~7 (recognition closure)}

\newcommand{\OptFP}{framework-primitive option}
\newcommand{\OptSAU}{SAU non-descent option}
\newcommand{\OptVDE}{V-Differential elevation option}
\newcommand{\SigPattern}{$\mathrm{(ii)}/\mathrm{(iii)}/\mathrm{(ii)}$}

\newcommand{\thmTemplate}{Conditional schema~\ref{thm:template-applicability} (Recognition-Mode Template Applicability)}
\newcommand{\thmVerdict}{Conditional schema~\ref{thm:verdict-signature} (Cross-Track Verdict Signature)}
\newcommand{\thmCCT}{Conditional schema~\ref{thm:closure-content-thesis} (Closure-Content Thesis)}
\newcommand{\thmAFvFive}{Conditional schema~\ref{thm:attack-foreclosure-v5} (Attack Foreclosure v5)}

\newcommand{\RML}{recognition-mode landing}
\newcommand{\RMmode}{recognition mode}
\newcommand{\DMmode}{derivation mode}
\newcommand{\PPGrade}{paper-proposal-grade}
\newcommand{\TwoInstance}{two-instance verified, corpus-pending}
\newcommand{\FwInternal}{framework-internal}
\newcommand{\FwExternal}{framework-external}

\newcommand{\Cvobs}{\mathcal{V}_{\text{obs}}}
\newcommand{\CMI}{\mathcal{M}_{\text{I}}}
\newcommand{\CRE}{\mathcal{R}_{\text{E}}}
\newcommand{\CCh}{\mathcal{C}_{\text{h}}}
\newcommand{\CQamp}{\mathcal{Q}_{\text{amp}}}
\newcommand{\CIalg}{\mathcal{I}_{\text{alg}}}
\newcommand{\CRreal}{\mathcal{R}_{\text{real}}}
\newcommand{\COCstar}{\mathcal{O}_{C^*}}
\newcommand{\CFexistsSO}{\mathcal{F}_{\exists \mathrm{SO}}}
\newcommand{\CSHO}{\mathcal{S}_{\mathrm{HO}}}
\newcommand{\CPPH}{\mathcal{P}_{\mathrm{PH}}}
\newcommand{\CAVNP}{\mathcal{A}_{\mathrm{VNP}}}
\newcommand{\CMmot}{\mathcal{M}_{\text{mot}}}

\newcommand{\fourAxis}{four-axis calibration profile}
\newcommand{\TypedHole}{typed coverage hole}
\newcommand{\xiNCone}{\xi_{\text{NC}^{1}\text{-coverage}}}
\newcommand{\xiAdaptive}{\xi_{\text{adaptive-trace}}}
\newcommand{\FType}{framework-type}
\newcommand{\TDS}{typed-deposit species}
\newcommand{\TypedBridgeCalibration}{calibration-side typed bridge}

\newcommand{\thmCFS}{Theorem~\ref{thm:calibration-family-structure} (Calibration-Family Structure)}
\newcommand{\thmRebadge}{Theorem~\ref{thm:rebadge-discipline} (Rebadge Discipline)}
\newcommand{\thmBaT}{Theorem~\ref{thm:bridges-as-theorems} (Bridges-as-Theorems)}
\newcommand{\thmBTD}{Theorem~\ref{thm:bridge-type-direction} (Bridge Type-Direction)}
\newcommand{\thmCFG}{Theorem~\ref{thm:coverage-family-geometry} (Coverage-Family Geometry)}
\newcommand{\thmFTSC}{Theorem~\ref{thm:framework-type-self-char} (Framework-Type Self-Characterization)}
\newcommand{\thmGates}{Catalog~\ref{cat:named-functorial-gates} (Named Functorial Gates)}

\newcommand{\GateFunctionDegree}{\Xi_{\text{function-degree-to-NS-certificate-degree}}}
\newcommand{\GatePlantedStatistical}{\Xi_{\text{planted-statistical-to-NP-hard-language-lift}}}
\newcommand{\GateUniformTime}{\Xi_{\text{uniform-time-hierarchy-realization-functor}}}

\newcommand{\ObligationRef}[1]{(Appendix~D, Obligation \texttt{#1})}
\newcommand{\LeanDecl}[1]{\texttt{\detokenize{#1}}}
 
\newtheorem{theorem}{Theorem}[section]
\newtheorem{lemma}[theorem]{Lemma}
\newtheorem{proposition}[theorem]{Proposition}
\newtheorem{corollary}[theorem]{Corollary}
\newtheorem{schema}[theorem]{Conditional schema}
\Crefname{schema}{Conditional schema}{Conditional schemas}
\crefname{schema}{conditional schema}{conditional schemas}
\theoremstyle{definition}
\newtheorem{definition}[theorem]{Definition}
\newtheorem{remark}[theorem]{Remark}
\floatstyle{ruled}
\newfloat{algorithm}{tbp}{loa}
\floatname{algorithm}{Algorithm}

\urlstyle{same}
\captionsetup{
  font=small,
  labelfont=bf,
  labelsep=period
}
\captionsetup[table]{skip=6pt,position=top}
\captionsetup[figure]{skip=6pt}
\captionsetup[algorithm]{skip=6pt,position=top}
\crefname{algorithm}{algorithm}{algorithms}
\Crefname{algorithm}{Algorithm}{Algorithms}
\crefname{theorem}{theorem}{theorems}
\Crefname{theorem}{Theorem}{Theorems}
\crefname{lemma}{lemma}{lemmas}
\Crefname{lemma}{Lemma}{Lemmas}
\crefname{proposition}{proposition}{propositions}
\Crefname{proposition}{Proposition}{Propositions}
\crefname{corollary}{corollary}{corollaries}
\Crefname{corollary}{Corollary}{Corollaries}
\crefname{definition}{definition}{definitions}
\Crefname{definition}{Definition}{Definitions}
\crefname{remark}{remark}{remarks}
\Crefname{remark}{Remark}{Remarks}
\setlist[enumerate]{leftmargin=*, itemsep=2pt, topsep=4pt}
\setlength{\droptitle}{-3.2em}
\pretitle{\begin{center}\LARGE\bfseries}
\posttitle{\par\end{center}\vspace{0.5em}}
\preauthor{\begin{center}\normalsize}
\postauthor{\par\end{center}\vspace{-0.4em}}
\predate{\begin{center}\small}
\postdate{\par\end{center}\vspace{0.6em}}
\setlength{\emergencystretch}{2em}
\clubpenalty=10000
\widowpenalty=10000
\displaywidowpenalty=10000
\interfootnotelinepenalty=10000
\allowdisplaybreaks[2]
\fancypagestyle{firstpage}{\fancyhf{}
  \fancyfoot[C]{\scriptsize Automorph Inc., Wilmington, DE, USA \quad \texttt{ioannis@automorph.io} \quad \textcopyright\ 2026 \quad CC-BY 4.0}
  \renewcommand{\headrulewidth}{0pt}
  \renewcommand{\footrulewidth}{0pt}
}

\title{One Meta-Theory, Three Clay-Problem Closures}
\author{Ioannis Tsiokos\,\orcidlink{0009-0009-7659-5964}}
\date{\today}
\hypersetup{
  pdftitle={One Meta-Theory, Three Clay-Problem Closures},
  pdfauthor={Ioannis Tsiokos},
  hidelinks
}

\begin{document}

\maketitle
\thispagestyle{firstpage}

\begin{abstract}
Six Birds Theory has produced three landed Clay-class conditional
closures: Navier--Stokes regularity, the Riemann Hypothesis, and
P~vs~NP, each under the SB closure assumption of Foundations~I.  This
paper does not re-prove those substrate results.  It extracts their
common meta-theoretic form for a reader who treats Six Birds as a
typed framework for organizing attempted closures of open conjectures.
The account has two complementary halves.  The negative half,
Foreclosure, studies a typed primary carrier, meaning an attempted
proof construction whose type signature targets a designated
conjecture, and proves that admissible carriers fall into a finite
state space: non-target-equivalent, target-equivalent tautological,
bridge-failed, or matrix-element-terminal.  The positive half,
Recognition-Mode Landing, gives a seven-stage route for
trace-state-only substrates: instead of deriving the load-bearing
source from framework primitives, it imports that source as named
recognition under closure formation.  The v5 refinement states the
joint thesis: derivation routes remain subject to foreclosure, while
recognition-mode landings are structurally distinct peer routes under
the same closure assumption.  A third part, the Calibration Family,
provides the typed bookkeeping that makes this comparison uniform
across substrates: calibrations, bridges-as-theorems, directed bridge
types, coverage geometry, framework-type mismatch, and named
functorial gates.  Layer-instantiation examples appear throughout: the
three substrate validation rows (NS, RH, and PvNP), a thirteen-entry
calibration seed catalog of typed detectors, two named coverage holes,
and three named functorial gates from the PvNP track.  The formal
catalog contains seventeen rows: three partial structural derivations,
eight primary theorem-grade rows, and six primary obligation rows.  The body gives the mathematical proofs or
obligation analyses, while Lean serves as the mechanization companion
and audit trail.  The nine honest obligation axioms are
recorded with opacity guards, so their assumptions are explicit rather
than silently instantiated; the full validator chain cross-checks the
inventory, manifest, and trust base.  The scope is deliberately
limited: BSD and Hodge are out of scope except as future work or
prediction, and no substrate closure is strengthened here to a
standard-ZFC unconditional claim.
\end{abstract}

\medskip
\noindent\textbf{Keywords.} Six Birds Theory; carrier dichotomy;
recognition-mode landing; closure-content thesis; calibration family;
meta-theoretic foreclosure.

\section{Introduction}\label{sec:intro}

Three recent papers in this series treat Navier--Stokes regularity,
the Riemann Hypothesis, and P versus NP under a common conditional
assumption.  This paper does not add a fourth problem and does not
reargue those three case studies.  It asks why the three arguments
have the same broad shape.

The three case studies are the empirical starting point.  The
conditional closure of Navier--Stokes regularity landed by \PaperNS\
under the SB closure assumption of \FOne\
\cite{Tsiokos2026NavierStokes,Tsiokos2026EmergenceCalculus}; the
conditional closure of the Riemann Hypothesis landed by \PaperRH\
under the SB closure assumption of \FOne\ \cite{Tsiokos2026RH}; and
the conditional closure of P versus NP landed by \PaperPvNP\ under
the SB closure assumption of \FOne\ \cite{Tsiokos2026PvNP}.  The
official Clay Mathematics Institute statements of these three
problems are
\cite{Fefferman2000NavierStokes,Bombieri2000Riemann,Cook2000PvsNP};
the corresponding volume is
\cite{CarlsonJaffeWiles2006MillenniumPrize}.  Here the
SB closure assumption is the closure-formation working axiom of
\FOne: it is the common hypothesis under which the substrate landings
are stated, not a claim that those landings have been upgraded to
standard-ZFC-unconditional theorems.  The Six Birds framework, as the
term is used in this paper, is the formal program developed in the
Foundations monographs for studying structural closure of mathematical
problems; we use \FOne\ for closure formation, \FTwo\ for the
admissibility background, and \FThree\ for the typed audit discipline
\cite{Tsiokos2026EmergenceCalculus,Tsiokos2026FoundationsII,Tsiokos2026FoundationsIII}.

The meta-question is therefore not whether the three substrate
papers can be compressed into one direct proof.  It is whether their
shared shape has a stable mathematical explanation.  The answer
developed here has two halves.  The negative half says that, once a
target problem has been placed inside the framework's typed audit
discipline, internal attempts to close the target fall into a finite
set of carrier states.  A carrier is the chosen mathematical
presentation through which the framework tries to make the target
visible; a typed carrier is one whose target, admissible imports, and
residual obstruction are part of the data being audited.  The positive
half says that some closures are not obtained by forcing a derivation
through those internal carriers; instead, they land by naming the
structural source that already supplies the closure content under the
SB closure assumption.  The two halves are not competitors.  They are
the two sides of the same account: foreclosure explains why
framework-internal derivations run out, while recognition-mode landing
explains how a formed closure can nevertheless land at theorem grade
under the same foundational assumption.

This distinction matters because the three substrate landings do not
look like routine variations on one proof template.  The conditional
Navier--Stokes closure is organized around analytic carrier
discipline.  The conditional RH closure uses a spectral and trace
state structure.  The conditional P versus NP closure lives in a
complexity-theoretic landscape where presentation, average-case
structure, and nonuniformity all matter.  If the common shape were
only rhetorical, a meta-theory would add little.  What the present
paper argues is that the common shape is structural: each case has a
target, a native carrier, a residual obstruction, a distinction
between internal and external content, and an audit that prevents
unlicensed transport between them.

The first half is organized around three named spine claims.  The
Carrier Dichotomy Theorem says, at the level of typed carriers, that
a carrier either remains non-equivalent to the target or enters one
of three carrier-residual-equivalent modes.  The Coverage Theorem is
the exhaustiveness assertion behind that classification: every
admissible carrier aimed at the target is accounted for by the four
typed states.  Attack Foreclosure v4 then packages the consequence
for derivation routes: carrier pivoting, bridge import, and
cascade-internal computation do not provide a free internal route to
the target.  These previews are deliberately informal; the full
statements, hypotheses, and audit disclosures appear in
\S\ref{sec:foreclosure}.

The carrier language should be read in ordinary mathematical terms
first.  A failed attempt to close a target can fail because the
presentation was aimed at the wrong proposition; because it was aimed
at the right proposition but became trapped in a self-audit failure;
because a bridge to the target would itself require the target-level
content; or because the computation reaches a matrix-element terminus
whose decision is not supplied by the carrier.  The framework's
vocabulary records these alternatives as typed states rather than as
informal diagnoses.  The claim is not that there are only four ways a
mathematician can be stuck.  It is narrower: after the framework has
fixed the target, the carrier, and the admissible imports, the audit
sees only a finite door-space of typed outcomes.

That finiteness is the reason the negative half can be a foreclosure
statement rather than a case-by-case complaint.  Carrier pivoting is
the move from one presentation of the target to another.  Bridge
import is the move that tries to carry residual closure across a
transport from one carrier to the target.  Cascade-internal
computation is the move that keeps reducing inside the same carrier,
hoping the residual becomes structurally or numerically smaller.  The
foreclosure claim is that these moves do not become free merely
because they are internal to the framework.  Each must either land in
the finite carrier analysis or expose the external mathematical
content it needs.  This is why the paper treats foreclosure as an
audit theorem, not as a pessimistic slogan about proof search.

The second half begins from a different observation.  The RH and
P versus NP landings use a recognition-mode pattern rather than a
pure derivation-mode pattern.  A recognition-mode landing is a
structured landing in which the framework records a named
mathematical source of closure content instead of claiming that this
content has been generated internally.  A recognition source, in this
paper's usage, is such a named structural fact about the substrate
setting; when the CSL and SDTC sources are discussed later, they are
described as imported as named recognition per \FOne\ closure
formation, never as derived from framework primitives.  A formed
closure is the corresponding closure object: the target, the closure
content, and the admissible source are already organized as one
structural setting rather than assembled after the fact.

The Recognition-Mode Template Applicability Theorem is the spine
claim that the seven-stage recognition-mode landing template applies
to formed Six Birds closures whose substrates sit in the relevant
trace-state-only column.  In plain terms, a trace-state-only column is
a substrate class where the closure content appears as a trace-state
fact rather than as an internally descending derivation.  The template
itself is a disciplined seven-step audit: construct the closure
setting, translate it to the target, check applicability, rule out
smuggling, harden against tautology, sweep the available derivation
options, and land by recognition.  Attack Foreclosure v5 refines, and
does not replace, Attack Foreclosure v4: derivation routes remain
foreclosed, while recognition-mode landings are structurally distinct
and can be theorem-grade under the SB closure assumption.

This is the main safeguard against a common misunderstanding.  The
recognition-mode path is not a weaker substitute for a derivation-mode
proof, and it is not a license to call any useful outside theorem a
closure source.  It is a typed landing pattern.  The source must be
named; it must be admitted by the closure-formation discipline; the
attempt to derive it internally must be separated from the act of
recognizing it; and the resulting closure is conditional under the SB
closure assumption.  The positive half of the paper is therefore just
as restrictive as the negative half.  It says when recognition is a
valid landing mode, not that recognition removes the need for
mathematical content.

The present paper is not a survey of the analytic, arithmetic, or
complexity-theoretic details of the three substrate papers.  It
extracts from them the common structural question: when the framework
has produced a formed setting, which parts of the work are genuinely
internal, which parts are named external mathematics, and how does the
audit prevent the two from being conflated?  The resulting answer is
modest in one respect and ambitious in another.  It is modest because
it does not claim new unconditional results about the Clay problems.
It is ambitious because it claims that the framework's closure
behavior has a reusable meta-theoretic form.

Part III supplies the apparatus that makes this answer uniform across
substrates.  A calibration is a typed way of measuring an obstruction:
it chooses a detector class and the native carrier type in which that
detector lives.  The calibration family records the non-descending
translations used by the framework as a family rather than as one
universal instrument.  Different obstruction shapes call for
different detector classes and different native carrier types:
low-degree polynomial detectors, information-theoretic detectors,
stable or local algorithmic detectors, topological detectors, and
others.  The calibration apparatus explains why bridges between such
calibrations are typed and directional, why some gaps are named
mathematical gaps rather than free changes of notation, and why the
framework's natural reach on P versus NP is structurally different
from the standard Clay target type.  Without this third part, the
foreclosure and recognition-mode halves would remain a collection of
case-specific observations.  With it, they become one account of how
the framework organizes closure attempts across different
mathematical substrates.

The calibration family also explains why this paper does not try to
state a single universal obstruction measure.  A detector that sees a
low-degree polynomial shadow is not the same mathematical instrument
as one that sees an information-theoretic obstruction, a stable-local
algorithmic obstruction, or a topological obstruction.  Bridges
between these instruments can be basis changes in favorable cases,
but in general they are substantive theorems with direction.  This is
particularly visible on the P versus NP track, where the framework's
native type is nonuniform, average-case, scoped, and package-based,
while the standard Clay target is uniform, worst-case,
unconditional, and language-based.  The gap is not a matter of
notation; it is part of the mathematical bookkeeping that the
meta-theory must expose.

The payoff of this bookkeeping is methodological rather than
computational.  A mathematician reading only the substrate papers sees
three different mathematical terrains.  A mathematician reading this
paper should see why the same audit questions recur in all three:
what is the target, what is the carrier, what is the residual, what
counts as an accepted import, and where does target-strength content
enter?  These questions are ordinary questions about proof structure,
but the framework makes them explicit enough to support a finite
classification.  That explicitness is useful even when the reader does
not adopt the framework's vocabulary as their own.  It isolates the
point at which a proposed internal derivation stops being internal.

This also explains the precise sense in which the paper speaks of
foreclosure.  The paper does not claim that the space of mathematical
ideas is finite.  It claims that, after the framework fixes the type
of the target and the admissible role of a carrier, the audit sees a
finite set of state-types through which a derivation-mode attempt can
pass.
The word ``foreclosure'' is therefore not a statement about human
ingenuity.  It is a statement about typed bookkeeping: certain moves
cannot be counted as free progress because they reintroduce the target
content, miss the target, or terminate at content not supplied by the
carrier.  This is the level at which a meta-theorem can speak across
Navier--Stokes, RH, and P versus NP without pretending that the three
problems have the same mathematics.

The phrase ``recognition-mode landing'' carries the complementary
point.  Some target closures, on this account, are not best understood
as the endpoint of a descending internal calculation.  They are best
understood as the point where a formed closure setting exposes the
load-bearing source and the framework names that source under the
closure-formation discipline.  This is why the paper is careful about
the verb ``import.''  A recognition source is imported as named
recognition; it is not produced inside the framework by a hidden
derivation.  The import is still mathematical content, and the audit
is responsible for saying where it came from and why it is admissible
under the SB closure assumption.

These two phrases also delimit what the paper can and cannot do.  The
negative half cannot prove that no outside mathematics will ever close
a target; it classifies the framework-internal moves available once
the target and carrier discipline are fixed.  The positive half cannot
turn a recognition source into an internally derived theorem; it
records the source as named closure content under \FOne.  The
calibration half cannot remove every gap between framework-native
objects and standard problem statements; it names those gaps as typed
transport problems.  The three halves together give a coherent audit
of the existing conditional landings, not an unconditional solution
manual for future Clay-class problems.

This positioning fixes the relation to the rest of the series.
\PaperNS, \PaperRH, and \PaperPvNP\ are the substrate papers: they
contain the problem-specific conditional landings.  \PaperCSL\
\cite{Tsiokos2026CSL} and \PaperSDTC\ \cite{Tsiokos2026SDTC} supply
recognition-source material used by the recognition-mode track.
\PaperAdequacy\ \cite{Tsiokos2026XiCompanion} supplies the adequacy
residual language used by the carrier analysis, while
\PaperNeedles\ \cite{Tsiokos2026NeedleKiller} and the Foundations
volumes provide the surrounding audit and admissibility vocabulary.
This paper is downstream from those works.  It does not replace their
arguments; it gives a meta-theoretic account of the common shape that
their conditional closures exhibit.

The dependence on the substrate papers is therefore intentionally
asymmetric.  When this paper invokes a substrate landing, it does so
only as a conditional closure already landed in that substrate paper
under the SB closure assumption of \FOne.  When it invokes the
Foundations volumes, it uses them for the shared closure-formation,
admissibility, and audit vocabulary.  When it invokes the apparatus
papers, it uses them for residuals, recognition-source supply, and
calibration evidence.  None of these citations is meant to smuggle a
new Clay-class proof into the introduction.  They establish the record
whose common structure is being analyzed.

The paper is written for a reader who may know none of this internal
language.  For that reason the body proceeds in two registers.  Each
new framework-native term is first given an ordinary mathematical
paraphrase: a carrier is a presentation, a residual is the obstruction
left by that presentation, a bridge is a typed transport, a
recognition source is named closure content, and a calibration is an
obstruction detector together with its native carrier type.  After
that first translation, the shorter framework term is used because the
term does real work.  The reader should not need to know the internal
history of the Six Birds program to follow the argument, but the
reader should be able to see why the vocabulary is more than
decoration.

The formal catalog behind the paper has seventeen primary rows,
grouped into the three parts just described: six Foreclosure rows,
four Recognition-Mode rows, and seven Calibration-Family rows.  Some
rows are fully structural theorems, some are structural derivations
from explicitly stated same-paper premises, and some are candidate
meta-claims recorded for later audit.  The introduction only names
the spine.  Later sections give the statements, the proof modes, and
the exact caveats.

For a reader new to this series, the most important practical point is
that the paper separates three tasks that are often conflated.  First,
one must identify the target-relative carrier and the residual
obstruction.  Second, one must determine whether attempted internal
moves actually lower that obstruction or merely move the same
target-strength obligation to another place.  Third, if the closure
lands by recognition, one must name the imported structural content
and say exactly why it is accepted under the closure-formation
discipline.  The paper's terminology is heavy only because these
distinctions are load-bearing.  Whenever possible, the body translates
the framework-native vocabulary back into the ordinary mathematical
role it is playing.

The intended reading order is as follows.  Section~\ref{sec:apparatus}
sets up the minimal vocabulary: carriers, residuals, typed audit
discipline, recognition sources, and calibration language.  Section
\ref{sec:foreclosure} is the easiest doorway into the paper: it gives
Part I, the finite carrier-state analysis and the foreclosure account
for derivation-mode routes.  Section~\ref{sec:recognition} gives
Part II, the recognition-mode landing template and the v5 refinement
that places recognition-mode landings peer with derivation-mode
landings under the SB closure assumption.  Section
\ref{sec:validation} checks the account against the three landed
substrates.  Section~\ref{sec:calibration} gives Part III, the
calibration-family apparatus that makes the account portable across
substrates.  The last three sections record the cross-paper audit,
the predictions and open problems, and the nonclaims and honest
caveats.

The scope is intentionally narrow.  BSD \cite{Wiles2000BSD} and Hodge
\cite{Deligne2000Hodge} are outside the body claims of this paper:
BSD has no landing in the source corpus, and Hodge is treated only as
a future prediction in the appropriate section.  Nothing in this paper asserts a standard-ZFC-unconditional
proof of Navier--Stokes regularity, the Riemann Hypothesis, P versus
NP, BSD, Hodge, or any other Clay-class problem.
 \section{Apparatus}\label{sec:apparatus}

\subsection{Closure Formation}\label{sec:apparatus:closure-formation}

A closure formation is the background act of organizing an open
problem, its admissible source material, and the claimed closure event
inside one formal setting.  In this paper the framework-native name is
\FOne\ closure formation, and we use it only as a black-box import from
\FOne\ \cite{Tsiokos2026EmergenceCalculus}.  A formed Six Birds
closure is therefore a formal closure of an open problem in which the
closure event is itself treated as a structural fact about the
problem's mathematical setting, rather than as an informal label
attached after the fact to a proof.  The SB closure assumption is the
\FOne\ closure-formation working axiom under which all substrate
landings cited in this paper are stated.  It is not a
standard-ZFC-unconditional claim, and it is not re-proved here.

In the notation used later, a substrate \(S\) that has undergone this
closure-formation act is recorded by the structural predicate
\[
  \operatorname{FormedClosure}(S).
\]
The SB closure assumption is the ambient hypothesis under which this
predicate is used for the cited substrate landings.  Since \FOne\ fixes
that assumption as framework background, this paper does not introduce a
new independent symbol for it; it treats it as the standing
closure-formation hypothesis whenever a substrate landing is cited.

The same background also fixes the role of accepted imports.  An
accepted import is external mathematical content that the closure
formation permits the carrier to use under its stated type.  This does
not mean that every useful theorem may be silently inserted into the
argument.  The import must be named, must play a typed role in the
closure formation, and must be distinguished from content produced by
the carrier itself.  This is why the later recognition-mode sections
use the phrase ``imported as named recognition.''  It is a closure-
formation phrase, not a derivation claim.

The carrier arguments in Part I use the same convention in the form of
an accepted-imports base
\[
  B:\mathrm{Prop}\to\mathrm{Prop},\qquad
  h_{\mathrm{sound}}:\forall Y:\mathrm{Prop},\; B(Y)\Rightarrow Y.
\]
Thus an accepted import is not a loose citation.  It is a proposition
selected by \(B\), together with the soundness hypothesis that such a
selection contributes genuine accepted content.  The Carrier Dichotomy
uses this hypothesis explicitly in \Cref{sec:foreclosure}.

This convention is what keeps the conditional-closure language stable
across the paper.  When a substrate paper lands a target under the SB
closure assumption, the landing is being evaluated inside this
closure-formation discipline.  The meta-math paper then asks how the
landing is structured, not whether the SB closure assumption can be
removed.  In particular, a formed closure can be theorem-grade
relative to the SB closure assumption while still depending on named
source material.  Later sections use this distinction constantly:
they separate the carrier's native work from the accepted imports that
closure formation licenses.

\subsection{The BirdInt Judgment}\label{sec:apparatus:birdint}

A typed interaction judgment is a way of saying not merely that two
formal objects interact, but that the interaction belongs to a
specified structural scenario.  The framework-native form used here is
the BirdInt judgment of \FThree\
\cite{Tsiokos2026FoundationsIII}: the Foundations III judgment form
that classifies six-bird interactions into structurally typed cases.
The particular background fact needed later is the Same-Level
Self-Audit Failure theorem, Theorem 24 of \FThree.  We use that result
only as imported audit vocabulary: it is the structural reason why an
attempt that collapses into auditing itself at the same target level is
classified as a target-equivalent failure mode.  Section
\ref{sec:foreclosure} uses this black-box fact when explaining the
target-equivalent state in the Carrier Dichotomy.

At the level of notation, a BirdInt judgment records a typed interaction
of a source, target, and audit scenario:
\[
  \operatorname{BirdInt}(A,B;\sigma).
\]
Theorem~24 of \FThree\ says, in the form needed here, that a judgment
whose audit attempts to certify its own target-level claim at the same
level collapses to same-level self-audit failure:
\[
  \operatorname{BirdInt}(A,A;\mathrm{same\mbox{-}level\ audit})
  \Longrightarrow
  \operatorname{SelfAuditFail}(A).
\]
Part I abbreviates this imported conclusion as the self-audit-failure
condition appearing in the \CRETE{} state.

The admissibility background for these judgments is supplied by
\FTwo\ \cite{Tsiokos2026FoundationsII} together with \FThree.  The
point for the present paper is simply that the audit judgment is typed:
one cannot replace a source, target, or level by an informally similar
one and still count the interaction as the same judgment.

For a reader not using the Six Birds vocabulary, the BirdInt role can
be understood as a typed diagnostic layer.  It tells the paper which
kind of failure or interaction has occurred, rather than merely
recording that an attempted proof did not close.  In particular, a
self-audit failure at the same level is not treated as a harmless
restatement.  It says that the attempted closure has fallen back onto
the target-level content it was supposed to supply.  That diagnostic
role is what later makes the target-equivalent branch of the carrier
classification mathematically meaningful.

The paper therefore treats BirdInt as a source of audit judgments, not
as an additional proof engine.  Its contribution here is classificatory:
it supplies a typed language in which one can say that a carrier has
entered a same-level self-audit failure, or that a proposed interaction
belongs to the wrong level, without converting that diagnosis into a
new theorem about the target.  This is why the Carrier Dichotomy can
use the Same-Level Self-Audit Failure theorem without re-opening the
Foundations III proof.

\subsection{Adequacy Residuals}\label{sec:apparatus:adequacy-residual}

An adequacy residual is the obstruction left after a proposed proof
carrier has covered what it can cover.  The framework-native notation
comes from \PaperAdequacy\ \cite{Tsiokos2026XiCompanion}: for a
carrier \(C\), domain \(D\), and level \(L\), the quantity
\(\Xi_C(D\mid L)\) measures the gap between what \(C\) covers in
\(D\) and what is actually load-bearing for the claim at level \(L\).
This paper uses the residual as a black box.  It does not re-derive
the residual calculus of \PaperAdequacy.

The residual should not be read as a numerical error term unless a
particular substrate supplies such an interpretation.  In this paper
it is a structural obstruction: the part of the target-relevant load
that remains after the carrier's native resources have been credited.
Thus a residual can be closed natively, can be shown to track the
target itself, or can expose a bridge or matrix-element terminus.  The
Carrier Dichotomy in Section~\ref{sec:foreclosure} uses exactly this
residual role.  It does not need the internal construction of
\(\Xi_C(D\mid L)\); it needs the fact that \PaperAdequacy\ supplies a
typed residual notion that can be compared with the target and with
the carrier's accepted imports.

The phrase ``residual closure'' later means that this obstruction has
been closed relative to the carrier's type and accepted-imports base.
It does not by itself say that the target has been proved in ordinary
unconditional form.  The meta-theoretic question is what kind of
closure has occurred.  A residual closure can be native to the carrier,
can be target-equivalent in a way that exposes circularity, can require
a bridge to the target, or can terminate at a matrix-element family
whose decision is outside the carrier.  Those alternatives are the raw
material for the Part I carrier-state analysis.

\begin{definition}[Carrier residual]\label{def:carrier-residual}
For the purposes of this paper, the carrier's adequacy residual
\(\Xi_C(D\mid L)\) is the Paper~1 obstruction-residual measure of a
typed carrier \(C\) against its accepted-imports base, meaning the
external mathematical material licensed for that carrier by closure
formation: it records the remaining target-relevant gap after the
carrier's native coverage in domain \(D\) is accounted for at level
\(L\).
\end{definition}

Thus a residual closure is written schematically as
\[
  \Xi_C(D\mid L)=0,
\]
but the equality is interpreted inside the carrier's typed context and
accepted-imports base.  When Part I compares a residual closure with a
target \(X\), the comparison is not a numerical comparison; it is the
typed assertion that the residual-closure proposition is classically
equivalent to the target proposition under the accepted imports.

The carrier analysis also imports two pieces of black-box apparatus
from \PaperNeedles\ \cite{Tsiokos2026NeedleKiller}.  First, the
layer-dissolving membrane endpoint is the endpoint reached when a
dissolving probe breaks a carrier's native structure rather than
preserving it.  Second, the all-six taxonomy is the Paper~4 taxonomy
of the six endpoint types that carrier analyses inherit.  In ordinary
terms, Paper~4 supplies the endpoint language for saying whether a
probe stays native to the carrier or dissolves the layer in which the
carrier was supposed to operate.  Later sections use this vocabulary
to discuss carrier states and cascade endpoints, but the taxonomy
itself is not reconstructed here.

This black-box use of \PaperNeedles\ is deliberately conservative.
The present paper does not need to classify every possible probe in
the Paper~4 sense.  It needs only the endpoint distinction that
carrier analyses inherit: a probe may remain native to the carrier's
own type, or it may dissolve the layer and thereby stop being evidence
that the original carrier closed the target.  This is the apparatus
behind the later language of native closure, bridge failure, and
cascade terminality.

The layer-dissolving viewpoint is also the bridge to the recognition
mode.  If a probe dissolves the layer, the framework should not pretend
that the original carrier has performed a native closure.  It must
either record the dissolved endpoint as external content or explain how
the endpoint is accepted by the closure formation.  This is one of the
places where the negative and positive halves of the paper meet: the
same apparatus that blocks an internal derivation route can identify
the source that a recognition-mode landing must name.

\subsection{Typed-Context Discipline}\label{sec:apparatus:typed-context}

The typed-context discipline is the ambient typing discipline of Parts
I--III: every claim, carrier, bridge, calibration, and obligation
carries an explicit type signature, and no informal substitution is
permitted between types.  This discipline is inherited from the
\FTwo/\FThree\ admissibility and audit setting
\cite{Tsiokos2026FoundationsII,Tsiokos2026FoundationsIII}.  In Part I
it governs the four carrier-state classifications and the bridge
composition claims.  In Part II it governs the distinction between a
derivation route and a recognition-mode landing.  In Part III it
governs calibration bridges, including the fact that a bridge from one
typed source to another need not be reversible.  The purpose of the
discipline is to prevent an argument from silently moving between
different sources, targets, or carrier regimes merely because the
informal prose makes them sound related.

The judgment form used in the body is
\[
  C\models\text{\TCD},
\]
meaning that the carrier \(C\) is being audited in a context whose
source, target, carrier type, and admissible imports have all been fixed.
Together with the aimed-at relation \(C\leadsto X\), this is the typed
prefix for the carrier-state predicates of Part I:
\[
  C\leadsto X,\qquad C\models\text{\TCD}.
\]

Two consequences are used throughout the body.  First, a carrier aimed
at one target cannot be treated as a carrier aimed at another target
without a typed bridge or an explicit re-aiming step.  Second, a
bridge has direction: a transport from one regime to another does not
by itself supply the reverse transport.  These are ordinary typing
rules, but they do substantial work in the meta-theory.  They are what
let the paper distinguish a genuine basis change from an imported
theorem, and what prevent a recognition source from being rebranded as
an internally derived primitive.

This discipline is especially important because the paper moves across
different mathematical substrates.  Analytic regularity, trace-state
arithmetic, and complexity-theoretic hardness do not live in the same
native carrier type.  A statement that is meaningful in one setting may
have no typed reverse in another.  The meta-theory therefore refuses
to identify objects by analogy alone.  It asks for the source type, the
target type, and the direction of the transport.  Parts I--III apply
that same rule to carriers, bridges, calibrations, and named transport
gates.

\subsection{Conditional Schemas in the Body}\label{sec:apparatus:opacity-guards}

Several meta-mathematical claims in Parts~I, II, and~III remain open and
are stated in this paper as conditional schemas rather than as proved
theorems. Each conditional schema is presented in its own
\texttt{schema} environment in the body, with its conditions stated
explicitly. The formalization-side register of these schemas is
collected in Appendix~\ref{app:formalization}.
 \section{Part I --- Foreclosure}\label{sec:foreclosure}

\subsection{Setup}\label{sec:foreclosure:setup}

Part I gives the negative half of the paper: it explains why the
framework's internal routes do not freely close a target conjecture
once the target, the carrier, and the admissible imports have been
typed.  The entry point is deliberately elementary.  We begin with the
object being audited, then list the four ways the audit can classify
it.  Full theorem statements are deferred to the later subsections of
Part I.

A typed primary carrier aimed at \(X\) is an attempted proof
construction whose type signature targets a designated open conjecture
\(X\).  The carrier is not just a proof idea; it includes the target it
is aimed at, the residual obstruction it leaves, and the imports it is
allowed to use under the \TCD.  This is the level at which the
foreclosure account operates: not on all possible mathematical
thoughts about \(X\), but on carriers whose role in the framework has
already been made explicit.

Mathematically, the carrier used in this section is a typed object
\[
  C=(E_C,L_C,D_C,A_C,\rho_C)
\]
with carrier dimension data, audit matrix \(A_C\), native-probe matrix
\(L_C\), dissolving readout \(D_C\), and audit record \(\rho_C\).  Its
carrier residual is the Paper~1 residual
\[
  R_C := \Xi_C(D_C\mid L_C).
\]
The predicates \(C\leadsto X\) and \(C\models\text{\TCD}\) record,
respectively, that the carrier is aimed at \(X\) and that the
typed-context discipline applies to the carrier audit.

There are four carrier states in the Part I audit.  A \nonXequiv\
carrier is not classically routed to \(X\) at all: its closure, if it
has one, points to the wrong conjectural content rather than to the
target.  An \Xequiv\ carrier, by contrast, is one whose residual
closure tracks \(X\) itself.  Such a carrier can still fail in three
typed ways.  The \CRETE\ state is the target-equivalent fallback: the
carrier has become a tautology relative to \(X\) itself, a case whose
audit background is the Same-Level Self-Audit Failure theorem of
\FThree\ \cite{Tsiokos2026FoundationsIII}.  The \CREBF\ state is the
bridge-failed dead end: the carrier needs a typed bridge to the target
whose composition is unavailable or structurally incompatible.  The
\CRECTMT\ state is the cascade-terminal matrix-element case: the
carrier's internal reduction reaches a matrix-element terminus whose
decision is not supplied by the carrier.

In predicate form, with \(R_C=0\) denoting closure of the residual,
the four states used throughout Part~I are:
\[
\begin{aligned}
  \mbox{\Xequiv}_X(C)
    &\Longleftrightarrow (R_C=0 \leftrightarrow X),\\
  \mbox{\nonXequiv}_X(C)
    &\Longleftrightarrow
      \exists X'.\, (R_C=0 \leftrightarrow X')\wedge X'\neq X,\\
  \mbox{\CRETE}_X(C)
    &\Longleftrightarrow
      \mbox{\Xequiv}_X(C)\wedge
      \operatorname{SelfAuditFail}(C),\\
  \mbox{\CREBF}_X(C)
    &\Longleftrightarrow
      \operatorname{NoBridge}(C,X)\vee
      \operatorname{BridgeMismatch}(C,X),\\
  \mbox{\CRECTMT}_X(C)
    &\Longleftrightarrow
      \operatorname{MRed}(C)\wedge
      \neg\operatorname{IRDec}(C)\wedge
      \neg\operatorname{MDEqTarget}(C,X).
\end{aligned}
\]
Here \(\operatorname{SelfAuditFail}(C)\) abbreviates the
Foundations~III audit alternatives \(C\)'s same-level audit being
outside scope or undefined-circular.  The symbols
\(\operatorname{NoBridge}\), \(\operatorname{BridgeMismatch}\),
\(\operatorname{MRed}\), \(\operatorname{IRDec}\), and
\(\operatorname{MDEqTarget}\) abbreviate, respectively, the
target-relative bridge failures, matrix-family reduction,
inherited-records decision predicate, and target-equivalence test for
the terminal matrix family.

The operational moves reviewed later in Part I are equally few.  A
\fim\ is one of three internal moves licensed by the framework:
\carrierPivot, replacing the present carrier by another carrier aimed
at \(X\); \bridgeImport, trying to compose a carrier with a transport
bridge to \(X\); or \cascadeInternal, continuing the computation within
the fixed carrier in order to lower its residual.  The foreclosure
claim is not that these moves are syntactically impossible.  It is
that, under the typed audit, each move either lands back in the four
carrier-state classification or exposes named external content.
The move taxonomy is the three-constructor type
\[
  \mbox{\fim} ::= \mbox{\carrierPivot}
      \mid \mbox{\bridgeImport}
      \mid \mbox{\cascadeInternal}.
\]

The remaining subsections of Part I now make this setup precise.  The
Carrier Dichotomy Theorem gives the four-state classification.  The
Bridge Impossibility Corollary records what happens when a
non-\(X\)-equivalent carrier is bridged into an \(X\)-equivalent
composite.  The
Coverage Theorem supplies the exhaustiveness claim, and Type-Finite
Foreclosure converts that coverage into a finite door-space of carrier
outcomes.  Attack Foreclosure v4 packages the three framework-internal
moves into the negative-half conclusion.  CTMT Recursion explains why
the cascade-terminal state can occur at nested depths rather than only
as a one-step endpoint.  In reading order: §3.3 treats Carrier
Dichotomy, §3.4 Bridge Impossibility, §3.5 Coverage and Type-Finite
Foreclosure, §3.6 Attack Foreclosure v4, and §3.7 CTMT Recursion.

\subsection{Carrier Dichotomy}\label{sec:foreclosure:carrier-dichotomy}

The Carrier Dichotomy Theorem is the spine of Part I.  It turns the
informal question ``what happened to this attempted closure of \(X\)?''
into a typed audit question: once the target, the carrier, the
accepted-imports base, and the \TCD\ have been fixed, the carrier is
classified by its residual closure.  The theorem separates the
classical-residue case, where the carrier is not routed to \(X\), from
the three carrier-residual-equivalent foreclosure modes.

\begin{definition}[Typed primary carrier]\label{def:typed-primary-carrier}
A typed primary carrier aimed at \(X\) is the carrier-level datum
introduced in \S\ref{sec:foreclosure:setup}: an attempted proof
construction whose type signature targets the conjecture \(X\), carries
the residual \(\Xi_C(D\mid L)\) from \Cref{def:carrier-residual}, and
comes equipped with the accepted-imports base and audit setting in
which that residual is interpreted.
Formally, this section works with carriers satisfying
\[
  C\leadsto X,\qquad C\models\text{\TCD},\qquad
  R_C=\Xi_C(D_C\mid L_C),
\]
together with target-relative bridge predicates and carrier-internal
cascade predicates.  The target-relative predicates depend on \(X\);
the cascade predicates are intrinsic to \(C\)'s own residual cascade.
\end{definition}

\begin{definition}[Four typed carrier states]\label{def:four-typed-states}
For a typed primary carrier \(C\) aimed at \(X\), the four Part I
states are the following.  The \nonXequiv\ state says that the
residual closure of \(C\) is logically equivalent to some conjectural
content \(X'\) distinct from \(X\).  The three carrier-residual
equivalent states are \CRETE, the target-equivalent fallback;
\CREBF, the bridge-failed dead end; and \CRECTMT, the
cascade-terminal matrix-element case.  The last three are the
three typed foreclosure modes available once the residual closure is
\Xequiv.
Equivalently, writing \(\mathcal B\) for the accepted-imports base,
\[
\begin{aligned}
  \mbox{\nonXequiv}_X(C)
    &\Longleftrightarrow
      \exists X'.\, (R_C=0\leftrightarrow X')\wedge X'\neq X,\\
  \operatorname{OpenPending}_{\mathcal B}(X')
    &\Longleftrightarrow \neg \mathcal B(X'),\\
  \operatorname{NativeOutcome}_{\mathcal B}(C,X')
    &\Longleftrightarrow
      (\mathcal B(X')\wedge R_C=0)\vee
      \operatorname{OpenPending}_{\mathcal B}(X'),\\
  \mbox{\CRETE}_X(C)
    &\Longleftrightarrow
      (R_C=0\leftrightarrow X)\wedge
      \operatorname{SelfAuditFail}(C),\\
  \mbox{\CREBF}_X(C)
    &\Longleftrightarrow
      \operatorname{NoBridge}(C,X)\vee
      \operatorname{BridgeMismatch}(C,X),\\
  \mbox{\CRECTMT}_X(C)
    &\Longleftrightarrow
      \operatorname{MRed}(C)\wedge
      \neg\operatorname{IRDec}(C)\wedge
      \neg\operatorname{MDEqTarget}(C,X).
\end{aligned}
\]
\end{definition}

\begin{definition}[Coverage context]\label{def:coverage-ctx}
Fix a target conjecture \(X\) and an accepted-imports base
\(\mathcal B:\mathrm{Prop}\to\mathrm{Prop}\) with import soundness.  A
coverage context for \((X,\mathcal B)\), denoted
\(\operatorname{CoverageCtx}(X,C)\) when the typed carrier under audit is
\(C\), is an opaque admissibility predicate marking the audit setting in
which the four-state classification is intended to be exhaustive.  The
predicate is treated as primitive in the body of this paper.
\end{definition}

\begin{theorem}[Carrier Dichotomy]\label{thm:carrier-dichotomy}
Let \(\mathcal B:\mathrm{Prop}\to\mathrm{Prop}\) be an
accepted-imports base, and assume import soundness
\[
  \forall Y.\ \mathcal B(Y)\Longrightarrow Y .
\]
Let \(X\) be an open target conjecture and let \(C\) be a typed primary
carrier satisfying
\[
  C\leadsto X,\qquad C\models\text{\TCD},\qquad
  \operatorname{CoverageCtx}(X,C).
\]
Then the carrier audit has the following two-branch form:
\[
\begin{aligned}
  \mbox{\nonXequiv}_X(C)
  &\Longrightarrow
    \exists X'.\,
      (R_C=0\leftrightarrow X')\wedge X'\neq X\wedge
      \operatorname{NativeOutcome}_{\mathcal B}(C,X'),\\
  \mbox{\Xequiv}_X(C)
  &\Longrightarrow
    \operatorname{ExactlyOne}\bigl(
      \mbox{\CRETE}_X(C),\mbox{\CREBF}_X(C),\mbox{\CRECTMT}_X(C)
    \bigr).
\end{aligned}
\]
Here
\[
\begin{aligned}
\operatorname{ExactlyOne}(A,B,D)
  \Longleftrightarrow{}&
  (A\wedge\neg B\wedge\neg D)\vee
  (B\wedge\neg A\wedge\neg D)\\
  &{}\vee(D\wedge\neg A\wedge\neg B).
\end{aligned}
\]

\begin{enumerate}
\item If \(C\) is \nonXequiv, so that the residual closure of \(C\) is
  logically equivalent to some \(X'\neq X\), then the
  \NativeClosure\ of \(C\) is one of two kinds: either \(X'\) is
  already in the accepted-imports base and the carrier closes natively,
  or \(X'\) remains open-pending relative to that base.
\item If \(C\) is \Xequiv, so that the residual closure of \(C\) is
  logically equivalent to \(X\), then exactly one of the three
  carrier-residual-equivalent foreclosure modes applies:
  \[
    \textup{\CRETE},\qquad \textup{\CREBF},\qquad \textup{\CRECTMT}.
  \]
\end{enumerate}

Consequently, under the Coverage Theorem, every admissible typed
primary carrier aimed at \(X\) is exhausted by the four-state audit:
\[
  \mbox{\nonXequiv}_X(C)\vee
  \mbox{\CRETE}_X(C)\vee
  \mbox{\CREBF}_X(C)\vee
  \mbox{\CRECTMT}_X(C),
\]
with the three
\Xequiv\ modes disjoint in the audit classification.
\end{theorem}

\begin{proof}
Assume first that \(C\) is \nonXequiv.  By the definition of
\(\mbox{\nonXequiv}_X(C)\), there is a proposition \(X'\) such that
\[
  R_C=0\leftrightarrow X',
  \qquad
  X'\neq X .
\]
It remains only to prove the displayed native-outcome alternative for
this same \(X'\).  Since membership in the accepted-imports base is a
proposition, classical excluded middle gives the two cases
\(\mathcal B(X')\) and \(\neg\mathcal B(X')\).
\begin{itemize}
\item \emph{Case 1: \(\mathcal B(X')\).}  By import soundness,
  \(\mathcal B(X')\Rightarrow X'\), so \(X'\) holds.  Applying the
  implication \(X'\Rightarrow R_C=0\) from the equivalence
  \(R_C=0\leftrightarrow X'\), we obtain \(R_C=0\).  Therefore
  \[
    \mathcal B(X')\wedge R_C=0
  \]
  holds, which is the first disjunct of
  \(\operatorname{NativeOutcome}_{\mathcal B}(C,X')\).
\item \emph{Case 2: \(\neg\mathcal B(X')\).}  Then
  \(\operatorname{OpenPending}_{\mathcal B}(X')\) holds by definition.
  Hence the second disjunct of
  \(\operatorname{NativeOutcome}_{\mathcal B}(C,X')\) holds.
\end{itemize}
Combining the chosen \(X'\), the residual equivalence, the inequality
\(X'\neq X\), and the native-outcome disjunction proves the first
branch.

For the second branch, assume \(\mbox{\Xequiv}_X(C)\), i.e.
\[
  R_C=0\leftrightarrow X .
\]
The hypotheses \(C\leadsto X\), \(C\models\text{\TCD}\), and
\(\operatorname{CoverageCtx}(X,C)\) place the carrier in the admissible
coverage setting for the \(X\)-equivalent branch.  We now invoke the
same-paper coverage-exhaustiveness obligation for the
\(X\)-equivalent case.  That obligation says precisely that an
\(X\)-equivalent carrier in this context lies in exactly one of the
three carrier-residual-equivalent modes.  Therefore
\[
  \operatorname{ExactlyOne}\bigl(
    \mbox{\CRETE}_X(C),\mbox{\CREBF}_X(C),\mbox{\CRECTMT}_X(C)
  \bigr).
\]
The \CRETE\ alternative is the target-equivalent fallback whose audit
source is the Same-Level Self-Audit Failure theorem of
\FThree\ \cite{Tsiokos2026FoundationsIII}; the \CREBF\ and \CRECTMT\
alternatives are the bridge-failed and matrix-element-terminal cases
defined above.  This proves the second branch and hence the dichotomy.
\end{proof}\footnote{\raggedright The Lean theorem proves the structural
derivation in \texttt{SixBirdsMetaMath.MetaMath.\allowbreak
Foreclosure.CarrierDichotomy.\allowbreak carrier\_dichotomy}, applying
the same-paper obligation \texttt{SixBirdsMetaMath.MetaMath.\allowbreak
Foreclosure.CarrierDichotomy.\allowbreak carrier\_dichotomy\_t3\_pending}
under the coverage-context guard; the narrowing is recorded in
Appendix~\ref{app:formalization}.}

The next two foreclosure subsections use this classification in
different directions.  Bridge Impossibility asks what happens when a
non-\(X\)-equivalent carrier is enriched by a bridge into an
\(X\)-equivalent composite.  Type-Finite Foreclosure later uses the
Coverage Theorem to turn the carrier-state audit into a finite
door-space of possible outcomes.  Attack Foreclosure v4 then applies
that finite door-space to the three framework-internal moves introduced in
\S\ref{sec:foreclosure:setup}.

\subsection{Bridge Impossibility}\label{sec:foreclosure:bridge-impossibility}

The bridge case is the first immediate use of
\thmCarrierDichotomy.  A \nonXequiv\ carrier can sometimes be enriched
by adding a bridge to the target side; the question is what status that
bridge has if the enriched composite is now \Xequiv.  The corollary
below says that such a bridge cannot be treated as a free transport.
If the bridge is the one that composes the original carrier into the
\Xequiv\ composite, then the target-level strength has moved into the
bridge itself.

\begin{definition}[Carrier-side typed bridge]\label{def:carrier-side-typed-bridge}
A \TypedBridgeCarrier\ is a typed transport object used in the Part I
carrier audit: it is directed from a carrier-side comparison context
toward the target-side carrier data, and its composition with the
original carrier is part of the typed evidence under review.\footnote{This
is the carrier-side bridge of Part I.  It is distinct from the
calibration-side typed bridge used later in Part III for bridges
between calibrations.}
The composition relation is written
\[
  \operatorname{BridgeComp}_X(B,C,C_B),
\]
meaning that the bridge \(B\) composed with the original carrier \(C\)
produces the enriched composite carrier \(C_B\) for target \(X\).
\end{definition}

\begin{definition}[\(X\)-strength obligation]\label{def:x-strength-obligation}
An \XstrengthObl\ is the obligation carried by a bridge whose
composition with a non-\(X\)-equivalent carrier produces an
\(X\)-equivalent composite.  In standard terms, the bridge is no longer
merely a change of presentation; it must supply the target-strength
content needed for the composite to hit \(X\).
We write this predicate as \(\operatorname{XObl}(B,X)\).
\end{definition}

\begin{corollary}[Bridge Impossibility]\label{cor:bridge-impossibility}
Let \(\mathcal B\) be a sound accepted-imports base.  Let \(C\) be a
\nonXequiv\ typed primary carrier with \(R_C=0\).  Let \(B\) be a
carrier-side typed bridge, and let \(C_B\) be a typed composite carrier
such that
\[
  \operatorname{BridgeComp}_X(B,C,C_B),\qquad
  R_{C_B}=0,\qquad
  (R_{C_B}=0\leftrightarrow X),
\]
with \(C_B\leadsto X\), \(C_B\models\text{\TCD}\), and
\(\operatorname{CoverageCtx}(X,C_B)\).  Then
\[
  \operatorname{XObl}(B,X).
\]
Equivalently, \(B\) carries an \XstrengthObl: it cannot be a free
non-circular transport to \(X\).
\end{corollary}

\begin{proof}
The composite hypotheses give
\[
  \mbox{\Xequiv}_X(C_B),
\]
because \(R_{C_B}=0\leftrightarrow X\) is exactly the definition of
\(X\)-equivalence for the composite carrier.  The remaining hypotheses
on \(C_B\) give
\[
  C_B\leadsto X,\qquad C_B\models\text{\TCD},\qquad
  \operatorname{CoverageCtx}(X,C_B).
\]
Thus the second branch of \thmCarrierDichotomy applies to the composite
carrier \(C_B\), with the same sound accepted-imports base
\(\mathcal B\).  We obtain the three-way classification
\[
  \operatorname{ExactlyOne}\bigl(
    \mbox{\CRETE}_X(C_B),\mbox{\CREBF}_X(C_B),\mbox{\CRECTMT}_X(C_B)
  \bigr).
\]

The hypothesis \(\operatorname{BridgeComp}_X(B,C,C_B)\) is what ties
the bridge \(B\) to this particular composite; without it, the
classification of \(C_B\) would say nothing about an unrelated bridge.
Applying the same-paper bridge-composition obligation to this
composition relation and to the displayed three-way classification
therefore yields
\[
  \operatorname{XObl}(B,X).
\]
This is the required \(X\)-strength obligation for \(B\).
\end{proof}\footnote{\raggedright The Lean theorem proves the structural
derivation in \texttt{SixBirdsMetaMath.MetaMath.\allowbreak
Foreclosure.BridgeImpossibility.\allowbreak bridge\_impossibility},
applying the Carrier Dichotomy plus the same-paper obligation
\texttt{SixBirdsMetaMath.MetaMath.\allowbreak
Foreclosure.BridgeImpossibility.\allowbreak
bridge\_impossibility\_carries\_obligation} under the
bridge-composition guard; the narrowing is recorded in
Appendix~\ref{app:formalization}.}

This corollary is deliberately local to the carrier audit.  It is not a
general no-go theorem for external proof methods, and it does not claim
that every mathematical bridge into \(X\) is impossible.  It says only
that, inside the typed carrier setup of Part I, a bridge that turns a
non-\(X\)-equivalent carrier into an \(X\)-equivalent composite must
carry the target-strength content it appears to transport.  The next
subsections broaden the audit: the Coverage Theorem supplies
exhaustiveness, and Attack Foreclosure v4 applies the resulting
classification to the three framework-internal moves.

\subsection{Coverage}\label{sec:foreclosure:coverage}

The Coverage Theorem is the central structural content behind Part I.
The Carrier Dichotomy separates the branches once a carrier is known
to be \nonXequiv\ or \Xequiv; Coverage is the assertion that, under the
framework's admissible coverage context, no typed primary carrier aimed
at \(X\) escapes the four-state audit.  The coverage context is the
guard that marks the setting in which the target, carrier, accepted
imports, and typed audit discipline are all in force.

Coverage is therefore stronger than simply naming the four states.  It
asserts that the framework has no fifth admissible outcome for a typed
carrier aimed at \(X\).  A carrier whose residual closure points to
some \(X'\neq X\) is accounted for by the \nonXequiv\ branch.  A
carrier whose residual closure is equivalent to \(X\) must be audited
as \CRETE, \CREBF, or \CRECTMT.  The schema is the place where those
two sentences become an exhaustiveness claim rather than a vocabulary
choice.

\begin{schema}[Coverage Theorem]\label{thm:coverage}
Let \(X\) be an open target conjecture, and let \(C\) be a typed
primary carrier aimed at \(X\) under the \TCD.  Assume \(C\) lies in an
admissible coverage context for \(X\).  Then the four typed states are
exhaustive: \(C\) is \nonXequiv, or \(C\) is in \textup{\CRETE},
or \(C\) is in \textup{\CREBF}, or \(C\) is in \textup{\CRECTMT}.
Equivalently, under
\[
  C\leadsto X,\qquad C\models\text{\TCD},\qquad
  \operatorname{CoverageCtx}(X,C),
\]
every admissible typed primary carrier aimed at \(X\) satisfies the
four-way disjunction
\[
  \mbox{\nonXequiv}_X(C)\vee
  \mbox{\CRETE}_X(C)\vee
  \mbox{\CREBF}_X(C)\vee
  \mbox{\CRECTMT}_X(C).
\]
The assertion is exhaustiveness only; disjointness of the three
\Xequiv\ modes is the separate exactly-one conclusion used in
\thmCarrierDichotomy.
\end{schema}

\footnote{\raggedright Tracked in the formalization harness as a
conditional schema in \texttt{SixBirdsMetaMath.MetaMath.\allowbreak
Foreclosure.Coverage.\allowbreak coverage\_theorem}, conditioned on the
opaque coverage-context guard; see Appendix~\ref{app:formalization}.}

The mathematical burden of the schema is not the displayed disjunction
itself but the admissibility analysis that makes the disjunction
exhaustive.  A precise definition of ``open conjecture relative to an
accepted-imports base'' should start with a predicate
\(\mathcal B:\mathrm{Prop}\to\mathrm{Prop}\) and import soundness
\(\mathcal B(Y)\Rightarrow Y\).  Relative openness should then rule out
the target being already settled by the base: informally, neither
\(\mathcal B(X)\) nor \(\mathcal B(\neg X)\) should hold in the audit
context.  It must also be stable under the carrier's residual
equivalences, so that replacing \(X\) by a proposition \(X'\) detected
by \(R_C=0\leftrightarrow X'\) does not smuggle accepted content into
an allegedly open target.

The corresponding relative-consistency argument is the load-bearing
piece.  One must show that an admissible carrier cannot escape the
four-state classification by exploiting an inconsistent or
over-generous accepted-imports base.  If the residual closure is
equivalent to some \(X'\neq X\), the carrier is in the
\nonXequiv\ branch.  If the residual closure is equivalent to \(X\),
then the target-equivalent carrier must be audited as one of the
three carrier-residual-equivalent modes: same-level target-equivalence
failure, bridge failure, or CTMT terminality.  The missing proof is the
relative-consistency theorem saying that no admissible fifth case is
available once the coverage context has ruled out malformed targets,
malformed carriers, and base-inconsistent native closure.  Until that
argument is supplied, Coverage remains the central obligation rather
than a proved classification theorem.

This is why Coverage appears twice in the Part I dependency graph.
First, it supplies the exhaustiveness dependency used by the Carrier
Dichotomy's \Xequiv\ branch.  Second, it is the input for Type-Finite
Foreclosure: once every admissible carrier falls into one of the four
states, the infinite carrier space can be audited through a finite
door-space of typed outcomes.  The next subsection states that
finite-door-space consequence.

\subsection{Type-Finite Foreclosure}\label{sec:foreclosure:type-finite}

Coverage makes the carrier audit exhaustive; Type-Finite Foreclosure
extracts the finite shape of that audit.  The carrier space may be
infinite: there may be indefinitely many attempted carriers, pivots,
presentations, or reductions aimed at the same target \(X\).  The
door-space is different.  It is the space of typed audit outcomes, and
after Coverage that space has exactly four named entries.

\begin{theorem}[Type-Finite Foreclosure]\label{thm:type-finite-foreclosure}
Let \(X\) be an open target conjecture, and let \(C\) range over typed
primary carriers aimed at \(X\) under the \TCD\ and an admissible
coverage context.  Then the carrier-state door-space is type-finite:
every such \(C\) is classified by one of the four states
\nonXequiv, \textup{\CRETE}, \textup{\CREBF}, or \textup{\CRECTMT}.
Thus the codomain of audit classifications has exactly four entries,
and the image of any family of typed carriers aimed at \(X\) has
cardinality at most four, even when the carrier family itself is
infinite.
Formally, for
\[
  \mathsf{State}
  =\{\mathsf{nonX},\mathsf{TE},\mathsf{BF},\mathsf{CTMT}\},
\]
define
\[
\begin{aligned}
  \operatorname{ClassifiesAs}(\mathsf{nonX},X,C)
    &\Longleftrightarrow \mbox{\nonXequiv}_X(C),\\
  \operatorname{ClassifiesAs}(\mathsf{TE},X,C)
    &\Longleftrightarrow \mbox{\CRETE}_X(C),\\
  \operatorname{ClassifiesAs}(\mathsf{BF},X,C)
    &\Longleftrightarrow \mbox{\CREBF}_X(C),\\
  \operatorname{ClassifiesAs}(\mathsf{CTMT},X,C)
    &\Longleftrightarrow \mbox{\CRECTMT}_X(C).
\end{aligned}
\]
Then
\[
  C\leadsto X\wedge C\models\text{\TCD}\wedge
  \operatorname{CoverageCtx}(X,C)
  \Longrightarrow
  \exists s\in\mathsf{State}.\,
    \operatorname{ClassifiesAs}(s,X,C).
\]
\end{theorem}

\begin{proof}
Let
\[
  \mathsf{State}
  =\{\mathsf{nonX},\mathsf{TE},\mathsf{BF},\mathsf{CTMT}\}
\]
with \(\operatorname{ClassifiesAs}\) defined as in the statement.
Take any admissible carrier \(C\) aimed at \(X\).  By the Coverage
Theorem,
\[
  \mbox{\nonXequiv}_X(C)\vee
  \mbox{\CRETE}_X(C)\vee
  \mbox{\CREBF}_X(C)\vee
  \mbox{\CRECTMT}_X(C).
\]
We split on this disjunction.  In the first case choose
\(\mathsf{nonX}\); by definition
\(\operatorname{ClassifiesAs}(\mathsf{nonX},X,C)\).  In the second
case choose \(\mathsf{TE}\); in the third choose \(\mathsf{BF}\); in
the fourth choose \(\mathsf{CTMT}\).  Thus for every admissible
carrier \(C\) there exists \(s\in\mathsf{State}\) such that
\(\operatorname{ClassifiesAs}(s,X,C)\).

Consequently the audit map \(C\mapsto s\), whenever a representative
state is selected, has codomain \(\mathsf{State}\).  Since
\(\mathsf{State}\) has four constructors, the image of any family of
admissible \(X\)-aimed carriers has cardinality at most four.  This is
the asserted type-finiteness of the door-space; no finiteness of the
underlying carrier family is used.
\end{proof}\footnote{\raggedright The Lean theorem proves the structural
derivation in \texttt{SixBirdsMetaMath.MetaMath.\allowbreak
Foreclosure.TypeFiniteForeclosure.\allowbreak
type\_finite\_foreclosure}, applying the same-paper obligation
\texttt{SixBirdsMetaMath.MetaMath.\allowbreak Foreclosure.Coverage.\allowbreak
coverage\_theorem} under the coverage-context guard; the narrowing is
recorded in Appendix~\ref{app:formalization}.}

This theorem is the bridge from carrier classification to operational
foreclosure.  Attack Foreclosure v4 will use the finite door-space to
audit the three framework-internal moves: carrier pivot, bridge import,
and cascade-internal computation.  CTMT Recursion then explains why the
matrix-element terminal case can reappear at nested depths without
creating a fifth kind of carrier state.

\subsection{Attack Foreclosure v4}\label{sec:foreclosure:attack-foreclosure-v4}

Attack Foreclosure v4 is the operational meta-claim of Part I.  The
preceding subsections classified carriers and bridges; this subsection
packages that classification as a statement about moves.  The point is
not that mathematicians cannot make progress on \(X\).  The point is
that progress cannot be manufactured solely by the framework's internal
machinery.  Once \(X\) is in framework scope, every internal move is
audited as either staying inside the finite carrier-state door-space or
exposing mathematical content that must be supplied from outside that
machinery.

Two restrictions in this statement are essential.  First, \(X\) must be
in framework scope: the audit is not a universal claim about all
propositions, but a claim about targets for which the framework's
typing and closure-formation discipline is active.  Second, ``advances
\(X\) without external content'' means that the move would make
target-relevant progress using only the framework's internal
operations, without naming and accepting an external mathematical
source.  The schema says that the three internal move classes do not
have that power.

The move classes themselves are not banned.  A \carrierPivot\ may
replace an unhelpful presentation by a sharper one.  A \bridgeImport\
may identify the exact theorem needed to transport content into the
target type.  A \cascadeInternal\ may reduce a residual to a more
focused terminal family.  The foreclosure claim is narrower: after the
typed audit, each such move either remains inside the four-state
classification or points to content that must be supplied as
\FwExternal\ mathematics.

\begin{definition}[Framework-scope predicate]\label{def:in-scope}
For a target conjecture \(X\), the framework-scope predicate
\(\operatorname{InScope}(X)\) marks \(X\) as a target to which the
framework-internal-move audit applies.  The predicate is treated as
primitive in the body.
\end{definition}

\begin{schema}[Attack Foreclosure v4]\label{thm:attack-foreclosure-v4}
Let \(X\) be a target conjecture in framework scope.  Let \(m\) be a
\fim, where
\[
  m\in\{\mbox{\carrierPivot},\mbox{\bridgeImport},
  \mbox{\cascadeInternal}\}.
\]
Then \(m\) does not advance \(X\) without importing external
mathematical content:
\[
  \operatorname{InScope}(X)
  \Longrightarrow
  \neg\,\operatorname{AdvancesWithoutExternal}(m,X).
\]
Equivalently, a framework-internal move has one of two typed outcomes:
either it remains inside the four-state carrier classification from
\thmCarrierDichotomy\ and \thmCoverage, without closing \(X\) by
internal means, or it exposes \FwExternal\ content that must be
accepted through the closure-formation discipline rather than generated
by the framework-internal move itself.  The same statement applies
pointwise to chains by auditing each move in the chain.
\end{schema}
\footnote{\raggedright Tracked in the formalization harness as a
conditional schema in \texttt{SixBirdsMetaMath.MetaMath.\allowbreak
Foreclosure.AttackForeclosure.\allowbreak attack\_foreclosure\_v4},
conditioned on the framework-scope guard composed with opaque
internal-move semantics; see Appendix~\ref{app:formalization}.}

The intended derivation shape is nevertheless visible from the Part I
spine.  A \carrierPivot\ replaces the current carrier by another
carrier aimed at \(X\); the Carrier Dichotomy and Coverage then return
that alternative carrier to the four-state audit.  A \bridgeImport\
tries to turn a \nonXequiv\ carrier into an \Xequiv\ composite; Bridge
Impossibility says the bridge must carry the \(X\)-strength obligation
rather than transporting target content for free.  A
\cascadeInternal\ continues the reduction inside the carrier; the CTMT
analysis says that such a reduction either closes by accepted external
content or reaches a matrix-element terminus, possibly at a deeper
level.  The informal part is the universal packaging of these three
move analyses into a single all-chains foreclosure statement.
Symbolically, the intended case analysis is
\[
\begin{aligned}
  m=\mbox{\carrierPivot}
    &\Longrightarrow
    C'\leadsto X\ \wedge\
    \exists s\in\mathsf{State}.\,
      \operatorname{ClassifiesAs}(s,X,C'),\\
  m=\mbox{\bridgeImport}
    &\Longrightarrow
    \operatorname{BridgeComp}_X(B,C,C_B)
    \Rightarrow \operatorname{XObl}(B,X),\\
  m=\mbox{\cascadeInternal}
    &\Longrightarrow
    R_C \text{ closes or enters a CTMT tower.}
\end{aligned}
\]
The obligation in the schema is the uniform quantification over this
three-way case split under the framework-scope guard.

For a chain \(m_1,\ldots,m_k\) of internal moves, the intended
invariant is pointwise.  Suppose that, before \(m_i\), the current
object is a typed carrier, typed composite, or CTMT tower already under
the audit discipline.  If \(m_i\) is a pivot, the audit restarts on a
new carrier aimed at the same \(X\), and Type-Finite Foreclosure sends
that carrier back to the four-state door-space.  If \(m_i\) is a
bridge import, the bridge is either unrelated to the composite and so
does not apply, or it is the bridge in a typed composition and Bridge
Impossibility transfers the target strength into an
\(\operatorname{XObl}\) obligation.  If \(m_i\) is cascade-internal,
then either the residual is closed by accepted content or the carrier
enters the CTMT tower described in \Cref{thm:ctmt-recursion}.  Thus an
iterated internal chain preserves the same alternative: return to the
finite audit classification, or expose external content.  The
unproved part is the global theorem that this invariant captures every
admissible chain of framework-internal moves; that is why v4 is stated
as an obligation.

This is a carrier-classification theorem, not a counterexample
construction against a named class of ordinary proof methods.  The
classical barriers literature rules out particular families of
techniques by showing that they cannot distinguish the relevant
worlds.  Attack Foreclosure v4 instead partitions the framework's
typed attempts into a finite audit enumeration and says that the
framework-internal part of the attempt does not itself supply the
external mathematical content needed for closure.

Class-B Irremovability from \PaperPvNP\
\cite{Tsiokos2026PvNP} is positioned as a fourth methodological
barrier alongside (and complementary to, not a strengthening of) the
classical relativization, natural-proofs, and algebrization barriers
\cite{BakerGillSolovay1975Relativization,RazborovRudich1997NaturalProofs,AaronsonWigderson2009Algebrization}.
That positioning is methodological: it explains where the Six Birds
foreclosure account sits relative to familiar proof-barrier language,
without claiming that Attack Foreclosure v4 subsumes those barriers or
converts them into framework theorems.

The next subsection isolates the remaining cascade-specific issue.
CTMT Recursion shows why the matrix-element terminal case can fire at
nested depths without changing the four-state door-space.  The
recognition-mode section then supplies the v5 refinement: derivation
routes remain foreclosed, while recognition-mode landings are treated
as structurally distinct from framework-internal derivations.

\subsection{CTMT Recursion}\label{sec:foreclosure:ctmt-recursion}

The \CRECTMT\ state in the Carrier Dichotomy is not necessarily a
one-step endpoint.  A carrier's cascade may reduce the residual to a
carrier-native matrix-element family; the decision of that family may
then reduce to a deeper family; and the same pattern may continue.  The
CTMT Recursion Theorem records this depth nesting as a type-theoretic
object rather than treating every terminal matrix-element case as a
flat terminal event.

\begin{theorem}[CTMT Recursion]\label{thm:ctmt-recursion}
For any target \(X\) and typed primary carrier \(C\), the
\CRECTMT\ state admits a depth-indexed recursion tower.  The base case
is the depth-\(1\) tower: \(C\) is in \CRECTMT\ exactly when \(C\) has
a one-level matrix-element terminus.  If \(C\)'s cascade reduces to a
deeper carrier-native matrix-element family whose own decision carries
a CTMT tower of positive depth, then \(C\) carries the next tower
level.  Formally, for any cascade semantics \(\sigma\),
\[
  \mathsf{Tower}_{\sigma}(X,C,1)
  \Longleftrightarrow
  \mbox{\CRECTMT}_X(C),
\]
and for every \(n\geq 0\),
\[
  \mathsf{Tower}_{\sigma}(X,C',n+1)\wedge
  \operatorname{Reduces}_{\sigma}(X,C,C')
  \Longrightarrow
  \mathsf{Tower}_{\sigma}(X,C,n+2).
\]
Thus matrix-element terminality can fire at depth \(1\), depth \(2\),
and iteratively at any finite positive depth admitted by the
inherited-records base.
\end{theorem}

\begin{proof}
Fix a cascade semantics \(\sigma\), a target \(X\), and a carrier
\(C\).  We first prove the depth-\(1\) identification.  If
\(\mathsf{Tower}_{\sigma}(X,C,1)\) holds, inspect the last constructor
of this tower.  The step constructor always has depth \(n+2\), hence
positive depth strictly greater than \(1\).  Therefore a tower of depth
\(1\) can only have been built by the base constructor, and the base
constructor carries exactly a proof of \(\mbox{\CRECTMT}_X(C)\).
Thus
\[
  \mathsf{Tower}_{\sigma}(X,C,1)
  \Longrightarrow \mbox{\CRECTMT}_X(C).
\]
Conversely, if \(\mbox{\CRECTMT}_X(C)\) holds, applying the base
constructor gives \(\mathsf{Tower}_{\sigma}(X,C,1)\).  The two
directions prove
\[
  \mathsf{Tower}_{\sigma}(X,C,1)
  \Longleftrightarrow \mbox{\CRECTMT}_X(C),
\]
and also show uniqueness of the depth-\(1\) construction: there is no
step-form tower at depth \(1\).

For the recursive clause, fix \(n\geq 0\), carriers \(C\) and \(C'\),
a deeper tower
\(\mathsf{Tower}_{\sigma}(X,C',n+1)\), and a reduction
\(\operatorname{Reduces}_{\sigma}(X,C,C')\).  The step constructor of
the tower definition takes exactly these two pieces of data and
produces
\[
  \mathsf{Tower}_{\sigma}(X,C,n+2).
\]
This proves the one-step extension.  Iterating the same constructor
along a finite sequence of reductions yields towers of arbitrary
finite positive depth admitted by the sequence, with the depth-\(1\)
base already identified with ordinary \CRECTMT.
\end{proof}\footnote{\raggedright This theorem is verified in Lean as
\texttt{SixBirdsMetaMath.MetaMath.\allowbreak
Foreclosure.CTMTRecursion.\allowbreak ctmt\_recursion}; see
Appendix~\ref{app:formalization}.}

The
deepest documented cascade in the present evidence base is an
eight-layer chain on the RH Branch B recursion; a track-independent
depth bound is an open problem in \Cref{sec:predictions:ctmt-depth-bound}.

The theorem closes Part I by isolating the remaining internal
recursion in the negative half.  Carrier pivoting, bridge import, and
cascade-internal computation have now each been assigned their typed
role: pivoting returns to the carrier-state audit, bridge import
creates an \(X\)-strength obligation when it reaches an
\(X\)-equivalent composite, and cascade-internal computation can
descend through a CTMT tower without creating a new kind of state.  The
paper now turns to the positive half: recognition-mode landing, where
progress is not claimed to arise from a framework-internal derivation
but from a named external source accepted by the closure-formation
discipline.

\begin{table}[H]
  \centering
  \caption{Part~I foreclosure summary: the six T-rows of the carrier-state analysis.}
  \label{tab:foreclosure:summary}
  \scriptsize
  \renewcommand{\arraystretch}{1.08}
  \begin{tabularx}{\textwidth}{@{}>{\raggedright\arraybackslash}p{0.06\textwidth}>{\raggedright\arraybackslash}X>{\raggedright\arraybackslash}p{0.16\textwidth}>{\raggedright\arraybackslash}p{0.18\textwidth}>{\centering\arraybackslash}p{0.08\textwidth}@{}}
    \toprule
    T\# & Statement & Lean coverage & Opacity guard & Spine? \\
    \midrule
    T1 & Carrier Dichotomy: every admissible typed carrier aimed at \(X\) is classified by the non-\(X\)-equivalent branch or the \CRETE/\CREBF/\CRECTMT{} branch. & partial & coverage-context & yes \\
    T2 & Bridge Impossibility: a carrier-side bridge whose composition reaches an \(X\)-equivalent composite carries an \(X\)-strength obligation. & partial & bridge-composition & no \\
    T3 & Coverage Theorem: the four typed states are exhaustive for admissible carriers under the audit discipline. & obligation & coverage-context & yes \\
    T4 & Type-Finite Foreclosure: the carrier outcome door-space has at most four typed states. & partial & coverage-context & no \\
    T5 & Attack Foreclosure v4: framework-internal moves do not advance \(X\) without external mathematical content. & obligation & framework-scope; opaque internal-move semantics & yes \\
    T6 & CTMT Recursion: \CRECTMT{} terminality forms a depth-indexed tower under cascade reduction. & theorem & none & no \\
    \bottomrule
  \end{tabularx}
\end{table}
  \section{Part II --- Recognition-Mode Landing Template}\label{sec:recognition}

\subsection{Column classification}\label{sec:recognition:column-classification}

Part I explained the negative half of the meta-mathematical claim: the
framework's internal derivation moves do not by themselves create the
external mathematical content needed for a Clay-class landing. Part II is
the positive-half complement. It describes the \RML\ path, where the
closure of a target is achieved under the SB closure assumption by
recognizing a named structural source rather than by presenting that source
as the output of a framework-internal derivation.

The organizing distinction is the V-Differential column-membership
classification. In plain terms, this is a substrate classification that
asks whether the substrate's structural witness is exposed by the problem
or stays compressed into trace-state evidence. A substrate in the
\VDWCE\ has a witness object directly produced by the closure route's
substrate-specific apparatus. A substrate in the \VDTSO\ has witness data
compressed into trace-state evidence rather than directly produced as the
witness object. The trace-state-only column is the column to which the
\SevenStage\ applies.

\begin{definition}[V-Differential trace-state-only column]\label{def:vdiff-tso-column}
A substrate lies in the \VDTSO\ when the load-bearing witness for closure
is not exposed as a direct witness object by the problem, but is instead
available through trace-state evidence that must be named and accepted by
the closure-formation discipline.
Writing \(S\) for a substrate,
\[
\begin{aligned}
  \mbox{\VDTSO}(S)
    &\Longleftrightarrow
      \operatorname{WitnessCompressedAsTraceState}(S),\\
  \mbox{\VDWCE}(S)
    &\Longleftrightarrow
      \operatorname{WitnessExposedByProblem}(S).
\end{aligned}
\]
The first predicate is the column membership used by the
recognition-mode template.  The second records the contrasting case:
the substrate exposes its load-bearing witness through its own
problem-specific closure apparatus.
\end{definition}

The three substrate examples used throughout the paper sit on the two
sides of this distinction. The conditional closure of Navier--Stokes
regularity landed by \PaperNS\ under the SB closure assumption of
Foundations I is a witness-content-exposing row: its closure route is
carried by substrate-specific apparatus, including the layer-dissolving
membrane, rather than by a recognition-mode template
\cite{Tsiokos2026NavierStokes}. The conditional closure of the Riemann
hypothesis landed by \PaperRH\ under the SB closure assumption of
Foundations I is trace-state-only: the relevant recognition source is
\(\Gamma_{\mathrm{SDTC}}^{\mathrm{Selberg}}\), imported as named
recognition per Foundations I closure formation from the SDTC apparatus
and the RH landing papers \cite{Tsiokos2026SDTC,Tsiokos2026RH}, with
\cite{Selberg1956TraceFormula} as the classical source for the
underlying trace formula. The
conditional closure of \(P\) versus \(NP\) landed by \PaperPvNP\ under the
SB closure assumption of Foundations I is also trace-state-only: the
recognition source is \(\Gamma_{\mathrm{csl\mbox{-}sat\mbox{-}hidden}}\),
imported as named recognition per Foundations I closure formation from
the CSL source and the PvNP landing papers
\cite{Tsiokos2026PvNP,Tsiokos2026CSL}.

This terminology is not a hierarchy of rigor. Recognition-mode landings
are peer with derivation-mode landings under the SB closure assumption.
The \RMmode\ path and the \DMmode\ path are structurally distinct: the
former records named recognition of a load-bearing source, while the
latter records a derivation route through the framework's internal
machinery. The point of Part II is to state the recognition-mode path in a
form that is as typed and auditable as the foreclosure path of Part I, not
to make recognition mode subordinate to derivation mode.

The four T-rows in this part are arranged accordingly. The
Recognition-Mode Template Applicability Theorem says that the seven-stage
template applies to formed closures in the trace-state-only column. The
Cross-Track Verdict Signature records the common Stage-6 sweep signature
seen in the two recognition-mode rows. The Closure-Content Thesis names
the load-bearing structural content as formed-layer closure content rather
than framework-primitive output. The Attack Foreclosure v5 refinement then
combines the two halves: derivation routes remain subject to v4
foreclosure, while recognition-mode landings import their load-bearing
source as named recognition.

All four Part~II main claims carry paper-proposal-grade status.
They are two-instance verified on PvNP and RH, and remain
corpus-pending until a third-instance cross-track replication on a
V-Differential trace-state-only substrate (Hodge is the leading
candidate) and a structural argument from column-membership to
verdict pattern are supplied. See Appendix~\ref{app:formalization} for the per-theorem
opacity-guard register.

The reading order is therefore direct. The next subsection states the
template and its seven stages. The following subsections record the
cross-track verdict signature, the closure-content thesis, and the v5
refinement, before the paper turns in \Cref{sec:validation} to the three
substrate rows that motivate this column classification.

\subsection{Template applicability}\label{sec:recognition:template-applicability}

The \SevenStage\ is the positive-half spine of the paper. It packages the
recognition-mode route as a typed procedure for formed-layer closures:
first build a closure object under the SB closure assumption, then audit
that the closure is not smuggling the needed source, and finally name the
source as accepted recognition content. The template is anchored in
\FOne\ closure formation \cite{Tsiokos2026EmergenceCalculus} and in the
\FTwo\ seven-schema admissibility background
\cite{Tsiokos2026FoundationsII}; this paper treats those inputs as
apparatus and records the cross-track applicability claim as
\PPGrade, \TwoInstance.

\StageOne. The construction begins by forming a saturated closure layer.
In ordinary terms, this means building the closure object with enough
bookkeeping, type separation, level structure, activation thresholds, and
visibility data to make it admissible for the later recognition audit. The
stage also includes an anti-tautology exhibit: the closure must add reach
over the carriers available before the formed layer was introduced.

\StageTwo. The second stage supplies a translation theorem from the
framework-internal predicate attached to the closure to the standard
mathematical target statement. This theorem is not where the
load-bearing source enters. Its role is to ensure that once the
recognition source is accepted by closure formation, the accepted content
lands on the standard target rather than remaining only a framework-local
statement.

\StageThree. The applicability audit checks that the relevant master
theorem applies to the formed closure. For CSL-hiddenness substrates this
means the CSL biconditional is structurally applicable; for
self-dual-trace substrates it means the duality-confinement apparatus is
structurally applicable. The audit separates hypotheses already verified
from closure structure from residual substrate-specific content still
needed for the landing.

\StageFour. The smuggle audit asks whether the closure has secretly
assumed the recognition source through its admissibility hypotheses,
lawfulness checks, saturation conditions, or empirical-bridge
bookkeeping. A positive finding here forces a repair of the proof program.
A clean pass means only that the closure has not already smuggled the
source; it does not by itself derive the source.

\StageFive. Anti-tautology hardening strengthens the initial
anti-tautology exhibit at the operational-predicate level. The point is
to rule out a closure whose apparent gain decomposes into independent
claims about independent carriers. The recognition-mode template requires
the operational predicate to carry genuine dependence on the formed
closure's data.

\StageSix. The three-option derivation sweep tests whether the
load-bearing residual can be obtained by a framework-primitive route, by
a non-descent route, or by elevation inside V-Differential. This stage is
where the cross-track verdict signature of the next subsection is
recorded. Here we only name the stage: the detailed signature and its
two-instance status belong to the Cross-Track Verdict Signature schema.

\StageSeven. Recognition closure names the load-bearing structural source
as accepted under Foundations I closure formation. For PvNP this source is
\(\Gamma_{\text{csl-sat-hidden}}\), imported as named recognition per
Foundations I closure formation from the CSL apparatus. For RH it is
\(\Gamma_{\text{SDTC}}^{\text{Selberg}}\), imported as named recognition
per Foundations I closure formation from the SDTC and RH apparatus. In
both cases the source is imported as named recognition, not derived from
framework primitives.

The structured template value used in the schema is a seven-field datum
\[
  T(c,X)=\langle S_1,S_2,S_3,S_4,S_5,S_6,S_7\rangle .
\]
Here \(S_1,\ldots,S_7\) are, respectively, \StageOne, \StageTwo,
\StageThree, \StageFour, \StageFive, \StageSix, and \StageSeven.  The
separate guard \(\operatorname{Applies}(c,X,T)\) asserts
that these seven stages actually hold for the formed closure \(c\) and
target \(X\); the mere existence of the seven-field record is not by
itself the applicability claim.

\begin{definition}[Template-applies predicate]\label{def:applies}
For a substrate \(c\), a target \(X\), and the seven-stage template
\(T\), the predicate \(\operatorname{Applies}(c,X,T)\) marks \((c,X)\)
as a pair on which the template's seven-stage audit succeeds.  The
predicate is treated as primitive in the body.
\end{definition}

\begin{schema}[Recognition-Mode Template Applicability]\label{thm:template-applicability}
Let \(c\) be a formed Six Birds closure and let \(X\) be its target
statement.  Suppose
\[
  \begin{gathered}
    \mbox{\VDTSO}(c),\qquad \operatorname{SBClosure}(c),\\
    T:\mathsf{7Stage}(c,X),\qquad
    \operatorname{Applies}(c,X,T).
  \end{gathered}
\]
Then the closure lands at theorem-grade:
\[
  \operatorname{TheoremGrade}(c,X).
\]
The applicability claim is \PPGrade\ and \TwoInstance: the two verified
instances are PvNP via \(\Gamma_{\text{csl-sat-hidden}}\) and RH via
\(\Gamma_{\text{SDTC}}^{\text{Selberg}}\).
\end{schema}\footnote{\raggedright Tracked in the formalization
harness as a conditional schema in
\texttt{SixBirdsMetaMath.MetaMath.\allowbreak
RecognitionMode.Template.\allowbreak recognition\_template\_applicability},
conditioned on the template-applies guard composed with opaque
substrate-typing predicates (formed-layer closure, V-Differential
trace-state-only column, SB closure assumption, lands-at-theorem-grade);
paper-proposal-grade; see Appendix~\ref{app:formalization}.}

The logical shape is a guarded implication.  The column hypothesis
restricts the schema to trace-state-only substrates.  The SB closure
hypothesis supplies the Foundations~I closure assumption, the
seven-field template value supplies the stage data, and the opaque
applies guard asserts that the stage data has actually been satisfied.
Only with all four inputs does the theorem-grade conclusion follow.  In
the display, \(\mathsf{7Stage}\), \(\operatorname{Applies}\), and
\(\operatorname{TheoremGrade}\) abbreviate the seven-stage-template
type, the template-applies guard, and the theorem-grade landing
predicate.

The substantive content of the obligation is the claim that the guard
\(\operatorname{Applies}(c,X,T)\) is the right audit witness for
trace-state-only substrates.  The closure-content thesis of
\Cref{thm:closure-content-thesis} says that, in this column, the
load-bearing content is named recognition rather than
framework-primitive derivation.  Thus the recognition source is not an
extra decorative label; it is the structural content that lets the
formed closure land.  The seven stages are the audit proving that this
content is being used in the sanctioned way.  \StageOne\ constructs the
formed closure, \StageTwo\ connects the closure predicate to the target,
\StageThree--\StageSix\ run the applicability, smuggle,
anti-tautology, and derivation-sweep audits, and \StageSeven\ names the
recognition source as accepted closure content.  The obligation is that
when a TSO substrate supplies a formed closure, the SB closure
assumption, and a successful seven-stage audit, theorem-grade landing
follows.

The two verified rows are the trace-state-only substrate rows already
named in the column classification. The PvNP row uses the conditional
closure of \(P\) versus \(NP\) landed by \PaperPvNP\ under the SB closure
assumption of Foundations I, with \(\Gamma_{\text{csl-sat-hidden}}\)
imported as named recognition per Foundations I closure formation from
\PaperCSL\ \cite{Tsiokos2026PvNP,Tsiokos2026CSL}. The RH row uses the
conditional closure of the Riemann hypothesis landed by \PaperRH\ under
the SB closure assumption of Foundations I, with
\(\Gamma_{\text{SDTC}}^{\text{Selberg}}\) imported as named recognition
per Foundations I closure formation from \PaperSDTC\ and \PaperRH\
\cite{Tsiokos2026SDTC,Tsiokos2026RH}. These two rows justify the
\TwoInstance\ qualifier; theorem-grade universal lifting remains
corpus-pending for third-instance cross-track replication.

The next subsection isolates \StageSix. It records the common
cross-track verdict signature for the three derivation-sweep options,
without changing the status of the template applicability claim itself:
\Cref{thm:template-applicability} remains \PPGrade\ and
\TwoInstance, while \Cref{thm:verdict-signature} states the Stage-6
signature that both verified rows exhibit.

\subsection{Cross-track verdict signature}\label{sec:recognition:verdict-signature}

\StageSix\ is the template's stress test for derivation routes. The
formed closure is run against three independent options, each asking
whether the load-bearing source can be obtained without invoking the
recognition step. The cross-track observation in this subsection is
deliberately limited: the \SigPattern\ signature is two-instance verified
on PvNP and RH, and remains corpus-pending outside those two rows.

\begin{definition}[Three-option derivation sweep]\label{def:three-option-sweep}
The three-option derivation sweep is the Stage-6 sub-procedure that runs a
trace-state-only closure attempt against three options. The \OptFP\ asks
whether the closure is reachable by framework-primitive descent. The
\OptSAU\ asks whether the attempted derivation stops at non-descent at the
SAU boundary. The \OptVDE\ asks whether elevating the data inside
V-Differential resolves the closure route.
\end{definition}

The sweep records one of three verdicts for each option. Verdict (i)
accepts the option as the closure route. Verdict (ii) marks the option as
not the closure route but not itself blocking. Verdict (iii) elevates the
attempt to a V-Differential alternative. On the two verified
trace-state-only rows, the observed signature is this: the
framework-primitive option returns verdict (ii), the SAU non-descent
option returns verdict (iii), and the V-Differential elevation option
returns verdict (ii). Thus the \SigPattern\ signature is two-instance
verified on PvNP and RH.

Formally, Stage~6 has a three-element option type and a three-element
verdict type:
\[
  \mathsf{SweepOption}
    =\{\mathsf{FP},\mathsf{SAU},\mathsf{VDE}\},
  \qquad
  \mathsf{Verdict}=\{\mathrm{i},\mathrm{ii},\mathrm{iii}\}.
\]
For a seven-stage template value \(T\), the verdict reader is the
opaque function
\[
  v_T:\mathsf{SweepOption}\longrightarrow\mathsf{Verdict}.
\]
The two-instance signature asserted below is the assignment
\[
  v_T(\mathsf{FP})=\mathrm{ii},\qquad
  v_T(\mathsf{SAU})=\mathrm{iii},\qquad
  v_T(\mathsf{VDE})=\mathrm{ii}.
\]
The two observed rows are summarized in
\Cref{tab:verdict-signature}.

\begin{schema}[Cross-Track Verdict Signature]\label{thm:verdict-signature}
Let \(c\) be a formed closure in the \VDTSO, let \(X\) be its target,
and let \(T:\mathsf{7Stage}(c,X)\) satisfy
\(\operatorname{Applies}(c,X,T)\).  Then the Stage~6 verdict reader
for the two verified cross-track instances satisfies
\[
  v_T(\mathsf{FP})=\mathrm{ii}\ \wedge\
  v_T(\mathsf{SAU})=\mathrm{iii}\ \wedge\
  v_T(\mathsf{VDE})=\mathrm{ii}.
\]
In prose: the \OptFP\ returns verdict (ii), the \OptSAU\ returns
verdict (iii), and the \OptVDE\ returns verdict (ii).
\end{schema}\footnote{\raggedright Tracked in the formalization
harness as a conditional schema in
\texttt{SixBirdsMetaMath.MetaMath.\allowbreak
RecognitionMode.VerdictSignature.\allowbreak verdict\_signature},
conditioned on the template-applies guard composed with the opaque
sweep-verdict; paper-proposal-grade, two-instance verified; see
Appendix~\ref{app:formalization}.}

The discussion is per-option rather than by a single undifferentiated
verdict.  The framework-primitive option, the SAU non-descent option,
and the V-Differential elevation option are checked as separate cases.
The function \(v_T\) is opaque in the
formalization, so the displayed equalities are not definitional
reductions; they are exactly the \SigPattern\ content recorded for the
two verified TSO rows.

The structural reason for the \SigPattern\ shape is supplied by the
closure-content thesis, but only as a pending argument rather than as a
complete proof.  The \OptFP\ returns verdict (ii) because, on a
trace-state-only substrate, a verdict (i) here would say that the
load-bearing source is derivable from framework primitives; that is
exactly the pattern ruled out by the closure-content thesis.  The
\OptSAU\ returns verdict (iii) because non-descent at the SAU boundary
does not block the closure: the source has already been named as
recognition content, so the audit elevates to the V-Differential
trace-state classification.  The \OptVDE\ then returns verdict (ii)
because this elevation is a structural relabeling of the route, not a
new derivation of the load-bearing source.  A future theorem-grade
lifting of the verdict-signature schema would have to turn this
per-option explanation into a general argument from V-Differential
column membership to verdict pattern.

The two-instance evidence has the same shape in both verified rows. In
PvNP, the framework-primitive option supplies the grammar of the content
but not the substrate-specific non-currentization content, the SAU
non-descent route becomes circular through the translation theorem's
biconditional, and V-Differential elevation classifies the substrate
without turning that classification into a derivation. In RH, the same
two-instance \SigPattern\ signature appears with the corresponding
trace-decay and Selberg-trace inputs. This is evidence for the
recognition-mode reading of TSO substrates, not a universal theorem about
all future tracks.

\begin{table}[H]
  \centering
  \caption{Cross-track verdict signature \SigPattern{} per substrate; two-instance verified, with universal lifting deferred to third-instance replication.}
  \label{tab:verdict-signature}
  \small
  \begin{tabularx}{\textwidth}{@{}>{\raggedright\arraybackslash}p{0.14\textwidth}>{\centering\arraybackslash}X>{\centering\arraybackslash}X>{\centering\arraybackslash}X@{}}
    \toprule
    Substrate & Option (i): framework-primitive & Option (ii): SAU non-descent & Option (iii): V-Differential elevation \\
    \midrule
    PvNP & verdict (ii) & verdict (iii) & verdict (ii) \\
    RH & verdict (ii) & verdict (iii) & verdict (ii) \\
    \bottomrule
  \end{tabularx}
\end{table}
 
The next subsection explains why the signature has this form. The
Closure-Content Thesis states the structural reason: in a
trace-state-only substrate, the load-bearing content needed for closure is
formed-layer closure content named as an accepted recognition source, not
content separately derived from the framework's primitives.

\subsection{Closure-content thesis}\label{sec:recognition:closure-content-thesis}

The Cross-Track Verdict Signature points to a semantic distinction. In a
trace-state-only substrate, the sweep does not find the load-bearing
source by descending through framework primitives. The load-bearing source
is structural content of the formed layer itself: the closure names it,
records its provenance, and imports it as accepted recognition under
Foundations I closure formation. The Closure-Content Thesis is the
meta-theoretic claim that this is not an accident of PvNP or RH, but the
closure shape of the trace-state-only column.

We use three predicates to state the thesis. A structural fact is
\emph{formed-layer closure content} when it is content supplied by the
formed closure under Foundations I closure formation, rather than a
free-standing theorem produced later. A structural fact is
\emph{derivable from framework primitives} when it can be obtained purely
from the framework's internal primitive machinery, without accepting it as
part of closure formation. A structural fact is \emph{named as an accepted
recognition source} when the landing explicitly records that fact as the
source imported under the closure-formation discipline. The thesis says
that the load-bearing structural content of a trace-state-only landing
satisfies the first and third predicates and fails the second.

Let
\[
  s_{c,X}:=\operatorname{LoadBearingFor}(c,X)
\]
be the load-bearing structural content attached to the formed closure
\(c\) and target \(X\).  The three predicates are
\[
\begin{aligned}
  \operatorname{FLC}(s,c),\qquad
  \operatorname{DFP}(s),\qquad
  \operatorname{NAS}(s,c).
\end{aligned}
\]
These abbreviate formed-layer closure content, derivability from
framework primitives, and named-as-accepted-source.
The thesis asserts the following guarded conjunction:
\[
  \mbox{\VDTSO}(c)
  \Longrightarrow
  \operatorname{FLC}(s_{c,X},c)\wedge
  \neg\operatorname{DFP}(s_{c,X})\wedge
  \operatorname{NAS}(s_{c,X},c).
\]
The load-bearing assignment \(s_{c,X}\) is opaque: the conjunction is
not forced by expanding the definition of \(s_{c,X}\).

\begin{definition}[Load-bearing predicate]\label{def:load-bearing}
For a structural content \(s\), a substrate \(c\), and a target \(X\),
the predicate \(\operatorname{LoadBearingFor}(s,c,X)\) marks \(s\) as
load-bearing for the closure of \(X\) on \(c\), that is, \(s\) is what
closure formation must accept as named recognition in order to land
\(X\).  The predicate is treated as primitive in the body.
\end{definition}

\begin{schema}[Closure-Content Thesis]\label{thm:closure-content-thesis}
Let \(c\) be a formed closure whose substrate lies in the \VDTSO, and let
\(X\) be its target statement. The load-bearing structural content needed
for \(c\) to close \(X\) is formed-layer closure content per Foundations I
and is named as an accepted recognition source. It is not separately
derivable from the currently available framework primitives. Equivalently,
every recognition-mode landing on a trace-state-only substrate names a
source imported as accepted recognition under closure formation, rather
than a source produced by the framework's internal machinery.
In symbols,
\[
  \operatorname{FLC}(s_{c,X},c)\wedge
  \neg\operatorname{DFP}(s_{c,X})\wedge
  \operatorname{NAS}(s_{c,X},c).
\]
\end{schema}\footnote{\raggedright Tracked in the formalization
harness as a conditional schema in
\texttt{SixBirdsMetaMath.MetaMath.\allowbreak
RecognitionMode.ClosureContentThesis.\allowbreak closure\_content\_thesis},
conditioned on the opaque load-bearing predicate; paper-proposal-grade;
see Appendix~\ref{app:formalization}.}

The displayed conjunction is read componentwise.  The first component
places \(s_{c,X}\) in the formed-layer closure content of \(c\); the
second rules out derivability from current framework primitives; the
third records \(s_{c,X}\) as a named accepted source for the landing.
This is why the thesis is tied to the load-bearing predicate rather
than to an arbitrary structural fact \(s\).

The three predicates in the conjunction are independent pieces of the
claim.  A structural source can be content of a formed closure without
being explicitly named; it can be named in a paper without being
load-bearing; and it can be derivable from primitives without being a
recognition source.  The thesis asserts the specific TSO pattern
\[
  \operatorname{FLC}\wedge\operatorname{NAS}\wedge\neg\operatorname{DFP}.
  \]
This is the right meta-claim because the witness in a trace-state-only
substrate is not exposed as a direct object of the problem.  The
load-bearing content is instead the recognition source that makes the
trace-state evidence function as closure content.  Foundations~I
closure formation \cite{Tsiokos2026EmergenceCalculus} is the discipline
under which such a source is imported as named recognition, and the CSL
source \cite{Tsiokos2026CSL} is the PvNP instance of that pattern.  The
thesis ties these two facts together: load-bearing TSO content is named
recognition content, not framework-primitive output.

In plain mathematical terms, ``imported as named recognition'' means that
the paper does not pretend the recognition source has been manufactured
inside the framework. The recognition source
\(\Gamma_{\text{csl-sat-hidden}}\) is imported as named recognition per
Foundations I closure formation from \PaperCSL, and the recognition source
\(\Gamma_{\text{SDTC}}^{\text{Selberg}}\) is imported as named recognition
per Foundations I closure formation from the SDTC and RH apparatus. The
closure records those sources as load-bearing structural content and then
uses the translation stage to land the standard target statement. This is
why the framework-primitive sweep option can supply grammar without
supplying the substrate-specific source itself.

The next subsection uses this thesis to state the v5 refinement of Attack
Foreclosure. The point of v5 is not that recognition mode escapes the v4
foreclosure of derivation routes. It is that recognition-mode landings are
structurally distinct from derivation routes: they import their
load-bearing source as named recognition under the SB closure assumption,
and so remain peer with derivation-mode landings rather than subordinate
to them.

\subsection{Attack Foreclosure v5}\label{sec:recognition:attack-foreclosure-v5}

The v5 refinement is the joint statement of the negative and positive
halves. The negative half remains \Cref{thm:attack-foreclosure-v4}:
derivation routes internal to the framework do not advance the target
without named external content. The positive half is supplied by the
recognition-mode path: for trace-state-only substrates, a formed closure
may land by importing its load-bearing source as named recognition under
the SB closure assumption. The refinement records that these are two
typed paths, not a collapse of one into the other.

The refinement uses three typed predicates.  A framework-internal move
\(m\) may be a derivation route, written
\(\operatorname{DerivationRoute}(m)\).  A recognition-mode landing for
\((c,X)\) is written \(l:\operatorname{RMLanding}(c,X)\).  Its source
import condition is
\(\operatorname{ImportsAsNamedRecognition}(l,s_{c,X})\), where
\(s_{c,X}\) is the load-bearing content from
\Cref{thm:closure-content-thesis}.

\begin{schema}[Attack Foreclosure v5]\label{thm:attack-foreclosure-v5}
Let \(X\) be a target in framework scope, and let \(c\) be a formed
closure whose substrate lies in the \VDTSO\ under the SB closure
assumption.  Then the v5 refinement has the following three-conjunct
form:
\[
\begin{aligned}
  &\forall m.\,
    \operatorname{DerivationRoute}(m)
    \Longrightarrow
    \neg\operatorname{AdvancesWithoutExternal}(m,X),\\
  &\forall l:\operatorname{RMLanding}(c,X).\,
    \operatorname{ImportsAsNamedRecognition}(l,s_{c,X}),\\
  &\forall l:\operatorname{RMLanding}(c,X).\,
    \operatorname{LandsAtTheoremGrade}(c,X).
\end{aligned}
\]
The first conjunct is the v4 foreclosure conclusion restricted to
derivation routes.  The second says that recognition-mode landings
import the load-bearing source as named recognition.  The third says
that such landings reach theorem-grade under the SB closure assumption.
Thus v5 refines, does not replace, v4: derivation routes remain
foreclosed, while recognition-mode landings are structurally distinct
and peer with derivation-mode landings under the SB closure assumption.
\end{schema}\footnote{\raggedright Tracked in the formalization
harness as a conditional schema in
\texttt{SixBirdsMetaMath.MetaMath.\allowbreak
RecognitionMode.AttackForeclosureV5.\allowbreak attack\_foreclosure\_v5},
conditioned on the framework-scope guard composed with opaque
derivation-route, recognition-mode-landing, and
imports-as-named-recognition predicates; paper-proposal-grade; see
Appendix~\ref{app:formalization}.}

The case structure is explicit.  If a path is represented by a
framework-internal move \(m\) satisfying
\(\operatorname{DerivationRoute}(m)\), the first conjunct applies the
v4 foreclosure predicate to \(m\).  If the path is represented by a
recognition-mode landing \(l\), the second and third conjuncts record
its named-source import and theorem-grade landing.  The two readings
are not identified; that separation is the formal content behind the
``refines, does not replace'' phrasing.

The intended closure-path dichotomy is that an admissible path is either
a derivation route or a recognition-mode landing; there is no tertium
quid in the v5 taxonomy.  On the derivation-route side, v4 remains in
force: the move is still subject to the statement that it does not
advance \(X\) without external content.  On the recognition-mode side,
\Cref{thm:template-applicability} supplies the theorem-grade landing
under the SB closure assumption and \Cref{thm:closure-content-thesis}
identifies the load-bearing source as named recognition.  Thus
the recognition case is new relative to v4, but it is not a
counterexample to v4.  It is a different typed route whose content is
imported under closure formation rather than derived by a framework
internal move.

The phrase ``refines, does not replace'' is essential. Attack Foreclosure
v5 is not a claim that the v4 foreclosure theorem was too strong, nor a
claim that recognition mode refutes the v4 analysis. It says that v4
continues to hold for derivation routes, while the recognition-mode
path of Part~II is not a derivation route from framework primitives
alone. The
recognition-mode landing imports a named structural source with
provenance and audit discipline; it does not reclassify that source as
framework-internal derivation.

This is also the peer-path framing of Part II. Recognition-mode landings
are peer with derivation-mode landings under the SB closure assumption:
they are not weaker, not contingent on derivation-mode results, and not
subordinate to them. The difference is structural. Derivation mode earns a
landing by deriving the needed content through an admissible derivation
route. Recognition mode earns a landing by naming the load-bearing source
as accepted recognition content under closure formation.

This closes the Part II schema chain.
\Cref{thm:template-applicability} states the template applicability
claim, \Cref{thm:verdict-signature} records the two-instance Stage-6
signature, \Cref{thm:closure-content-thesis} identifies the
closure-content source, and \Cref{thm:attack-foreclosure-v5} gives
the joint v5 refinement. All four claims are \PPGrade\ and \TwoInstance; the
validation rows in \Cref{sec:validation} record the PvNP and RH
instances that motivate the chain, while theorem-grade universal
lifting remains outside the scope of this section.

\begin{table}[H]
  \centering
  \caption{Part~II recognition-mode summary: the four paper-proposal-grade T-rows.}
  \label{tab:recognition:summary}
  \scriptsize
  \renewcommand{\arraystretch}{1.08}
  \begin{tabularx}{\textwidth}{@{}>{\raggedright\arraybackslash}p{0.06\textwidth}>{\raggedright\arraybackslash}X>{\raggedright\arraybackslash}p{0.14\textwidth}>{\raggedright\arraybackslash}p{0.20\textwidth}>{\raggedright\arraybackslash}p{0.20\textwidth}@{}}
    \toprule
    T\# & Statement & Lean coverage & Opacity guard & Status (per M5 \S10.6) \\
    \midrule
    T7 & The seven-stage recognition-mode landing template applies to formed TSO closures under SB closure assumption. & obligation & template-applies; opaque substrate-typing predicates & \PPGrade; \TwoInstance \\
    T8 & The Stage-6 sweep has the two-instance \SigPattern{} cross-track signature on PvNP and RH. & obligation & template-applies; opaque sweep-verdict & \PPGrade; \TwoInstance \\
    T9 & TSO load-bearing closure content is imported as named recognition, not derived from framework primitives. & obligation & opaque load-bearing predicate & \PPGrade; \TwoInstance \\
    T10 & Attack Foreclosure v5 refines, does not replace, v4: derivation routes remain foreclosed while recognition-mode landings are peer theorem-grade paths. & obligation & framework-scope; opaque recognition predicates & \PPGrade; \TwoInstance \\
    \bottomrule
  \end{tabularx}
\end{table}
  \section{Cross-track validation}\label{sec:validation}

\subsection{Introduction}\label{sec:validation:intro}

This section validates rather than derives.  The substrate papers are
used here as cited evidence for the meta-mathematical architecture
developed in \Cref{sec:recognition}: the column split, the
\SevenStage, the two-instance Stage-6 signature, and the v5 refinement.
We do not restate any substrate theorem, and we do not claim derivation
responsibility for any Clay-class result in this paper.

The conditional-closure qualifier is binding throughout the validation
rows. The first row uses the conditional closure of Navier--Stokes
regularity landed by \PaperNS\ under the SB closure assumption of
Foundations I \cite{Tsiokos2026EmergenceCalculus,Tsiokos2026NavierStokes};
the official problem statement is \cite{Fefferman2000NavierStokes},
and recent landmark work on the regularity question includes
\cite{Tao2016NSAveragedBlowup,BuckmasterVicol2019NS}.
The second row uses the
conditional closure of the Riemann hypothesis landed by \PaperRH\ under
the SB closure assumption of Foundations I, together with the SDTC
apparatus of \PaperSDTC\ \cite{Tsiokos2026RH,Tsiokos2026SDTC}; the
official problem statement is \cite{Bombieri2000Riemann}, with
\cite{Conrey2003RH,IwaniecKowalski2004ANT} as standard surveys.
The third row uses the conditional closure of \(P\) versus \(NP\)
landed by \PaperPvNP\ under the SB closure assumption of Foundations~I,
together with the CSL recognition source of \PaperCSL\
\cite{Tsiokos2026PvNP,Tsiokos2026CSL}; the official problem statement
is \cite{Cook2000PvsNP}, with
\cite{Cook1971NPCompleteness,Karp1972NPComplete21} for the foundational
{NP}-completeness results. These citations locate the
substrate evidence; they are not re-proofs inside meta-math.

The validation rows also exhibit the WCE--TSO split from
\Cref{sec:recognition:column-classification}. The NS row is the
witness-content-exposing row: its landing proceeds through
substrate-specific apparatus rather than the recognition-mode template.
The RH and PvNP rows are trace-state-only rows: their load-bearing sources
are imported as named recognition per Foundations I closure formation,
respectively \(\Gamma_{\text{SDTC}}^{\text{Selberg}}\) and
\(\Gamma_{\text{csl-sat-hidden}}\). This is the cross-track evidence for
the Part~II recognition-mode chain, not a claim that this paper
derives those recognition sources.

The reading order is correspondingly simple. The next subsection records
the NS witness-content-exposing row. The following two subsections record
the RH and PvNP trace-state-only rows, including their named recognition
sources. A final validation summary collects the column placement and the
role each row plays in supporting the Part II meta-claims.

\subsection{Navier--Stokes: witness-content-exposing row}\label{sec:validation:ns}

The conditional closure of Navier--Stokes regularity landed by \PaperNS\
under the SB closure assumption of Foundations I is the validation
section's witness-content-exposing row
\cite{Tsiokos2026NavierStokes}. This row is included by citation only:
the substrate theorem is not restated here, and this paper does not claim
derivation responsibility for the Navier--Stokes landing.

Its role is structural. The NS landing sits in the \VDWCE, where the
closure route proceeds through substrate-specific apparatus rather than
through the \SevenStage. In the language of \Cref{sec:apparatus}, the
landing is mediated by the layer-dissolving membrane endpoint and the
associated substrate apparatus, not by importing a trace-state-only
recognition source.

This makes the NS row a useful boundary case for Part II. It shows by
example that the recognition-mode landing template does not exhaust
Clay-class closures. The Part~II schemas describe the
trace-state-only recognition path; the NS row witnesses the other
column, where closure is witness-content-exposing and
substrate-specific.

The contrast with the next two rows is the validation point. RH and PvNP
are trace-state-only rows whose load-bearing sources are imported as
named recognition per Foundations I closure formation. NS is not. It is
therefore not evidence for the Part~II recognition-mode chain, but
it is evidence that the V-Differential column distinction is doing
real classification work.

\subsection{Riemann hypothesis: trace-state-only row}\label{sec:validation:rh}

The conditional closure of the Riemann hypothesis landed by \PaperRH\
under the SB closure assumption of Foundations I is the validation
section's first trace-state-only row, with the SDTC apparatus supplied by
\PaperSDTC\ \cite{Tsiokos2026RH,Tsiokos2026SDTC}. This row is cited as
substrate evidence only: the substrate theorem is not restated here, and
this paper does not claim derivation responsibility for the RH landing.

The recognition source for this row is
\(\Gamma_{\text{SDTC}}^{\text{Selberg}}\), imported as named recognition
per Foundations I closure formation from \PaperRH\ and \PaperSDTC. It is
not derived from framework primitives in this paper. In the column
language of Part II, the RH landing is trace-state-only because the
load-bearing content is carried by trace-state evidence and must be named
as an accepted recognition source.

This makes RH one of the two cross-track instances supporting the Part II
recognition-mode chain. It is positive evidence for the
Recognition-Mode Template Applicability Theorem because the landing
follows the \SevenStage. It is positive evidence for the Cross-Track
Verdict Signature because the Stage-6 sweep exhibits the two-instance
\SigPattern\ pattern. It is positive evidence for the Closure-Content
Thesis because the load-bearing source is accepted as formed-layer
closure content, and it supports the v5 refinement because the landing is
recognition-mode rather than a framework-internal derivation route.

The next row records the PvNP instance. Together, RH and PvNP are the two
trace-state-only rows behind the \TwoInstance\ qualifier used throughout
Part II.

\subsection{P versus NP: trace-state-only row}\label{sec:validation:pvnp}

The conditional closure of \(P\) versus \(NP\) landed by \PaperPvNP\
under the SB closure assumption of Foundations I is the validation
section's second trace-state-only row, with the CSL recognition source
supplied by \PaperCSL\ \cite{Tsiokos2026PvNP,Tsiokos2026CSL}. As in the
other validation rows, this is citation-only evidence: the substrate
theorem is not restated here, and this paper does not claim derivation
responsibility for the PvNP landing.

The recognition source for this row is
\(\Gamma_{\text{csl-sat-hidden}}\), imported as named recognition per
Foundations I closure formation from \PaperCSL\ in the context of the
\PaperPvNP\ landing. It is not derived from framework primitives in this
paper. In the column language of Part II, PvNP is trace-state-only because
the closure-relevant hiddenness content is carried as trace-state evidence
and must be named as accepted recognition content.

PvNP is the second of the two cross-track instances supporting the Part II
recognition-mode chain. It is positive evidence for the
Recognition-Mode Template Applicability Theorem because the landing uses
the \SevenStage. It is the other instance for the Cross-Track Verdict
Signature: the Stage-6 sweep exhibits the same two-instance
\SigPattern\ pattern as RH. It supports the Closure-Content Thesis
because the load-bearing hiddenness source is imported as formed-layer
closure content, and it supports the v5 refinement because the landing is
recognition-mode rather than a framework-internal derivation route.

The PvNP paper also carries the Class-B Irremovability gate-context used
later in Part III's named-functorial-gates discussion. That context
returns in \Cref{sec:calibration}; here it remains a forward pointer, not
a new proof claim. For this section, the role of PvNP is the
second trace-state-only instance paired with RH.

\begin{table}[H]
  \centering
  \caption{Cross-track validation summary for the three landed conditional Clay-class substrates.}
  \label{tab:cross-track-validation}
  \scriptsize
  \renewcommand{\arraystretch}{1.08}
  \begin{tabularx}{\textwidth}{@{}>{\raggedright\arraybackslash}p{0.08\textwidth}>{\raggedright\arraybackslash}p{0.13\textwidth}>{\raggedright\arraybackslash}X>{\raggedright\arraybackslash}p{0.24\textwidth}>{\raggedright\arraybackslash}p{0.15\textwidth}@{}}
    \toprule
    Substrate & V-Differential column & Landing route & Recognition source (when applicable) & SB closure assumption invocation \\
    \midrule
    NS & \VDWCE & Substrate-specific apparatus via the layer-dissolving membrane endpoint; not a recognition-mode landing. & Not applicable. & Conditional closure under Foundations~I. \\
    RH & \VDTSO & Recognition-mode landing through the \SevenStage. & \(\Gamma_{\text{SDTC}}^{\text{Selberg}}\), imported as named recognition per Foundations~I closure formation. & Conditional closure under Foundations~I. \\
    PvNP & \VDTSO & Recognition-mode landing through the \SevenStage. & \(\Gamma_{\text{csl-sat-hidden}}\), imported as named recognition per Foundations~I closure formation. & Conditional closure under Foundations~I. \\
    \bottomrule
  \end{tabularx}
\end{table}
  \section{Part III --- Calibration Family}\label{sec:calibration}

\subsection{Introduction}\label{sec:calibration:intro}

Parts~I and II separated two questions.  Part~I classified what happens
when the framework tries to advance a target conjecture by its own
carrier moves.  Part~II classified a second path, recognition-mode
landing, in which the load-bearing source is named and imported under
closure formation rather than produced by an internal derivation.  Part~III
supplies the apparatus that lets those two accounts be stated uniformly
across substrates: the calibration family, its bridges, its coverage
geometry, and the typed gates that mark where framework-native reach stops.

A \emph{calibration} is a typed object recording a named detector class
together with the kind of carrier it operates on.  The point of the
calibration language is not to collect analogies, but to type them.  A
low-degree detector, a mutual-information detector, a stable-local
detector, and an algebraic certificate-degree detector do not live in a
single undifferentiated method space.  They have different native carrier
types, different descent invariants, and different bridge requirements.
The calibration family records this information as a seed catalog rather
than as an informal list of techniques.  Later subsections make this
precise by listing the seed family, imposing the \fourAxis, and separating
new calibrations from rebadged copies.

The original Non-Descending Translation meta-theorem motivated the
calibration-family apparatus historically; its translation/handoff content
has been folded into Part~III as the structural background for the named
functorial gates of \Cref{sec:calibration:functorial-gates}.  The
vobs-arc evidence from \PaperPvNP\ and the 11+5 calibration sprint
supplied the concrete catalog from which the three gates of
\Cref{cat:named-functorial-gates} were extracted.  In this
paper that material is background: it is
not restated as a separate theorem, and the formal body begins with the
calibration-family structure itself.

The seven rows of Part~III then proceed as follows.  The
Calibration-Family Structure theorem introduces the seeded calibration
family.  The Rebadge Discipline states the four-axis criterion for
when a candidate calibration is genuinely new rather than a relabeled
clone.  The Bridges-as-Theorems result treats a
\TypedBridgeCalibration{} as either a basis-change dictionary or a
substantive theorem, and the Bridge Type-Direction theorem records the
directed source-to-target asymmetry of such bridges.  The
Coverage-Family Geometry describes the covered carrier regions, the
uncovered regions, and the named typed holes between them.  The
Framework-Type Self-Characterization compares the framework-native
type reached on the \(P\) versus \(NP\) track with the standard Clay
target type.  Finally, the Named Functorial Gates Catalog records the
third \TDS{}: typed functorial gates with source, target, direction,
and stated theorem obligation.

The reading order is therefore:
\Cref{sec:calibration:structure} gives the seed family,
\Cref{sec:calibration:rebadge} separates genuine additions from
rebadges,
\Cref{sec:calibration:bridges-as-theorems,sec:calibration:bridge-type-direction}
analyze calibration-side bridges,
\Cref{sec:calibration:coverage-family-geometry} gives the coverage
geometry, \Cref{sec:calibration:framework-type-self-char} states the
framework-type comparison, and
\Cref{sec:calibration:functorial-gates} names the functorial gates
that separate the framework-type from the standard target type.

\subsection{Calibration-family structure}\label{sec:calibration:structure}

The seed catalog is the typed form of the Part~III apparatus.  It records
that non-descending translation is not a single instrument, but a family
of typed constructions, each attached to an obstruction shape and to the
native carrier on which that obstruction is visible.  The family is a
seed rather than a closure condition: additional calibrations may be
added when they pass the typing and earning-their-place tests developed
below.

\begin{definition}[Calibration]\label{def:calibration}
A calibration is a record consisting of a name, a detector class, and a
native carrier type.  The name identifies the calibration in the family;
the detector class states what kind of obstruction it reads; and the
native carrier type states the mathematical carrier in which that
detector is meaningful.
Symbolically, the record shape is
\[
  \mathsf{Calibration}
  =
  \{\mathsf{name}:\mathsf{String},\
    \mathsf{detectorClass}:\mathsf{String},\
    \mathsf{nativeCarrierType}:\mathsf{String}\}.
\]
\end{definition}

The seed family is a finite list
\[
  \mathcal C_0:\mathsf{List}(\mathsf{Calibration}),
  \qquad |\mathcal C_0|=13,
\]
whose entries are displayed in \Cref{tab:calibration:seed}.  The list is
finite, but the record type is not closed: a later candidate is another
record of the same type, to be tested by the rebadge discipline below.

\begin{theorem}[Calibration-Family Structure]\label{thm:calibration-family-structure}
There is a non-descending typed calibration family with a thirteen-entry
seed catalog.  Each seed entry is a calibration in the sense of
\Cref{def:calibration}, and the seed is open-ended: the framework may
adjoin further typed calibrations without changing the fact that the
initial catalog is a typed family rather than a single construction.
In the formal finite statement,
\[
  |\mathcal C_0|\ge 13
  \quad\text{and}\quad
  \forall c:\mathsf{Calibration},\ c.\mathsf{name}\ne\varnothing
  \Rightarrow c\ \text{is an admissible candidate record}.
\]
\end{theorem}

\begin{proof}
Define \(\mathcal C_0\) to be the thirteen records displayed in
\Cref{tab:calibration:seed}.  Each entry has the three fields of
\Cref{def:calibration}: a name, a detector class, and a native carrier
type.  Thus every listed row is a well-typed calibration record.  By
direct enumeration of the list, there are thirteen entries, so
\[
  |\mathcal C_0|=13,
  \qquad\text{hence}\qquad
  |\mathcal C_0|\ge 13 .
\]

The non-descending character of the seed is also read from the typed
record data.  No row is introduced as a quotient, specialization, or
descendant of another row; each has its own detector class and native
carrier type.  This is the structural, seed-level meaning needed here:
the family is a collection of typed calibration records rather than a
single instrument presented under thirteen names.  The sharper
``not merely renamed'' test is supplied by the Rebadge Discipline in
\Cref{thm:rebadge-discipline}.

Finally, the family is open-ended because the type
\(\mathsf{Calibration}\) is a record type, not a closed enumeration.
Given any further record \(c\) with nonempty name, one may form the
extended list \(\mathcal C_0\cup\{c\}\) and then test it by the
subsequent profile and bridge disciplines.  The theorem asserts the
existence and typed structure of the seed, not a claim that the seed is
the universe of all calibrations.
\end{proof}\footnote{\raggedright This theorem is verified in Lean as
\texttt{SixBirdsMetaMath.MetaMath.\allowbreak
CalibrationFamily.Structure.\allowbreak calibration\_family\_structure};
see Appendix~\ref{app:formalization}.}

The evidence base for the seed is the PvNP track calibration sprint (11
valid calibrations plus 5 honest retracts).  The sprint belongs to the
\PaperPvNP{} evidence base \cite{Tsiokos2026PvNP}; the theorem here
records the structural catalog, not a new substrate closure claim.
Background literature for the individual detector classes includes
\cite{Lovasz1978Kneser} for the topological/cohomological line,
\cite{MulmuleySohoni2001GCT} for the arithmetic-complexity / GCT line,
\cite{Krajicek2019ProofComplexity} for proof complexity and
Nullstellensatz certificates, and
\cite{AlonKrivelevichSudakov1998PlantedClique,Kearns1998SQ} for the
planted-clique and statistical-query backgrounds adjacent to the
hiddenness-side calibrations.

\begin{table}[H]
  \centering
  \caption{Thirteen-entry calibration seed catalog for \Cref{thm:calibration-family-structure}.}
  \label{tab:calibration:seed}
  \scriptsize
  \renewcommand{\arraystretch}{1.06}
  \begin{tabularx}{\textwidth}{@{}>{\raggedright\arraybackslash}p{0.10\textwidth}>{\raggedright\arraybackslash}p{0.26\textwidth}>{\raggedright\arraybackslash}X>{\raggedright\arraybackslash}p{0.22\textwidth}>{\raggedright\arraybackslash}p{0.12\textwidth}@{}}
    \toprule
    Symbol & Name & Native carrier type & Detector class & Source \\
    \midrule
    \(\Cvobs\) & volumetric observables & \(AC^0\) / \(AC^0[p]\)-adjacent carriers & low-degree polynomial detector & \PaperPvNP \\
    \(\CMI\) & information calibration & bounded-information packages & mutual information / KL divergence & \PaperPvNP \\
    \(\CRE\) & stability calibration & random-restriction / coupling-stable carriers & stable, local, overlap algorithms & \PaperPvNP \\
    \(\CCh\) & cohomological calibration & cohomology-lifted carriers with Lovasz/Kneser reach & topological obstruction & \PaperPvNP \\
    \(\CQamp\) & amplitude calibration & black-box and query carriers & quantum amplitude / query detector & \PaperPvNP \\
    \(\CIalg\) & algebraic calibration & proof-system and certificate-degree carriers & ideal / Nullstellensatz detector & \PaperAOR\ \cite{Tsiokos2026AOR} \\
    \(\CRreal\) & realizability calibration & lambda-calculus / typed-program carriers & realizability detector & \PaperPvNP \\
    \(\COCstar\) & \(C^*\)-operator calibration & operator-algebraic carriers & \(C^*\)-operator detector & \PaperPvNP \\
    \(\CFexistsSO\) & existential-second-order calibration & logical / locality carriers & descriptive-complexity detector & \PaperPvNP \\
    \(\CSHO\) & higher-order strategy calibration & higher-order game carriers & game-semantic strategy detector & \PaperPvNP \\
    \(\CPPH\) & persistent-homology calibration & scale-dependent topological carriers & persistent-homology detector & \PaperPvNP \\
    \(\CAVNP\) & arithmetic-complexity calibration & GCT / polynomial-family carriers & arithmetic-complexity / VNP detector & \PaperPvNP \\
    \(\CMmot\) & computational motive (META) & cross-calibration transport carriers & motive-level transport detector & \PaperPvNP \\
    \bottomrule
  \end{tabularx}
\end{table}
 
The next issue is duplication.  A large seed catalog is useful only if it
distinguishes genuinely new calibrations from renamed copies of existing
ones.  The Rebadge Discipline of \Cref{thm:rebadge-discipline} supplies that
discriminator by testing candidates along the \fourAxis{}: algebraic
identities, detector class, descent invariants, and carrier
orientation.

\subsection{Rebadge discipline}\label{sec:calibration:rebadge}

The calibration family is deliberately open-ended, but open-endedness
requires a guard against relabeling.  A candidate should not count as a
new calibration merely because it has a new name.  It must differ from
the existing seed family in the structural information by which a
calibration earns its place.

\begin{definition}[Calibration four-axis profile]\label{def:four-axis-profile}
The \fourAxis{} of a calibration consists of its algebraic identities,
detector class, descent invariants, and carrier orientation.  A candidate
calibration is a rebadge of an existing calibration when the two profiles
match on all four axes.  It is genuinely new relative to a family when it
is not a rebadge of any existing family member.
Thus a profile is the product
\[
  \mathsf{Profile}
  =
  \mathsf{AlgId}\times
  \mathsf{DetectorClass}\times
  \mathsf{DescentInv}\times
  \mathsf{CarrierOrient}.
\]
For profiles \(p\) and \(q\),
\[
  \mathsf{Rebadge}(p,q)
  \iff
  p_{\mathsf{alg}}=q_{\mathsf{alg}}\wedge
  p_{\mathsf{det}}=q_{\mathsf{det}}\wedge
  p_{\mathsf{desc}}=q_{\mathsf{desc}}\wedge
  p_{\mathsf{orient}}=q_{\mathsf{orient}},
\]
and
\[
  \mathsf{GenuineNew}(p,\mathcal F)
  \iff
  \forall q\in\mathcal F,\ \neg\mathsf{Rebadge}(p,q).
\]
\end{definition}

\begin{theorem}[Rebadge Discipline]\label{thm:rebadge-discipline}
Within a calibration family, a candidate calibration is genuinely new if
and only if, for every existing family member, it differs on at least one
axis of the \fourAxis{}.  Equivalently, a candidate matching an existing
member on algebraic identities, detector class, descent invariants, and
carrier orientation is a rebadge; on attempted earning-its-place, it
collapses to the existing member's content.
Equivalently, for a candidate profile \(p\) and family \(\mathcal F\),
\[
  \mathsf{GenuineNew}(p,\mathcal F)
  \iff
  \forall q\in\mathcal F,\,
  p_{\mathsf{alg}}\ne q_{\mathsf{alg}}\vee
  p_{\mathsf{det}}\ne q_{\mathsf{det}}\vee
  p_{\mathsf{desc}}\ne q_{\mathsf{desc}}\vee
  p_{\mathsf{orient}}\ne q_{\mathsf{orient}}.
\]
\end{theorem}

\begin{proof}
Fix a candidate profile \(p\) and a family \(\mathcal F\) of existing
profiles.  By definition,
\(\mathsf{Rebadge}(p,q)\) is the conjunction of the four componentwise
equalities
\[
  p_{\mathsf{alg}}=q_{\mathsf{alg}},\quad
  p_{\mathsf{det}}=q_{\mathsf{det}},\quad
  p_{\mathsf{desc}}=q_{\mathsf{desc}},\quad
  p_{\mathsf{orient}}=q_{\mathsf{orient}}.
\]
Thus \(\mathsf{Rebadge}(p,q)\) is exactly equality of the four-axis
profiles.  Since the axes in this abstraction are string-valued
coordinates, equality and inequality on each coordinate are decidable
componentwise.

Assume first that \(p\) is genuinely new relative to \(\mathcal F\).
Let \(q\in\mathcal F\).  If all four axes of \(p\) and \(q\) were
equal, then \(\mathsf{Rebadge}(p,q)\) would hold, contradicting
\(\mathsf{GenuineNew}(p,\mathcal F)\).  Therefore at least one of the
four coordinate equalities fails, which is precisely the displayed
four-way disjunction.

Conversely, assume that for every \(q\in\mathcal F\) at least one of
the four axes differs.  If \(p\) were a rebadge of some \(q\in\mathcal F\),
then all four axes would be equal.  This contradicts the assumed
four-way disjunction for that same \(q\).  Hence no existing member is
a rebadge target for \(p\), so \(p\) is genuinely new.
\end{proof}\footnote{\raggedright This theorem is verified in Lean as
\texttt{SixBirdsMetaMath.MetaMath.\allowbreak
CalibrationFamily.RebadgeDiscipline.\allowbreak rebadge\_discipline};
see Appendix~\ref{app:formalization}.}

This theorem is the first place where Part~III turns the catalog into a
discipline.  The next subsection asks what it takes to move information
between two genuinely typed calibrations: such a bridge is not free
transport, but a basis-change dictionary or a substantive theorem.

\subsection{Bridges as theorems}\label{sec:calibration:bridges-as-theorems}

Once calibrations are typed, a bridge between them is no longer an
informal analogy.  It must say what is being transported, from which
calibration to which calibration, and why the transport is legitimate.
The Bridges-as-Theorems principle is the calibration-family answer: a
bridge is either a genuine basis-change identity, or it is a theorem with
content.

\begin{definition}[Calibration-side typed bridge]\label{def:calibration-side-typed-bridge}
A \TypedBridgeCalibration{} is a directed typed transport between two
calibrations.  It has a source calibration, a target calibration, and a
bridge kind indicating whether the transport is a basis-change identity
or substantive theorem.\footnote{This is the calibration-side bridge of
Part~III.  It is distinct from the carrier-side typed bridge of
\Cref{def:carrier-side-typed-bridge}, which appears in the Bridge
Impossibility Corollary of Part~I.}
The record form is
\[
  B=\{\mathsf{source}:\mathsf{Calibration},\
      \mathsf{target}:\mathsf{Calibration},\
      \mathsf{kind}:\mathsf{BridgeKind}\}.
\]
\end{definition}

\begin{definition}[Calibration-bridge kind dichotomy]\label{def:bridge-kind}
The kind of a \TypedBridgeCalibration{} has exactly two possibilities:
a basis-change kind, which identifies two presentations of the same
detector-and-carrier content, and a substantive kind, which asserts a
nontrivial theorem relating structurally different calibration
invariants.
As a two-constructor type,
\[
  \mathsf{BridgeKind}
  ::= \mathsf{basisChange}\mid \mathsf{substantive}.
\]
\end{definition}

\begin{theorem}[Bridges-as-Theorems]\label{thm:bridges-as-theorems}
Every calibration-side typed bridge between calibrations is either a
basis-change identity or a substantive theorem.  Free amalgamation
between calibrations, in which two calibrations are merely placed side by
side without a typed relation between them, is not a bridge.
In symbols,
\[
  \forall B:\TypedBridgeCalibration,\quad
  \mathsf{kind}(B)=\mathsf{basisChange}\vee
  \mathsf{kind}(B)=\mathsf{substantive}.
\]
\end{theorem}

\begin{proof}
Let \(B\) be a calibration-side typed bridge.  By definition,
\(\mathsf{kind}(B)\) is an element of the two-constructor type
\(\mathsf{BridgeKind}\).  There are therefore exactly two cases.

If \(\mathsf{kind}(B)=\mathsf{basisChange}\), then \(B\) is a typed
dictionary between equivalent presentations of the same calibration
content, and the left disjunct holds.  If
\(\mathsf{kind}(B)=\mathsf{substantive}\), then \(B\) is typed theorem
content relating its source and target calibrations, and the right
disjunct holds.  Since the bridge-kind type has no third constructor,
there is no third bridge case.  Free amalgamation is excluded because it
does not supply a value of \(\mathsf{BridgeKind}\) at all.
\end{proof}\footnote{\raggedright This theorem is verified in Lean as
\texttt{SixBirdsMetaMath.MetaMath.\allowbreak
CalibrationFamily.BridgesAsTheorems.\allowbreak bridges\_as\_theorems};
see Appendix~\ref{app:formalization}.}

This dichotomy is still not the whole story.  Even when a bridge exists,
its direction matters: a bridge from one calibration to another need not
provide the reverse bridge.  The next subsection isolates that
type-direction asymmetry.

\subsection{Bridge type-direction}\label{sec:calibration:bridge-type-direction}

Calibration-side bridges carry direction as part of their type.  A
transport from one calibration to another has a source and a target; it
is not automatically a transport in the reverse direction.  This matters
whenever the host theorem has its own direction, such as worst-case to
average-case, nonuniform to uniform, or statistical to algorithmic.

\begin{definition}[Calibration-bridge reversal]\label{def:calibration-bridge-reversal}
Given a \TypedBridgeCalibration{} \(B\) with source \(S\) and target
\(T\), its formal reversal is the bridge record obtained by swapping the
source and target fields.  Symmetry is the stronger condition that the
original bridge and its reversal agree as typed bridge data.
Writing \(B:S\to T\), the reversed record \(B^{\mathrm{op}}\) has type
\[
  B^{\mathrm{op}}:T\to S,
  \qquad
  \mathsf{kind}(B^{\mathrm{op}})=\mathsf{kind}(B).
\]
\end{definition}

\begin{theorem}[Bridge Type-Direction]\label{thm:bridge-type-direction}
Calibration-side typed bridges carry a directed source-to-target type
signature.  For a bridge \(B\) with source \(S\) and target \(T\), the
bridge is symmetric with its reversal if and only if \(S=T\).  In
particular, when \(S\ne T\), the existence of a typed bridge \(S\to T\)
does not by itself supply a usable typed bridge \(T\to S\).
Formally,
\[
  B=B^{\mathrm{op}}
  \iff
  \mathsf{source}(B)=\mathsf{target}(B).
\]
\end{theorem}

\begin{proof}
Let \(B\) have source \(S\), target \(T\), and bridge kind \(k\).  By
definition,
\[
  B=(S,T,k),
  \qquad
  B^{\mathrm{op}}=(T,S,k).
\]
Suppose first that \(B=B^{\mathrm{op}}\).  Comparing the source fields
of these two bridge records gives \(S=T\), which is
\(\mathsf{source}(B)=\mathsf{target}(B)\).

Conversely, suppose \(S=T\).  Substituting \(T\) for \(S\) in the
record \((S,T,k)\), both \(B\) and \(B^{\mathrm{op}}\) become the same
record \((T,T,k)\).  Hence \(B=B^{\mathrm{op}}\).  The equivalence
shows that direction is encoded in the source and target fields: a
reverse record is always definable, but equality with the original is
available only in the endo-bridge case.
\end{proof}\footnote{\raggedright This theorem is verified in Lean as
\texttt{SixBirdsMetaMath.MetaMath.\allowbreak
CalibrationFamily.BridgeTypeDirection.\allowbreak bridge\_type\_direction};
see Appendix~\ref{app:formalization}.}

The next subsection uses this directionality in the coverage accounting.
Some carrier regions are reached by the calibration family, some remain
unreached, and the directed failures between them are recorded as typed
holes rather than silently filled by reverse transport.

\subsection{Coverage-family geometry}\label{sec:calibration:coverage-family-geometry}

The calibration family covers regions, not just named techniques.  Its
geometry is therefore not the same thing as a union of complexity classes:
some calibrations overlap on a shared carrier region, some regions remain
outside the present family coverage, and some gaps are important enough
to be named as typed holes.  This is the coverage meaning of the
calibration apparatus.

Let \(\operatorname{Cov}(c)\) be the native carrier-region coverage of a
calibration \(c\).  The family coverage is the finite union
\[
  \operatorname{Cov}(\mathcal C_0)
  =
  \bigcup_{c\in\mathcal C_0}\operatorname{Cov}(c).
\]
For the seeded family, the covered and uncovered region sets are recorded
as
\[
\begin{aligned}
  R_{\mathrm{cov}}
  &=
  \{\mathsf{boundedDepth},\mathsf{lowDegree},\mathsf{stableLocal},
    \mathsf{query},\mathsf{topological},\mathsf{lowNSDegree}\},\\
  R_{\mathrm{gap}}
  &=
  \{\mathsf{NC^1},\mathsf{NC^2},\mathsf{P/poly},
    \mathsf{adaptivePolyTime}\}.
\end{aligned}
\]
The standard references for the uncovered regions are
\cite{Cook1985NCTaxonomy} for the \(\mathsf{NC}\) hierarchy and
\cite{KarpLipton1982Advice} for \(\mathsf{P/poly}\).
The geometric assertions are \(R_{\mathrm{cov}}\subseteq
\operatorname{Cov}(\mathcal C_0)\) and
\[
  R_{\mathrm{gap}}\cap\operatorname{Cov}(\mathcal C_0)=\varnothing.
\]
The named typed holes are the records
\[
  \xiNCone:\TypedHole,
  \qquad
  \operatorname{region}(\xiNCone)=\mathsf{NC^1},
\]
and
\[
  \xiAdaptive:\TypedHole,
  \qquad
  \operatorname{region}(\xiAdaptive)=\mathsf{adaptivePolyTime}.
\]

\begin{table}[H]
  \centering
  \caption{Covered and uncovered carrier regions, with named typed holes, for \Cref{thm:coverage-family-geometry}.}
  \label{tab:coverage-regions}
  \small
  \begin{tabularx}{\textwidth}{@{}>{\raggedright\arraybackslash}p{0.30\textwidth}>{\raggedright\arraybackslash}X>{\raggedright\arraybackslash}p{0.22\textwidth}@{}}
    \toprule
    Region & Membership & Status \\
    \midrule
    bounded-depth & in family coverage & covered \\
    low-degree & in family coverage & covered \\
    stable-local & in family coverage & covered \\
    query & in family coverage & covered \\
    topological & in family coverage & covered \\
    low-Nullstellensatz-degree & in family coverage & covered \\
    \textsc{NC}\(^{1}\) & outside family coverage & uncovered \\
    \textsc{NC}\(^{2}\) & outside family coverage & uncovered \\
    \textsc{P/poly} & outside family coverage & uncovered \\
    adaptive polynomial time & outside family coverage & uncovered \\
    \(\xiNCone\) & names the \textsc{NC}\(^{1}\) coverage gap & \TypedHole{} \\
    \(\xiAdaptive\) & names the adaptive-trace coverage gap & \TypedHole{} \\
    \bottomrule
  \end{tabularx}
\end{table}
 
\begin{theorem}[Coverage-Family Geometry]\label{thm:coverage-family-geometry}
The calibration-family coverage union has the following typed geometry.
The bounded-depth, low-degree, stable-local, query, topological, and
low-Nullstellensatz-degree carrier regions lie inside the family
coverage.  The \textsc{NC}\(^{1}\), \textsc{NC}\(^{2}\),
\textsc{P/poly}, and adaptive polynomial-time regions lie outside it.
Moreover, \(\xiNCone\) and \(\xiAdaptive\) are named typed coverage holes,
and adding a calibration with empty native coverage does not enlarge the
family coverage.
The non-monotone clause is
\[
  \operatorname{Cov}(c)=\varnothing
  \Longrightarrow
  \operatorname{Cov}(\mathcal F\cup\{c\})
  =
  \operatorname{Cov}(\mathcal F).
\]
\end{theorem}

\begin{proof}
The proof is a composition of the four finite-geometry lemmas attached
to the seeded family.

First, applying the covered-regions lemma gives
\[
  \mathsf{boundedDepth},\mathsf{lowDegree},\mathsf{stableLocal},
  \mathsf{query},\mathsf{topological},\mathsf{lowNSDegree}
  \in \operatorname{Cov}(\mathcal C_0).
\]
These are the six inclusions in the theorem.  Second, applying the
uncovered-regions lemma gives
\[
  \mathsf{NC^1},\mathsf{NC^2},\mathsf{P/poly},
  \mathsf{adaptivePolyTime}
  \notin \operatorname{Cov}(\mathcal C_0),
\]
which proves the four exclusions.  In particular, the two named typed
holes have uncovered regions:
\[
\begin{aligned}
  \operatorname{region}(\xiNCone)&=\mathsf{NC^1},\\
  \operatorname{region}(\xiAdaptive)&=\mathsf{adaptivePolyTime}.
\end{aligned}
\]
Thus both named typed holes are genuine.

Third, the bounded-depth-overlap lemma supplies two distinct seed
calibrations, \(\Cvobs\) and \(\CMI\), whose native coverage both
contain the bounded-depth region.  This proves the shared-subregion
clause of the geometry.  Fourth, the empty-coverage lemma says that
whenever \(\operatorname{Cov}(c)=\varnothing\),
\[
  \operatorname{Cov}(\mathcal F\cup\{c\})
  =
  \operatorname{Cov}(\mathcal F).
\]
Thus increasing the number of calibration names does not monotonically
increase closure power.  The four lemmas together give the asserted
coverage geometry.
\end{proof}\footnote{\raggedright This theorem is composed from the Lean
declarations \texttt{SixBirdsMetaMath.MetaMath.\allowbreak
CalibrationFamily.CoverageFamilyGeometry.\allowbreak
coverage\_family\_geometry} (the six-covered / four-uncovered /
two-named-holes core),
\texttt{all\_covered\_regions\_in\_family},
\texttt{all\_uncovered\_regions\_not\_in\_family},
\texttt{bounded\_depth\_overlap}, and
\texttt{empty\_coverage\_adds\_nothing} (the non-monotone clause);
see Appendix~\ref{app:formalization}.}

The non-monotone observation is the practical warning carried by the
theorem.  Adding names to the calibration family does not automatically
increase closure power.  A new calibration changes the geometry only if
its native carrier type reaches a previously uncovered region, or if a
separate bridge theorem transports content into that region.  Otherwise
the family has become larger without having filled a typed hole.

This geometry prepares the final two Part~III steps.
\Cref{sec:calibration:framework-type-self-char} turns the geometry
into a framework-type comparison on the \(P\) versus \(NP\) track,
and \Cref{sec:calibration:functorial-gates} names the functorial
gates that separate that framework-type from the standard Clay
target type.

\subsection{Framework-type self-characterization}
\label{sec:calibration:framework-type-self-char}

The coverage geometry describes where the calibration family reaches.
The framework-type packages that reach into a four-axis type signature.
For the \(P\) versus \(NP\) track, the resulting comparison is not merely
that some details are missing; it is that the native framework type and
the standard Clay target type differ along the axes that matter for the
target.

\begin{definition}[Framework-type]\label{def:framework-type}
A \FType{} is a four-axis record consisting of uniformity,
worst-caseness, scoping, and packaging.  On the \(P\) versus \(NP\)
track, the cascade framework-type is nonuniform, average-case, scoped,
and package-based.  The standard Clay target framework-type is uniform,
worst-case, unconditional, and language-based.  The uniform /
nonuniform and worst-case / average-case axes are the framework-type
analogues of the distinctions surveyed in \cite{Impagliazzo1995FiveWorlds}.
Thus
\[
  \FType
  =
  \mathsf{Uniformity}\times
  \mathsf{WorstCaseness}\times
  \mathsf{Scoping}\times
  \mathsf{Packaging}.
\]
For the \(P\) versus \(NP\) row,
\[
\begin{aligned}
  F_{\mathrm{PvNP}}
  &=
  (\mathsf{nonuniform},\mathsf{averageCase},
    \mathsf{scoped},\mathsf{packageBased}),\\
  T_{\mathrm{Clay}}
  &=
  (\mathsf{uniform},\mathsf{worstCase},
    \mathsf{unconditional},\mathsf{languageBased}).
\end{aligned}
\]
The matching predicate is componentwise equality on these four axes.
\end{definition}

\begin{theorem}[Framework-Type Self-Characterization]\label{thm:framework-type-self-char}
For the \(P\) versus \(NP\) substrate, the cascade framework-type does
not match the standard Clay target framework-type.  Concretely, the
nonuniform / average-case / scoped / package-based framework type differs
from the uniform / worst-case / unconditional / language-based target
type.
Equivalently,
\[
  \neg\operatorname{typesMatch}(F_{\mathrm{PvNP}},T_{\mathrm{Clay}})
  \quad\text{and}\quad
  |\mathcal G_{\mathrm{named}}|=3,
\]
where \(\mathcal G_{\mathrm{named}}\) is the named functorial-gate
catalog of \Cref{cat:named-functorial-gates}.
\end{theorem}

\begin{proof}
The two framework types are concrete four-tuples:
\[
\begin{aligned}
  F_{\mathrm{PvNP}}
  &=
  (\mathsf{nonuniform},\mathsf{averageCase},
    \mathsf{scoped},\mathsf{packageBased}),\\
  T_{\mathrm{Clay}}
  &=
  (\mathsf{uniform},\mathsf{worstCase},
    \mathsf{unconditional},\mathsf{languageBased}).
\end{aligned}
\]
The predicate \(\operatorname{typesMatch}\) is componentwise equality
on these four coordinates.  If the two types matched, their first
coordinates would be equal:
\[
  \mathsf{nonuniform}=\mathsf{uniform}.
\]
But these are distinct labels.  Hence
\[
  \neg\operatorname{typesMatch}(F_{\mathrm{PvNP}},T_{\mathrm{Clay}}).
\]
The same comparison also differs on the remaining three axes, but the
first coordinate already proves the mismatch.

It remains to identify the gate catalog size appearing in the theorem.
\Cref{cat:named-functorial-gates} supplies the three-entry named
functorial-gate catalog; taking its cardinality clause gives
\[
  |\mathcal G_{\mathrm{named}}|=3.
\]
Combining the componentwise PvNP mismatch with this cardinality clause
proves the theorem.
\end{proof}\footnote{\raggedright This theorem is verified in Lean as
\texttt{SixBirdsMetaMath.MetaMath.\allowbreak
CalibrationFamily.FrameworkTypeSelfChar.\allowbreak
framework\_type\_self\_characterization}; see
Appendix~\ref{app:formalization}.}

\begin{schema}[Framework-general candidate]\label{sch:framework-type-pending}
Beyond the PvNP-specific mismatch of
\Cref{thm:framework-type-self-char}, the framework-general candidate
asserts that for every Clay-class substrate lying in the V-Differential
trace-state-only column, the cascade framework-type and the standard
Clay target framework-type differ. Formally, for substrates \(c\) with
\(\mbox{\VDTSO}(c)=\mathsf{True}\),
\[
  \operatorname{cascadeFrameworkType}(c)\neq
  \operatorname{standardClayTargetType}(c).
\]
The candidate is conditional on the V-Differential column-membership
gate plus opaque per-substrate framework-type predicates, and is not
applied by the proof of \Cref{thm:framework-type-self-char}; it captures
the conjectural extension to substrates beyond PvNP, which is left as
future work.\footnote{\raggedright Tracked in the formalization harness
as a conditional schema in \texttt{SixBirdsMetaMath.MetaMath.\allowbreak
CalibrationFamily.FrameworkTypeSelfChar.\allowbreak
framework\_type\_self\_char\_pending}, conditioned on the V-Differential
trace-state-only column-membership gate plus opaque per-substrate
framework-type predicates; see Appendix~\ref{app:formalization}.}
\end{schema}

Thus \Cref{sch:framework-type-pending} is the framework-general
candidate, not an input to the PvNP-specific theorem just proved.

The conclusion is a self-characterization, not a claim that the standard
target has been reached by the calibration family alone.  It says where
the framework natively lives, and therefore which typed gates must be
supplied to move from the framework-type to the standard target type.  The
next subsection names those gates explicitly.

\subsection{Named functorial gates}\label{sec:calibration:functorial-gates}

The final Part~III object is neither a carrier state nor a calibration.
It is a specified transport problem: a typed source, a typed target, a
direction, and a theorem obligation saying what external mathematics
would have to supply.  These named functorial gates form the third
\TDS{} in the paper's calibration vocabulary.

\begin{definition}[Named functorial gate]\label{def:functorial-gate}
A named functorial gate consists of a name, a typed source, a typed
target, a direction, and a stated theorem obligation.  It is a typed
bridge-as-content: the gate does not assert that the transport has been
supplied, but records exactly what a supplied transport would have to
prove.
The record shape is
\[
  \mathsf{FunctorialGate}
  =
  \{\mathsf{name},\mathsf{source},\mathsf{target},
    \mathsf{direction},\mathsf{theoremObligation}\}.
\]
Here source and target are typed-source descriptions: calibration,
regime, or framework/target type labels, depending on the gate.
\end{definition}

\begin{table}[H]
  \centering
  \caption{Named functorial gates recorded by \Cref{cat:named-functorial-gates}.}
  \label{tab:functorial-gates}
  \scriptsize
  \renewcommand{\arraystretch}{1.08}
  \begin{tabularx}{\textwidth}{@{}>{\raggedright\arraybackslash}X@{}}
    \toprule
    Gate specification \\
    \midrule
    \(\GateFunctionDegree\). Source: function-degree / \(\Cvobs\) side; target: Nullstellensatz certificate-degree side; direction: source to target; obligation: relate detector-shadow degree to ideal-membership certificate degree. \\
    \addlinespace
    \(\GatePlantedStatistical\). Source: average-case statistical regime; target: worst-case NP-hard language regime; direction: average to worst; obligation: lift average-case detection content to worst-case language hardness. \\
    \addlinespace
    \(\GateUniformTime\). Source: nonuniform-family side; target: uniform time-hierarchy host layer; direction: nonuniform to uniform; obligation: realize the nonuniform cascade in a Williams-style uniform time-hierarchy setting \cite{Williams2014NonUniformACC}. \\
    \bottomrule
  \end{tabularx}
\end{table}
 
\begin{theorem}[Named Functorial Gates]\label{cat:named-functorial-gates}
The named functorial gates catalog is a \TDS{} consisting of the three
gates listed in Table~\ref{tab:functorial-gates}.  Each gate is
parameterized by typed source, typed target, direction, and stated theorem
obligation.  The catalog is the
species-level result: it records the three named gates as typed deposits,
while leaving each gate's individual theorem obligation to external
mathematics.
In symbols,
\[
  \mathcal G_{\mathrm{named}}
  =
  [G_1,G_2,G_3],
\]
with
\[
  |\mathcal G_{\mathrm{named}}|=3,
\]
where \(G_1,G_2,G_3\) are the three rows of
\Cref{tab:functorial-gates}, in the displayed order.  Every
\(g\in\mathcal G_{\mathrm{named}}\) has nonempty name and nonempty
theorem-obligation fields.
\end{theorem}

\begin{proof}
Define \(\mathcal G_{\mathrm{named}}\) to be the three records in
\Cref{tab:functorial-gates}.  Each record has the five fields in
\Cref{def:functorial-gate}: name, source, target, direction, and theorem
obligation.  Inspection of the displayed catalog gives exactly three
entries, so
\[
  |\mathcal G_{\mathrm{named}}|=3.
\]
Moreover, each displayed row has a nonempty name and a nonempty theorem
obligation field.  Therefore every \(g\in\mathcal G_{\mathrm{named}}\)
satisfies the species-level well-formedness condition.  This proves the
catalog theorem.  No individual gate obligation is solved by this
argument; those obligations are exactly the source-to-target transports
named in the fifth field of each gate record.
\end{proof}\footnote{\raggedright This theorem is verified in Lean as
\texttt{SixBirdsMetaMath.MetaMath.\allowbreak
CalibrationFamily.FunctorialGates.\allowbreak named\_functorial\_gates};
see Appendix~\ref{app:formalization}.}

The gates are the third typed-deposit species: typed bridges with stated
theorem obligation, parameterized by source / target / direction. Each
gate's individual theorem obligation awaits external mathematics; the
catalog row is the species-level result. Class-B Irremovability
(\cite{Tsiokos2026PvNP}) is positioned as a fourth methodological barrier
alongside relativization, natural-proofs, and algebrization
\cite{BakerGillSolovay1975Relativization,RazborovRudich1997NaturalProofs,AaronsonWigderson2009Algebrization}
for the gate-context.

Class-B Irremovability from \PaperPvNP\ \cite{Tsiokos2026PvNP} is
therefore complementary to the classical barriers, not a strengthening
of them.  Relativization, natural-proofs, and algebrization identify
well-known ways in which proof strategies fail to cross certain typed
boundaries.  The functorial-gate vocabulary records the corresponding
positive specification: if a future proof is to cross from the
framework-type to the standard target type, it must supply one of the
named transports as theorem-level content.

This closes Part~III.  The three gates remain open mathematical
obligations at the gate level, while the catalog that names them is part
of the verified calibration-family apparatus.  Section~8 returns to
predictions about further gates and further track replication.
 \section{Cross-paper consistency audit}\label{sec:audit}

The present section adds no formal claims to Parts~I--III and
restates no substrate theorem.  Its purpose is narrower: to check
that the framing decisions used by the meta-math paper are consistent
across the substrate corpus, and that the recognition sources used in
the two trace-state-only examples are recorded rather than silently
re-derived.

\subsection{Cross-paper consistency}\label{sec:audit:cross-paper-consistency}

(a) \emph{Recognition-mode and derivation-mode are peer paths under the
SB closure assumption.}  The first audit concerns the framing of
recognition-mode and derivation-mode landings on the \(P\) versus \(NP\)
track.  The recognition-mode landing should not be described as a
weaker fallback when a derivation route fails.  Under the SB closure
assumption, it is a peer closure path: the substrate contributes named
structural content through closure formation, while the framework
records that content and its target-facing readout.  The \PaperPvNP\
landing \cite{Tsiokos2026PvNP} treats recognition-mode closure as
structurally peer with derivation-mode closure, not as an inferior
substitute for it.

This peer-path audit is also what keeps \Cref{thm:attack-foreclosure-v5}
from being misread as a refutation of \Cref{thm:attack-foreclosure-v4}.
The v4 theorem continues to govern derivation routes.  The v5 refinement
separates those routes from recognition-mode landings, where the
load-bearing source is named as accepted recognition under the closure
discipline.  The audit conclusion is therefore a consistency statement
about framing, not a new theorem: the positive recognition path and the
negative derivation-foreclosure path occupy distinct typed roles.

(b) \emph{V-Differential column placement is consistent across the two
trace-state-only substrates.}  The second audit checks that the
trace-state-only classification is not being assigned ad hoc.  The
\PaperRH\ landing \cite{Tsiokos2026RH} and the \PaperPvNP{} landing both
place their substrates in the same V-Differential column: the structural
witness is not exposed as ordinary witness-content, but remains compressed
into trace-state evidence.  This is why both rows validate the
recognition-mode template rather than the witness-content-exposing
substrate pattern represented by Navier--Stokes.

The two examples are mathematically unlike each other: one is analytic
and trace-dual, the other is computational and hiddenness-based.  The
column-membership classification nevertheless uses the same criterion in
both cases.  In each case, the substrate supplies a trace-state-only
recognition source, the three-option sweep records the same
two-instance signature, and the closure-content thesis treats the source
as formed-layer content rather than framework-primitive output.  The
audit therefore supports the column classification used by Part~II, while
leaving corpus-wide generalization to the later predictions section.

(c) \emph{Recognition-closure architecture is consistent with Foundations
I closure formation.}  The third audit checks that the substrate papers do
not introduce private closure rules.  In the PvNP row, the CSL source
\cite{Tsiokos2026CSL} is named and imported as recognition content under
the same closure-formation discipline used elsewhere in the paper.  In
the RH row, the SDTC source \cite{Tsiokos2026SDTC} plays the same
architectural role: it is recorded as a named source, not manufactured by
a local derivation inside the meta-math framework.

This architecture is exactly the distinction made in
\Cref{thm:closure-content-thesis}.  A recognition source may have a
target-equivalent readout, but the source record is broader than the
readout: it carries provenance, closure identity, substrate placement,
and the structural law by which the closure is recognized.  The audit
therefore checks compatibility with Foundations I closure formation at
the level of source handling.  It does not certify the substrate
mathematics again; it verifies that the meta-math paper is using the same
closure architecture when it compares the two recognition landings.

\subsection{Source-record discipline}\label{sec:audit:source-record-discipline}

The source-record discipline is the bookkeeping rule that prevents a
recognition source from becoming an unnamed assumption or a disguised
derivation.  A source must be named, its substrate origin must be
identified, and its role in the closure must be separated from its
target-facing readout.  For PvNP this means the CSL hiddenness source is
recorded as named recognition; for RH this means the SDTC-Selberg source
is recorded as named recognition.  In both cases, the body of the paper
uses the language ``imported as named recognition per Foundations I
closure formation,'' not ``derived from framework primitives.''

This discipline also controls how the meta-math paper cites its
substrates.  \Cref{sec:validation} uses the substrate papers as
validation citations, not as objects to be reproved.  The present
audit then checks that the same source identity is preserved when the
framework-level conclusions are compared across tracks.  That is why the audit distinguishes a source record from
a readout: the readout may be the proposition-facing predicate that enters
the closure argument, while the source record is the wider structural
object whose provenance and closure status must remain visible.

The audit closes with a negative rule.  If a future recognition-mode
landing cannot name its source, cannot identify its column placement, or
cannot separate the source record from the target readout, it should not
be counted as an instance of the template in \Cref{thm:template-applicability}.
Those unresolved cases belong in the predictions and nonclaims registers,
not in the validated corpus.  \Cref{sec:predictions} turns to such
predictions.
 \section{Predictions and open problems}\label{sec:predictions}

This section records predictions and open problems only.  None of the
paragraphs below is used as evidence for the formal statements of
Parts~I--III, and none is a claim that a missing substrate landing has
already been supplied.  The common theme is that the present paper has
identified a small number of testable structural continuations: a
third trace-state-only substrate for the recognition-mode template, a
track-independent bound for CTMT recursion, and functorial gates beyond
the three named in the current typed-deposit catalog.

\subsection{Hodge-substrate column membership}
\label{sec:predictions:hodge}

The Hodge substrate
\cite{Deligne2000Hodge,Voisin2003HodgeTheoryII} is conjectured to lie
in the V-Differential trace-state-only column by structural reasoning;
no Hodge cascade has been constructed yet in the Six Birds corpus.
Conditional on a future Hodge cascade satisfying that column-membership
prediction, the recognition-mode landing template of
\Cref{sec:recognition} should apply with structurally identical
seven-stage shape.  In that event, the
\SigPattern{} verdict signature would be the third-instance corpus
test for the closure-content thesis.

The prediction is deliberately column-theoretic rather than
substrate-theoretic.  It does not assert a Hodge closure, and it does
not assert that the needed recognition source has already been named.
It says only that, if a future Hodge cascade exposes the same
trace-state-only behavior seen in the two-instance verified PvNP and
RH recognition-mode landings, then the seven-stage apparatus of
\Cref{thm:template-applicability} should be the correct template to
audit it.

This is also the cleanest proposed test of whether the \SigPattern{}
pattern is predictive rather than accidental.  The current evidence is
two-instance verified on PvNP and RH; a Hodge landing in the
trace-state-only column would be the leading candidate third instance.
Failure of the Hodge cascade to enter the trace-state-only column
would not refute the recognition-mode template, but it would redirect
the prediction to a different substrate class.

\subsection{Yang--Mills generalization conjecture}
\label{sec:predictions:ym}

Yang--Mills \cite{JaffeWitten2000YangMills} is a secondary
generalization candidate.  No Yang--Mills substrate exists in the
present corpus, and this paper makes no claim about a Yang--Mills
closure.  The open question is whether a future
Yang--Mills cascade should be expected to behave like a
trace-state-only substrate, like a witness-content-exposing substrate,
or like a new column not yet represented in the current
V-Differential split.

The conservative conjecture is therefore methodological: a future
Yang--Mills analysis should first be run through the column-membership
audit before any recognition-mode or derivation-mode template is
selected.  If the substrate is conjectured to be trace-state-only,
then the seven-stage recognition-mode template should become the
candidate landing apparatus.  If it is conjectured to expose witness
content directly, then the correct comparison should be with the
witness-content-exposing side of \Cref{sec:validation} rather than
with the Part~II schemas.

\subsection{Theorem-grade lifting of the Part~II schemas}
\label{sec:predictions:theorem-grade-lifting}

Lifting the four Part~II schemas to theorem-grade across the framework
requires two future inputs: (a) a third-instance substrate landing on
a V-Differential trace-state-only substrate beyond PvNP and RH, with
Hodge as the leading candidate; and (b) a structural argument from
the V-Differential column-membership predicate to the \SigPattern{}
verdict pattern.  Until both inputs exist, each of the
four Part~II universal claims remains a candidate meta-theorem.

The first input is empirical in the corpus sense: it asks for a third
landed trace-state-only example, audited by the same seven-stage
template and by the same two-instance signature comparison already
used for PvNP and RH.  The second input is structural: it asks for a
reason why trace-state-only column membership should force the
framework-primitive, SAU non-descent, and V-Differential elevation
sweep to display the \SigPattern{} pattern.  Without the structural
argument, a third example would still be suggestive rather than
theorem-grade.

This lifting problem is the main open problem left by Part~II.  The
recognition-mode path should remain peer with derivation-mode under
the SB closure assumption, but the paper-proposal-grade universal
form of these four claims should not be treated as final until both
future inputs have been supplied.

\subsection{Track-independent CTMT depth bound}
\label{sec:predictions:ctmt-depth-bound}

The CTMT Recursion Theorem of \Cref{thm:ctmt-recursion} supplies the
depth-indexed tower structure, but it does not supply a
track-independent numerical bound on tower depth.  The deepest
documented cascade in the current evidence base is the eight-layer RH
Branch B chain, and that observation should be read as a data point
rather than as a bound.

The open problem is to identify a substrate-independent invariant
that would bound CTMT depth across cascades.  Candidate invariants
could measure inherited-record complexity, descent height, native
matrix-family rank, or some future refinement of the carrier's
adequacy residual.  Any such proposal would need to be compatible with
the recursive constructor of \Cref{thm:ctmt-recursion}: a depth bound
cannot merely rule out further reductions by definition; it must
explain why the next deeper matrix-element family cannot arise under
the typed-context discipline.

A successful bound would sharpen the Part~I foreclosure account.  It
would not replace CTMT recursion; it would turn the present
open-ended tower into a finitely bounded tower for the relevant
substrate class.  Until such an invariant is found, CTMT depth remains
an open problem.

\subsection{Open functorial gates}
\label{sec:predictions:functorial-gates}

The Named Functorial Gates Catalog of
\Cref{cat:named-functorial-gates} records three species-level gates:
\[
  \begin{gathered}
    \GateFunctionDegree,\\
    \GatePlantedStatistical,\\
    \GateUniformTime .
  \end{gathered}
\]
Each individual gate still awaits the external mathematics that would
realize its stated theorem obligation.  Beyond those three, the
typed-deposit species is open-ended: further gates should be expected
whenever a calibration-to-calibration, regime-to-regime, or
framework-type-to-target-type transport is needed and cannot be
treated as a mere basis-change.

Two kinds of future gate work are predicted.  First, the three named
PvNP gates should be attacked individually as theorem obligations:
function-degree to Nullstellensatz certificate-degree, planted
statistical detection to worst-case language hardness, and
nonuniform-family realization to uniform time hierarchy.  Second,
new substrate tracks should expose additional functorial gates if
their calibration-family geometry has named typed holes analogous to
\(\xiNCone\) or \(\xiAdaptive\).

The open-gate program is therefore not a claim that the present
catalog is complete.  It is a prediction that the typed-deposit
species introduced in \Cref{sec:calibration:functorial-gates} should
continue to organize future handoff problems: a gap is either a
basis-change, a substantive theorem, or a named functorial gate with
an explicit theorem obligation still to be supplied.
 \section{Nonclaims and honest caveats}\label{sec:scope}

\subsection{Nonclaims register}\label{sec:scope:nonclaims}

This final section is the paper's binding nonclaim register.  It
collects the exclusions and honesty caveats that govern the earlier
sections: what substrate material is out of scope, which claims remain
conditional under the SB closure assumption, which recognition sources
are imported rather than derived, and which universal meta-claims are
candidate future work rather than results of this paper.

\begin{table}[H]
  \centering
  \caption{Nonclaim register: binding exclusions and caveats for the meta-math paper.}
  \label{tab:nonclaims}
  \scriptsize
  \renewcommand{\arraystretch}{1.08}
  \begin{tabularx}{\textwidth}{@{}>{\raggedright\arraybackslash}p{0.23\textwidth}>{\raggedright\arraybackslash}p{0.20\textwidth}>{\raggedright\arraybackslash}X@{}}
    \toprule
    Nonclaim & Source / scope & Required body wording \\
    \midrule
    BSD-track exclusion & M6 BSD material; no landed substrate paper & BSD-track results are out of scope; BSD has no landed substrate paper in the Tsiokos corpus and is left to future work. \\
    Hodge prediction only & M5 Hodge prediction; §8 only & Hodge is conjectural: conditional on a future Hodge cascade, the recognition-mode template should apply and the \SigPattern{} signature would be the third-instance corpus test. \\
    No standard-ZFC strengthening & all Clay-class references & Every referenced substrate landing is the conditional closure landed by its substrate paper under the SB closure assumption of Foundations~I. \\
    Recognition sources imported, not derived & TSO recognition sources & \(\Gamma_{\text{csl-sat-hidden}}\) and \(\Gamma_{\text{SDTC}}^{\text{Selberg}}\) are imported as named recognition per Foundations~I closure formation, not derived from framework primitives. \\
    Universal lifting deferred & T7--T10 & T7--T10 are paper-proposal-grade, two-instance verified, and corpus-pending until third-instance replication and a column-to-signature argument exist. \\
    M6 Theorem 7 folded into Part III & calibration-family background & M6's Non-Descending Translation content is historical motivation folded into Part~III, not a primary theorem of this paper. \\
    \bottomrule
  \end{tabularx}
\end{table}
 
The paragraphs below are the authoritative body register.  They govern
how the preceding sections should be read.

\subsection{BSD-track exclusion}\label{sec:scope:bsd-out}

BSD-track results are out of scope for this paper.  The meta-math
scope was fixed to substrates with landed three-Clay-class closures
(NS, RH, PvNP); BSD \cite{Wiles2000BSD} has no landed substrate paper
in the Tsiokos corpus and is left to future work.

This exclusion is not a negative statement about BSD.  It is a scope
decision: the present paper extracts the meta-mathematical structure
from the three landed conditional closures and does not use BSD as
evidence, motivation, or a partial validation row.

\subsection{Hodge prediction only}\label{sec:scope:hodge-prediction-only}

Hodge has no landed substrate paper in the present corpus.  It is
retained only as a structural prediction: the Hodge substrate is
conjectured to lie in the V-Differential trace-state-only column, so
the recognition-mode landing template should apply once a Hodge
cascade is built.  This prediction appears only in
\Cref{sec:predictions:hodge}.

Accordingly, this paper does not claim that Hodge is established, that
Hodge already satisfies the recognition-mode template, or that a Hodge
recognition source has been supplied.  The Hodge paragraph in
\Cref{sec:predictions} is a third-instance corpus test proposal, not
a result.

\subsection{No standard-ZFC strengthening}\label{sec:scope:standard-zfc}

Every Clay-class landing referenced by meta-math is conditional under
the SB closure assumption of Foundations~I.  The meta-math paper does
not strengthen any substrate landing to standard-ZFC-unconditional
status, and it does not claim that the substrate landings are
standard-ZFC-unconditional.

The binding form is: the conditional closure of NS, RH, or PvNP landed
by its substrate paper under the SB closure assumption of
Foundations~I.  Subsequent same-section references may abbreviate to
the substrate paper's conditional closure.  The paper never uses
``Paper~\(N\) proves the problem,'' ``Six Birds Theory proves the
problem,'' or any equivalent unconditional phrasing.

\subsection{Recognition-source imports}\label{sec:scope:recognition-source-imports}

Recognition sources are imported, not derived.  The recognition source
\(\Gamma_{\mathrm{csl\text{-}sat\text{-}hidden}}\) is imported as
named recognition per Foundations~I closure formation (per \PaperCSL),
and the recognition source
\(\Gamma_{\mathrm{SDTC}}^{\mathrm{Selberg}}\) is imported as named
recognition per Foundations~I closure formation (per \PaperRH\ and
\PaperSDTC).

The closure-content thesis is the corresponding meta-mathematical
claim: load-bearing closure content is named and accepted under the
closure-formation discipline, not produced by framework primitives.
Thus the paper never treats these sources as derived from the
framework, constructed internally by the framework, or manufactured by
the recognition-mode template itself.

\subsection{Universal lifting deferred}\label{sec:scope:universal-lifting-deferred}

The four Part~II schemas are paper-proposal-grade meta-claims,
two-instance verified on PvNP and RH and corpus-pending for
third-instance cross-track replication.  Theorem-grade lifting of these universal claims requires
two future inputs: a third-instance substrate landing on a
V-Differential trace-state-only substrate beyond PvNP and RH, with
Hodge as the leading candidate, and a structural argument from
V-Differential column-membership to the \SigPattern{} verdict pattern.

Until both inputs exist, the four Part~II schemas remain candidate
meta-theorems rather than universal theorem-grade results.  The v5 refinement continues to
refine, and not replace, v4: derivation routes remain subject to
Attack Foreclosure v4, while recognition-mode landings are peer with
derivation-mode under the SB closure assumption.

\subsection{Non-Descending Translation folded into Part III}
\label{sec:scope:m6-theorem-7-folded}

The original Non-Descending Translation meta-theorem motivated the
calibration-family apparatus historically; its translation/handoff
content has been folded into Part~III as the structural background
for the named functorial gates of
\Cref{sec:calibration:functorial-gates}.  The vobs-arc evidence from
\PaperPvNP\ and the 11+5 calibration sprint supplied the concrete
catalog from which the three gates of
\Cref{cat:named-functorial-gates} were extracted.

This paper does not restate the Non-Descending Translation
meta-theorem as a theorem, does not mechanize it as a primary target,
and does not cite it as supplying a theorem-grade result.  Its
retained role is historical and structural: the relevant formal
content lives in the Part~III calibration-family statements and in
the open functorial-gate program of \Cref{sec:predictions}.
 
\appendix
\section{Lean formalization}\label{app:formalization}



The mechanization companion for this paper is written in Lean~4
\cite{deMouraUllrich2021Lean4} on top of Mathlib
\cite{MathlibCommunity2020Lean}, in the broader tradition of
computer-assisted mathematical proofs that includes the four-color
theorem \cite{Gonthier2008FourColor}, the Kepler conjecture
\cite{HalesEtAl2017Flyspeck}, and recent large-scale formalizations
such as the Liquid Tensor Experiment \cite{CommelinTopaz2024LTE}; see
\cite{HalesNotices2008FormalProof} for a general overview of the
formal-proof program.

\subsection{Wording-discipline table}\label{app:formalization:wording-table}

This subsection is the canonical reference for how Lean-side status is
allowed to surface in body prose.  The body of the paper uses the
paper-prose wording below; raw Lean identifiers are confined to this
appendix.

\begin{table}[H]
\centering
\small
\begin{tabularx}{\textwidth}{@{}
  >{\raggedright\arraybackslash}p{0.16\textwidth}
  >{\raggedright\arraybackslash}X
  >{\raggedright\arraybackslash}X@{}}
\toprule
Lean status & Lean-facing disclosure wording & Forbidden disclosure wording \\
\midrule
\texttt{theorem}
  & ``Lean proves''; ``is formally verified''; ``Lean's proof goes by \dots''
  & ``verified informally''; ambiguous authorial ``we prove'' when the intended claim is the Lean result. \\
\texttt{partial}
  & ``Lean verifies the structural derivation applying the same-paper obligation''; ``Lean confirms the structural step.''
  & ``Lean proves'' without the partial-status qualifier; silent omission of the applied obligation. \\
\texttt{obligation}
  & ``Lean records this as a tagged obligation axiom under the opaque source-predicate \(\langle\)guard\(\rangle\); the obligation is honestly disclosed in the formalization trust base and is not silently accepted.''
  & ``Lean proves''; any indicative confirmation; bare Lean identifiers in body prose. \\
\texttt{substrate\_}\newline\texttt{pending}
  & Not used by the meta-math axis; included here for comparison with the general paper-writing policy.
  & Any claim that a substrate-pending item is theorem-grade. \\
\bottomrule
\end{tabularx}
\caption{Four-mode disclosure discipline for Lean-facing claims.}
\label{tab:appd:wording-discipline}
\end{table}

Meta-math adds an opacity-guard extension to the base policy.  Each
tagged obligation axiom is stated under an opaque source-predicate, so
the obligation cannot be collapsed by arbitrary instantiation in the
mechanization.  Body prose names the guard in paper language; this
appendix records the corresponding Lean-side predicate identifiers.
Across the axis there are nine opacity-guarded obligation axioms:
T1's carrier-dichotomy dependency, T2's bridge-impossibility
obligation, T3 Coverage, T5 Attack Foreclosure v4, the four
Recognition-Mode obligations T7--T10, and the framework-general T16
candidate.

\begin{table}[H]
\centering
\small
\begin{tabularx}{\textwidth}{@{}
  >{\raggedright\arraybackslash}p{0.35\textwidth}
  >{\raggedright\arraybackslash}X@{}}
\toprule
Paper-prose guard name & Lean predicate identifier(s) \\
\midrule
coverage-context guard
  & \LeanDecl{CoverageContext} \\
bridge-composition guard
  & \LeanDecl{BridgeComposition} \\
framework-scope guard
  & \LeanDecl{InFrameworkScope} \\
internal-move semantics (opaque)
  & \LeanDecl{internalMoveSemantics} \\
template-applies guard
  & \LeanDecl{TemplateApplies} \\
opaque substrate-typing predicates
  & \LeanDecl{FormedClosure}, \LeanDecl{VDifferentialTSOColumn}, \LeanDecl{SBClosureAssumption}, \LeanDecl{LandsAtTheoremGrade} \\
opaque sweep-verdict
  & \LeanDecl{sweepVerdict} \\
load-bearing predicate (opaque)
  & \LeanDecl{loadBearingFor} \\
derivation-route / recognition-landing / named-recognition predicates
  & \LeanDecl{IsDerivationRoute}, \LeanDecl{RecognitionModeLanding}, \LeanDecl{ImportsAsNamedRecognition} \\
per-substrate framework-type predicates (opaque)
  & \LeanDecl{cascadeFrameworkType}, \LeanDecl{standardClayTargetType} \\
\bottomrule
\end{tabularx}
\caption{Opacity-guard names used by obligation disclosures.}
\label{tab:appd:opacity-guards}
\end{table}

The three partial rows use the same discipline, but their proof terms
also perform structural work before applying a same-paper obligation.
T1 applies the coverage-context guarded carrier-dichotomy dependency;
T2 applies the Carrier Dichotomy and then the bridge-composition
obligation; T4 applies the Coverage Theorem under the coverage-context
guard.  Thus a partial row is not described as a theorem-grade Lean
proof, but it is also not merely a restatement of its obligation: the
body wording records both the verified structural derivation and the
specific obligation on which the derivation depends.

\subsection{Opacity-guard discussion}\label{app:formalization:opacity-guard-discussion}

The guard table above records names.  This subsection records the
mathematical role of each guard: what situation it marks, why the
formalization keeps it opaque, and which obligation or partial theorem
uses it.  Opacity is not a proof device for hiding content.  It is the
mechanized way of preventing a pending meta-mathematical source
condition from being replaced by an arbitrary user-chosen predicate.
Without this discipline, several obligation statements could be
trivialized or made inconsistent by instantiating their semantic
parameters with \(\top\) or \(\bot\).

\paragraph{Coverage-context guard.}
The coverage-context guard marks the admissible setting in which the
four-state coverage claim is meant to apply: a target \(X\), a typed
primary carrier \(C\) aimed at \(X\), and the surrounding audit
discipline that licenses the coverage question.  It is opaque because
not every syntactically constructible pair \((X,C)\) is an admissible
coverage instance; in particular, the reviewer-cycle instantiation
\(X=\Xi_C(D\mid L)=0\) with trivial carrier fields must not itself
produce coverage admissibility.  This guard is used by T1's
\LeanDecl{carrier_dichotomy_t3_pending}, by T3's
\LeanDecl{coverage_theorem}, and transitively by T4 when
\LeanDecl{type_finite_foreclosure} applies T3.

\paragraph{Bridge-composition guard.}
The bridge-composition guard marks the specific carrier-side bridge
whose composition with a carrier produces the composite carrier being
audited.  Its point is relational: the bridge \(B\) must be the bridge
that actually forms the composite \(C_B\), not an unrelated bridge
chosen after the fact.  The predicate is opaque because carrier-side
composition is abstract at the meta-math layer; concrete bridge
realizations live in substrate analyses.  It is used by T2's
\LeanDecl{bridge_impossibility_carries_obligation}.

\paragraph{Framework-scope guard.}
The framework-scope guard marks targets to which the internal-move
audit applies.  Attack Foreclosure v4 is not a universal statement over
all propositions; it is a statement about targets inside the framework's
typed audit scope.  The guard is opaque because framework scope is a
semantic admissibility condition supplied by the surrounding coverage,
bridge, and cascade audit discipline.  It is used by T5's
\LeanDecl{attack_foreclosure_v4} and by T10's
\LeanDecl{attack_foreclosure_v5}.

\paragraph{Opaque internal-move semantics.}
The opaque internal-move semantics carrier gives meaning to the three
framework-internal move classes: carrier pivots, bridge imports, and
cascade-internal computations.  It is opaque because a free semantic
parameter could be instantiated as the constant-true predicate, turning
the foreclosure axiom into a contradiction.  Fixing the semantics at
module scope closes that instantiation path.  This opaque semantic
carrier is used by T5's \LeanDecl{attack_foreclosure_v4}.

\paragraph{Template-applies guard.}
The template-applies guard marks substrates for which the seven-stage
recognition-mode landing audit has actually succeeded.  The seven stage
record documents the intended stage structure, but arbitrary
propositional fields are not by themselves evidence that the audit has
been carried out.  The guard is therefore opaque: it cannot be built
merely by setting stage placeholders to true.  It is used by T7's
\LeanDecl{recognition_template_applicability} and by T8's
\LeanDecl{verdict_signature}.

\paragraph{Opaque substrate-typing predicates.}
The four substrate-typing predicates record the conditions under which
the recognition template is meant to apply.  They mark, respectively,
that a closure has been formed, that its substrate lies in the
trace-state-only column, that the SB closure assumption is in force, and
that the landing reaches theorem-grade.  They remain opaque because
these are substrate-facing conditions, not definitional computations
inside the meta-math harness.  They compose with the template-applies
guard in T7; their Lean identifiers are listed in
\Cref{tab:appd:opacity-guards}.

\paragraph{Opaque sweep-verdict.}
The sweep-verdict function assigns a verdict to each of the three
Stage-6 sweep options.  It is opaque so that the
\SigPattern{} signature is not merely the result of a definition that
hard-codes the desired answers.  The obligation must assert the
cross-track verdict pattern; it must not be provable by unfolding the
verdict reader.  This guard is used by T8's
\LeanDecl{verdict_signature}.

\paragraph{Opaque load-bearing predicate.}
The load-bearing predicate selects the structural content that carries
the closure for a trace-state-only substrate.  It is opaque because a
transparent definition could bake into the returned object the three
properties that the Closure-Content Thesis is supposed to assert:
formed-layer closure content, non-derivability from framework
primitives, and named accepted-source status.  It is used by T9's
\LeanDecl{closure_content_thesis}.

\paragraph{Derivation-route, recognition-landing, and named-recognition predicates.}
The v5 refinement uses three opaque predicates to express its case
split.  They mark, respectively, the derivation-route branch, the
recognition-mode landing branch, and the act of importing a
load-bearing source as named recognition.  They are opaque because
transparent constructors or true-valued predicates would make the v5
recognition clauses vacuous.  They are used by T10; their Lean
identifiers are listed in \Cref{tab:appd:opacity-guards}.

\paragraph{Opaque per-substrate framework-type predicates.}
The predicates \LeanDecl{cascadeFrameworkType} and
\LeanDecl{standardClayTargetType} assign framework-type records to a
Clay-class trace-state-only substrate and to its standard Clay target
side.  They are opaque because the framework-general claim is
track-dependent: the PvNP instance is proved directly in T16, but the
universal Clay-class statement is not.  These predicates are used by
T16-pending, \LeanDecl{framework_type_self_char_pending}.

\subsection{T-row coverage map}\label{app:formalization:t-row-coverage}

The following tables give the formalization coverage map for the
seventeen T-rows, the three same-paper obligation sub-rows, and the
five supporting lemmas that are part of the meta-math axis audit.  All
module names below are relative to the prefix
\LeanDecl{SixBirdsMetaMath.MetaMath}.

\begin{table}[H]
  \centering
  \caption{Lean coverage summary by Part (manifest cross-check 2026-05-27).}
  \label{tab:lean-coverage-by-part}
  \small
  \begin{tabularx}{\textwidth}{@{}>{\raggedright\arraybackslash}p{0.17\textwidth}>{\centering\arraybackslash}X>{\centering\arraybackslash}X>{\centering\arraybackslash}X>{\centering\arraybackslash}X>{\centering\arraybackslash}X@{}}
    \toprule
    Part & Primary theorems & Supporting lemmas & Partials & Obligations & Definitions \\
    \midrule
    Part I --- Foreclosure & 1 & 0 & 3 & 4 & \(\sim 15\) \\
    Part II --- Recognition-Mode & 0 & 0 & 0 & 4 & \(\sim 17\) \\
    Part III --- Calibration Family & 7 & 5 & 0 & 1 & \(\sim 30\) \\
    \bottomrule
  \end{tabularx}
\end{table}
 
\begin{table}[H]
\centering
\scriptsize
\renewcommand{\arraystretch}{1.08}
\begin{tabularx}{\textwidth}{@{}
  >{\raggedright\arraybackslash}p{0.08\textwidth}
  >{\raggedright\arraybackslash}p{0.25\textwidth}
  >{\raggedright\arraybackslash}p{0.10\textwidth}
  >{\raggedright\arraybackslash}p{0.17\textwidth}
  >{\raggedright\arraybackslash}X@{}}
\toprule
Label & Paper-prose name & Status & Guard & Notes \\
\midrule
T1 & Carrier Dichotomy & partial & coverage-context & Structural branch proof applies same-paper obligation. \\
\multicolumn{5}{@{}p{\textwidth}@{}}{\emph{Lean:} module \LeanDecl{Foreclosure.CarrierDichotomy}; decl \LeanDecl{carrier_dichotomy}.} \\
T1-obl. & Carrier-dichotomy T3 dependency & obligation & coverage-context & X-equivalent branch classification. \\
\multicolumn{5}{@{}p{\textwidth}@{}}{\emph{Lean:} module \LeanDecl{Foreclosure.CarrierDichotomy}; decl \LeanDecl{carrier_dichotomy_t3_pending}.} \\
T2 & Bridge Impossibility & partial & bridge-composition & Applies T1 plus bridge obligation. \\
\multicolumn{5}{@{}p{\textwidth}@{}}{\emph{Lean:} module \LeanDecl{Foreclosure.BridgeImpossibility}; decl \LeanDecl{bridge_impossibility}.} \\
T2-obl. & Bridge X-strength obligation & obligation & bridge-composition & Conceptual ``therefore'' step for carrier-side bridges. \\
\multicolumn{5}{@{}p{\textwidth}@{}}{\emph{Lean:} module \LeanDecl{Foreclosure.BridgeImpossibility}; decl \LeanDecl{bridge_impossibility_carries_obligation}.} \\
T3 & Coverage Theorem & obligation & coverage-context & Four-state exhaustiveness. \\
\multicolumn{5}{@{}p{\textwidth}@{}}{\emph{Lean:} module \LeanDecl{Foreclosure.Coverage}; decl \LeanDecl{coverage_theorem}.} \\
T4 & Type-Finite Foreclosure & partial & coverage-context & Total finite carrier-state classification. \\
\multicolumn{5}{@{}p{\textwidth}@{}}{\emph{Lean:} module \LeanDecl{Foreclosure.TypeFiniteForeclosure}; decl \LeanDecl{type_finite_foreclosure}.} \\
T5 & Attack Foreclosure v4 & obligation & framework-scope & Opaque internal-move semantics. \\
\multicolumn{5}{@{}p{\textwidth}@{}}{\emph{Lean:} module \LeanDecl{Foreclosure.AttackForeclosure}; decl \LeanDecl{attack_foreclosure_v4}.} \\
T6 & CTMT Recursion & theorem & -- & Depth-indexed tower structure. \\
\multicolumn{5}{@{}p{\textwidth}@{}}{\emph{Lean:} module \LeanDecl{Foreclosure.CTMTRecursion}; decl \LeanDecl{ctmt_recursion}.} \\
\bottomrule
\end{tabularx}
\caption{Part I formalization coverage.}
\label{tab:appd:coverage-part-i}
\end{table}

\begin{table}[H]
\centering
\scriptsize
\renewcommand{\arraystretch}{1.08}
\begin{tabularx}{\textwidth}{@{}
  >{\raggedright\arraybackslash}p{0.08\textwidth}
  >{\raggedright\arraybackslash}p{0.25\textwidth}
  >{\raggedright\arraybackslash}p{0.10\textwidth}
  >{\raggedright\arraybackslash}p{0.21\textwidth}
  >{\raggedright\arraybackslash}X@{}}
\toprule
Label & Paper-prose name & Status & Guard & Notes \\
\midrule
T7 & Template Applicability & obligation & template-applies & Paper-proposal-grade; two-instance verified. \\
\multicolumn{5}{@{}p{\textwidth}@{}}{\emph{Lean:} module \LeanDecl{RecognitionMode.Template}; decl \LeanDecl{recognition_template_applicability}.} \\
T8 & Verdict Signature & obligation & template-applies; sweep-verdict & Records two-instance \SigPattern{} sweep signature. \\
\multicolumn{5}{@{}p{\textwidth}@{}}{\emph{Lean:} module \LeanDecl{RecognitionMode.VerdictSignature}; decl \LeanDecl{verdict_signature}.} \\
T9 & Closure-Content Thesis & obligation & load-bearing predicate & Recognition sources imported, not derived. \\
\multicolumn{5}{@{}p{\textwidth}@{}}{\emph{Lean:} module \LeanDecl{RecognitionMode.ClosureContentThesis}; decl \LeanDecl{closure_content_thesis}.} \\
T10 & Attack Foreclosure v5 & obligation & framework-scope; recognition predicates & v5 refines, does not replace, v4. \\
\multicolumn{5}{@{}p{\textwidth}@{}}{\emph{Lean:} module \LeanDecl{RecognitionMode.AttackForeclosureV5}; decl \LeanDecl{attack_foreclosure_v5}.} \\
\bottomrule
\end{tabularx}
\caption{Part II formalization coverage.}
\label{tab:appd:coverage-part-ii}
\end{table}

\begin{table}[H]
\centering
\scriptsize
\renewcommand{\arraystretch}{1.08}
\begin{tabularx}{\textwidth}{@{}
  >{\raggedright\arraybackslash}p{0.08\textwidth}
  >{\raggedright\arraybackslash}p{0.25\textwidth}
  >{\raggedright\arraybackslash}p{0.10\textwidth}
  >{\raggedright\arraybackslash}p{0.17\textwidth}
  >{\raggedright\arraybackslash}X@{}}
\toprule
Label & Paper-prose name & Status & Guard & Notes \\
\midrule
T11 & Calibration-Family Structure & theorem & -- & Thirteen-entry open-ended seed catalog. \\
\multicolumn{5}{@{}p{\textwidth}@{}}{\emph{Lean:} module \LeanDecl{CalibrationFamily.Structure}; decl \LeanDecl{calibration_family_structure}.} \\
T12 & Rebadge Discipline & theorem & -- & Four-axis profile discriminator. \\
\multicolumn{5}{@{}p{\textwidth}@{}}{\emph{Lean:} module \LeanDecl{CalibrationFamily.RebadgeDiscipline}; decl \LeanDecl{rebadge_discipline}.} \\
T13 & Bridges-as-Theorems & theorem & -- & Basis-change or substantive theorem. \\
\multicolumn{5}{@{}p{\textwidth}@{}}{\emph{Lean:} module \LeanDecl{CalibrationFamily.BridgesAsTheorems}; decl \LeanDecl{bridges_as_theorems}.} \\
T14 & Bridge Type-Direction & theorem & -- & Reverse symmetry iff source and target coincide. \\
\multicolumn{5}{@{}p{\textwidth}@{}}{\emph{Lean:} module \LeanDecl{CalibrationFamily.BridgeTypeDirection}; decl \LeanDecl{bridge_type_direction}.} \\
T15 & Coverage-Family Geometry & theorem & -- & Composed from supporting lemmas. \\
\multicolumn{5}{@{}p{\textwidth}@{}}{\emph{Lean:} module \LeanDecl{CalibrationFamily.CoverageFamilyGeometry}; decl \LeanDecl{coverage_family_geometry}.} \\
T16 & Framework-Type Self-Characterization & theorem & -- & PvNP-specific mismatch plus T17 length. \\
\multicolumn{5}{@{}p{\textwidth}@{}}{\emph{Lean:} module \LeanDecl{CalibrationFamily.FrameworkTypeSelfChar}; decl \LeanDecl{framework_type_self_characterization}.} \\
T16-pend. & Framework-general candidate & obligation & V-Differential; per-substrate types & Not applied by T16 proof. \\
\multicolumn{5}{@{}p{\textwidth}@{}}{\emph{Lean:} module \LeanDecl{CalibrationFamily.FrameworkTypeSelfChar}; decl \LeanDecl{framework_type_self_char_pending}.} \\
T17 & Named Functorial Gates & theorem & -- & Three-gate species-level catalog. \\
\multicolumn{5}{@{}p{\textwidth}@{}}{\emph{Lean:} module \LeanDecl{CalibrationFamily.FunctorialGates}; decl \LeanDecl{named_functorial_gates}.} \\
\bottomrule
\end{tabularx}
\caption{Part III primary T-row formalization coverage.}
\label{tab:appd:coverage-part-iii}
\end{table}

\begin{table}[H]
\centering
\scriptsize
\renewcommand{\arraystretch}{1.08}
\begin{tabularx}{\textwidth}{@{}
  >{\raggedright\arraybackslash}p{0.12\textwidth}
  >{\raggedright\arraybackslash}p{0.28\textwidth}
  >{\raggedright\arraybackslash}p{0.10\textwidth}
  >{\raggedright\arraybackslash}X@{}}
\toprule
Label & Supporting lemma & Status & Notes \\
\midrule
T15-lem. & Covered regions & theorem & Six covered regions in family coverage. \\
\multicolumn{4}{@{}p{\textwidth}@{}}{\emph{Lean:} module \LeanDecl{CalibrationFamily.CoverageFamilyGeometry}; decl \LeanDecl{all_covered_regions_in_family}.} \\
T15-lem. & Uncovered regions & theorem & Four uncovered regions outside coverage. \\
\multicolumn{4}{@{}p{\textwidth}@{}}{\emph{Lean:} module \LeanDecl{CalibrationFamily.CoverageFamilyGeometry}; decl \LeanDecl{all_uncovered_regions_not_in_family}.} \\
T15-lem. & Bounded-depth overlap & theorem & Shared sub-region overlap. \\
\multicolumn{4}{@{}p{\textwidth}@{}}{\emph{Lean:} module \LeanDecl{CalibrationFamily.CoverageFamilyGeometry}; decl \LeanDecl{bounded_depth_overlap}.} \\
T15-lem. & Non-monotone coverage & theorem & Empty calibration adds no coverage. \\
\multicolumn{4}{@{}p{\textwidth}@{}}{\emph{Lean:} module \LeanDecl{CalibrationFamily.CoverageFamilyGeometry}; decl \LeanDecl{empty_coverage_adds_nothing}.} \\
T16-lem. & PvNP mismatch & theorem & Direct record-field mismatch. \\
\multicolumn{4}{@{}p{\textwidth}@{}}{\emph{Lean:} module \LeanDecl{CalibrationFamily.FrameworkTypeSelfChar}; decl \LeanDecl{pvnp_framework_type_mismatch}.} \\
\bottomrule
\end{tabularx}
\caption{Part III supporting-lemma coverage.}
\label{tab:appd:coverage-part-iii-lemmas}
\end{table}

\subsection{Master mapping}\label{app:formalization:master-mapping}

This table is the canonical one-place claim-to-Lean map for the
meta-math paper.  It is cross-checkable against
\texttt{paper/meta\_math/notes/statements-of-record.yml}; the
Part-by-Part tables above remain the expanded local views.

\begin{landscape}
\begin{table}[p]
\centering
\fontsize{4.6}{4.8}\selectfont
\renewcommand{\arraystretch}{0.74}
\newcommand{\MasterDecl}[2]{\LeanDecl{SixBirdsMetaMath.MetaMath.}\newline\LeanDecl{#1.}\newline\LeanDecl{#2}}
\begin{tabularx}{\linewidth}{@{}
  >{\raggedright\arraybackslash}p{0.06\linewidth}
  >{\raggedright\arraybackslash}p{0.15\linewidth}
  >{\raggedright\arraybackslash}p{0.32\linewidth}
  >{\raggedright\arraybackslash}p{0.12\linewidth}
  >{\raggedright\arraybackslash}p{0.12\linewidth}
  >{\raggedright\arraybackslash}X@{}}
\toprule
Paper label & Paper-prose name & Lean declaration path & Coverage / alignment status & Footnote wording rule & Body location \\
\midrule
T1 & Carrier Dichotomy & \MasterDecl{Foreclosure.CarrierDichotomy}{carrier_dichotomy} & partial & narrowed & \S3.2, \Cref{thm:carrier-dichotomy} \\
T2 & Bridge Impossibility & \MasterDecl{Foreclosure.BridgeImpossibility}{bridge_impossibility} & partial & narrowed & \S3.3, \Cref{cor:bridge-impossibility} \\
T3 & Coverage Theorem & \MasterDecl{Foreclosure.Coverage}{coverage_theorem} & obligation/projection-schema & projection-schema & \S3.4, \Cref{thm:coverage} \\
T4 & Type-Finite Foreclosure & \MasterDecl{Foreclosure.TypeFiniteForeclosure}{type_finite_foreclosure} & partial & narrowed & \S3.5, \Cref{thm:type-finite-foreclosure} \\
T5 & Attack Foreclosure v4 & \MasterDecl{Foreclosure.AttackForeclosure}{attack_foreclosure_v4} & obligation/projection-schema & projection-schema & \S3.6, \Cref{thm:attack-foreclosure-v4} \\
T6 & CTMT Recursion & \MasterDecl{Foreclosure.CTMTRecursion}{ctmt_recursion} & theorem & faithful & \S3.7, \Cref{thm:ctmt-recursion} \\
T7 & Recognition-Mode Template Applicability & \MasterDecl{RecognitionMode.Template}{recognition_template_applicability} & obligation/projection-schema & projection-schema & \S4.2, \Cref{thm:template-applicability} \\
T8 & Cross-Track Verdict Signature & \MasterDecl{RecognitionMode.VerdictSignature}{verdict_signature} & obligation/projection-schema & projection-schema & \S4.3, \Cref{thm:verdict-signature} \\
T9 & Closure-Content Thesis & \MasterDecl{RecognitionMode.ClosureContentThesis}{closure_content_thesis} & obligation/projection-schema & projection-schema & \S4.4, \Cref{thm:closure-content-thesis} \\
T10 & Attack Foreclosure v5 & \MasterDecl{RecognitionMode.AttackForeclosureV5}{attack_foreclosure_v5} & obligation/projection-schema & projection-schema & \S4.5, \Cref{thm:attack-foreclosure-v5} \\
T11 & Calibration-Family Structure & \MasterDecl{CalibrationFamily.Structure}{calibration_family_structure} & theorem & faithful & \S6.2, \Cref{thm:calibration-family-structure} \\
T12 & Rebadge Discipline & \MasterDecl{CalibrationFamily.RebadgeDiscipline}{rebadge_discipline} & theorem & faithful & \S6.3, \Cref{thm:rebadge-discipline} \\
T13 & Bridges-as-Theorems & \MasterDecl{CalibrationFamily.BridgesAsTheorems}{bridges_as_theorems} & theorem & faithful & \S6.4, \Cref{thm:bridges-as-theorems} \\
T14 & Bridge Type-Direction & \MasterDecl{CalibrationFamily.BridgeTypeDirection}{bridge_type_direction} & theorem & faithful & \S6.5, \Cref{thm:bridge-type-direction} \\
T15 & Coverage-Family Geometry & \MasterDecl{CalibrationFamily.CoverageFamilyGeometry}{coverage_family_geometry} & theorem & faithful & \S6.6, \Cref{thm:coverage-family-geometry} \\
T16 & Framework-Type Self-Characterization & \MasterDecl{CalibrationFamily.FrameworkTypeSelfChar}{framework_type_self_characterization} & theorem & faithful & \S6.7, \Cref{thm:framework-type-self-char} \\
T16-pending & Framework-general candidate & \MasterDecl{CalibrationFamily.FrameworkTypeSelfChar}{framework_type_self_char_pending} & obligation/projection-schema & projection-schema & \S6.7, \Cref{sch:framework-type-pending} \\
T17 & Named Functorial Gates & \MasterDecl{CalibrationFamily.FunctorialGates}{named_functorial_gates} & theorem & faithful & \S6.8, \Cref{cat:named-functorial-gates} \\
T15-lem. & Covered regions & \MasterDecl{CalibrationFamily.CoverageFamilyGeometry}{all_covered_regions_in_family} & theorem & faithful & \Cref{tab:appd:coverage-part-iii-lemmas}; supports \Cref{thm:coverage-family-geometry} \\
T15-lem. & Uncovered regions & \MasterDecl{CalibrationFamily.CoverageFamilyGeometry}{all_uncovered_regions_not_in_family} & theorem & faithful & \Cref{tab:appd:coverage-part-iii-lemmas}; supports \Cref{thm:coverage-family-geometry} \\
T15-lem. & Bounded-depth overlap & \MasterDecl{CalibrationFamily.CoverageFamilyGeometry}{bounded_depth_overlap} & theorem & faithful & \Cref{tab:appd:coverage-part-iii-lemmas}; supports \Cref{thm:coverage-family-geometry} \\
T15-lem. & Non-monotone coverage & \MasterDecl{CalibrationFamily.CoverageFamilyGeometry}{empty_coverage_adds_nothing} & theorem & faithful & \Cref{tab:appd:coverage-part-iii-lemmas}; supports \Cref{thm:coverage-family-geometry} \\
T16-lem. & PvNP mismatch & \MasterDecl{CalibrationFamily.FrameworkTypeSelfChar}{pvnp_framework_type_mismatch} & theorem & faithful & \Cref{tab:appd:coverage-part-iii-lemmas}; supports \Cref{thm:framework-type-self-char} \\
\bottomrule
\end{tabularx}
\caption{Master claim-to-Lean mapping (one row per main theorem; cross-checkable against \texttt{paper/meta\_math/notes/statements-of-record.yml}).}
\label{tab:appd:master-mapping}
\end{table}
\end{landscape}

\subsection{Tagged-axiom register}\label{app:formalization:obligation-register}

The meta-math trust base contains exactly nine tagged obligation axioms.
Each one is guarded by an opaque source predicate or opaque semantic
carrier, and the conclusion shape is recorded explicitly so that the
reader can distinguish theorem-grade Lean proofs from same-paper
candidate content.  The full names below are the Lean declarations as
they occur under the prefix \LeanDecl{SixBirdsMetaMath.MetaMath}.

\begin{table}[H]
  \centering
  \caption{Tagged-axiom register for the nine meta-math \texttt{\# obligation} axioms.}
  \label{tab:tagged-axiom-register}
  \scriptsize
  \renewcommand{\arraystretch}{1.08}
  \begin{tabularx}{\textwidth}{@{}
    >{\raggedright\arraybackslash}p{0.14\textwidth}
    >{\raggedright\arraybackslash}p{0.23\textwidth}
    >{\raggedright\arraybackslash}X@{}}
    \toprule
    Tag & Opacity guard & Source / conclusion shape \\
    \midrule
    \multicolumn{3}{@{}p{\textwidth}@{}}{\LeanDecl{Foreclosure.CarrierDichotomy.carrier_dichotomy_t3_pending}} \\
    \texttt{\# obligation} & coverage-context & T1 dependency on T3: an \(X\)-equivalent admissible carrier is classified into exactly one of \CRETE, \CREBF, or \CRECTMT. \\
    \addlinespace
    \multicolumn{3}{@{}p{\textwidth}@{}}{\LeanDecl{Foreclosure.BridgeImpossibility.bridge_impossibility_carries_obligation}} \\
    \texttt{\# obligation} & bridge-composition & T2 ``therefore'' step: a bridge that composes a carrier into a classified composite carries the target \(X\)-strength obligation. \\
    \addlinespace
    \multicolumn{3}{@{}p{\textwidth}@{}}{\LeanDecl{Foreclosure.Coverage.coverage_theorem}} \\
    \texttt{\# obligation} & coverage-context & T3 central coverage claim: every admissible typed carrier aimed at \(X\) lies in one of the four typed states. \\
    \addlinespace
    \multicolumn{3}{@{}p{\textwidth}@{}}{\LeanDecl{Foreclosure.AttackForeclosure.attack_foreclosure_v4}} \\
    \texttt{\# obligation} & framework-scope; opaque internal-move semantics & T5 operational foreclosure: an in-scope framework-internal move does not advance \(X\) without external mathematical content. \\
    \addlinespace
    \multicolumn{3}{@{}p{\textwidth}@{}}{\LeanDecl{RecognitionMode.Template.recognition_template_applicability}} \\
    \texttt{\# obligation} & template-applies; opaque substrate-typing predicates & T7 applicability: a formed TSO closure satisfying the seven-stage template lands at theorem-grade under the SB closure assumption. \\
    \addlinespace
    \multicolumn{3}{@{}p{\textwidth}@{}}{\LeanDecl{RecognitionMode.VerdictSignature.verdict_signature}} \\
    \texttt{\# obligation} & template-applies; opaque sweep-verdict & T8 Stage-6 signature: the three sweep options return the two-instance \SigPattern{} verdict pattern. \\
    \addlinespace
    \multicolumn{3}{@{}p{\textwidth}@{}}{\LeanDecl{RecognitionMode.ClosureContentThesis.closure_content_thesis}} \\
    \texttt{\# obligation} & opaque load-bearing predicate & T9 closure-content thesis: load-bearing TSO closure content is formed-layer content, is not framework-primitive derivable, and is named as accepted source. \\
    \addlinespace
    \multicolumn{3}{@{}p{\textwidth}@{}}{\LeanDecl{RecognitionMode.AttackForeclosureV5.attack_foreclosure_v5}} \\
    \texttt{\# obligation} & framework-scope; opaque recognition predicates & T10 v5 refinement: derivation routes remain v4-foreclosed while recognition-mode landings import named recognition and land at theorem-grade. \\
    \addlinespace
    \multicolumn{3}{@{}p{\textwidth}@{}}{\LeanDecl{CalibrationFamily.FrameworkTypeSelfChar.framework_type_self_char_pending}} \\
    \texttt{\# obligation} & V-Differential TSO membership; opaque per-substrate framework types & T16 framework-general candidate: every Clay-class TSO substrate has cascade framework-type distinct from its standard Clay target type. \\
    \bottomrule
  \end{tabularx}
\end{table}
 
\paragraph{T1 carrier-dichotomy dependency.}
The T1 dependency axiom asserts the \(X\)-equivalent branch of the
Carrier Dichotomy: under the
typed-context discipline, aimed-at-\(X\) hypothesis, and coverage-context
guard, an \(X\)-equivalent admissible carrier lies in exactly one of
\CRETE, \CREBF, or \CRECTMT.  It is recorded as an obligation because
the three-way exhaustiveness of the \(X\)-equivalent branch is precisely
the same-paper coverage content, not a consequence of the local carrier
definitions alone.  The source statement is the same-paper Carrier
Dichotomy, and the narrowing to this branch is the result of the T1
substance and target-indexing reviewer cycles recorded in
\Cref{app:formalization:reviewer-log}.

\paragraph{T2 bridge-composition obligation.}
The T2 bridge-composition axiom asserts the ``therefore'' step for
carrier-side bridges: if a bridge \(B\) composes \(C\) into the
classified composite \(C_B\), then \(B\) carries the \(X\)-strength
obligation.  It is recorded as an obligation because the passage from
a classified composite carrier to an obligation on the bridge itself
is a conceptual bridge-content claim, not supplied by the Carrier
Dichotomy alone.  The source statement is the same-paper Bridge
Impossibility corollary.  The bridge-composition guard was added in
the axis re-review so that unrelated bridges cannot be certified by a
classified composite.

\paragraph{T3 Coverage Theorem.}
The T3 coverage axiom asserts the central four-state exhaustiveness
claim: every admissible typed primary
carrier aimed at \(X\) lies in at least one of the states
\(\nonXequiv\), \CRETE, \CREBF, or \CRECTMT.  It is recorded as an
obligation because the precise definition of ``open conjecture relative
to an accepted-imports base'' and the relative-consistency argument
that prevents an admissible carrier from escaping the four-state
classification are left open.  The source statement is the same-paper
Coverage Theorem together with its pending note.  The coverage-context
guard was introduced during the axis-level consistency review to block
non-admissible instantiations.

\paragraph{T5 Attack Foreclosure v4.}
The T5 v4 foreclosure axiom asserts that an in-scope framework-internal
move does not advance the
target \(X\) without external mathematical content.  Its mathematical
content is the operational foreclosure of the three internal move
classes: carrier pivots, bridge imports, and cascade-internal
computations.  It is recorded as an obligation because the derivation
from Carrier Dichotomy, Bridge Impossibility, and CTMT Recursion is
given informally rather than as a completed theorem-grade proof.  The
source statement is the same-paper Attack Foreclosure v4 schema.  The
framework-scope guard and opaque internal-move semantics were added
during the axis-level opacity review to prevent over-quantification over
arbitrary targets or arbitrary move semantics.

\paragraph{T7 Recognition-Mode Template Applicability.}
The T7 template-applicability axiom asserts that the seven-stage
recognition-mode landing template applies
to formed closures in the V-Differential trace-state-only column and
lands at theorem-grade under the SB closure assumption.  It is recorded
as an obligation because the universal template is paper-proposal-grade:
it is verified on the PvNP and RH instances but awaits corpus-level
lifting.  The source statement is the same-paper Recognition-Mode
Template Applicability schema.  The template-applies guard and opaque
substrate-typing predicates were introduced during the axis-level
opacity review.

\paragraph{T8 Cross-Track Verdict Signature.}
The T8 verdict-signature axiom asserts the Stage-6 \SigPattern{}
signature for the three sweep options:
framework-primitive, SAU non-descent, and V-Differential elevation.  It
is recorded as an obligation because the verdict reader is deliberately
opaque: the pattern must be a cross-track assertion, not a tautology
obtained by unfolding a definition.  The source statement is the
same-paper Cross-Track Verdict Signature schema.  The T8 opacity fix
replaced a hard-coded \LeanDecl{sweepVerdict} definition with an
opaque verdict function.

\paragraph{T9 Closure-Content Thesis.}
The T9 closure-content axiom asserts that the load-bearing content of a
trace-state-only closure is
formed-layer closure content, is named as an accepted recognition
source, and is not derivable from framework primitives.  It is recorded
as an obligation because this is the semantic thesis tying
recognition-mode landings to Foundations~I closure formation and the
\PaperCSL\ source naming, rather than a definitional property of a
transparent content object.  The source statement is the same-paper
Closure-Content Thesis.  The T9 opacity fix made
\LeanDecl{loadBearingFor} opaque so these conjuncts cannot be baked
into the selected content.

\paragraph{T10 Attack Foreclosure v5 refinement.}
The T10 v5 foreclosure axiom asserts the v5 refinement: derivation
routes remain subject to v4
Attack Foreclosure, while recognition-mode landings import their
load-bearing source as named recognition and land at theorem-grade under
the SB closure assumption.  It is recorded as an obligation because the
exhaustive recognition/derivation path distinction and the named-import
claim are paper-proposal-grade content, not derived from the Part~I
foreclosure theorem alone.  The source statement is the same-paper
Attack Foreclosure v5 schema.  The T10 opacity fix made the
derivation-route, recognition-landing, and imports-as-named-recognition
predicates opaque.

\paragraph{T16 framework-general candidate.}
The T16 pending axiom asserts the framework-general candidate: for every
Clay-class
trace-state-only substrate, the cascade framework-type differs from the
standard Clay target framework-type.  It is recorded as an obligation
because only the PvNP-specific mismatch is proved directly in T16; the
universal claim is gated by replication on at least one further
Clay-class track.  The source statement is the same-paper
Framework-Type Self-Characterization.  The T16 re-review narrowed the
axiom from a decidable PvNP mismatch to this genuinely
framework-general candidate.

For the four Part II obligations T7--T10, ``paper-proposal-grade''
means the following lifting condition.  The two verified instances,
PvNP and RH, are not yet a theorem-grade universal classification.  Universal lifting requires both a third cross-track
trace-state-only instance and a structural argument from
V-Differential column-membership to the \SigPattern{} verdict pattern.
Until those two inputs exist, the universal forms of T7--T10 remain
candidate meta-theorems, even though their two-instance evidence is part
of the source record.

\subsection{Reviewer-cycle log}\label{app:formalization:reviewer-log}

The mechanization axis closed only after the following reviewer-driven
cycles.  This log records the substantive changes, not every local edit.

\begin{enumerate}[label=\textbf{R\arabic*.}, leftmargin=*]
\item \textbf{T1 substance calibration.}  The first T1 review rejected
projection-packaging: the original carrier-dichotomy obligation had
packaged the full conclusion.  The fix added the explicit
accepted-imports soundness hypothesis, proved the non-\(X\)-equivalent
native-closure branch by classical case split, and narrowed the
same-paper obligation to the \(X\)-equivalent three-way classification.

\item \textbf{T1 target-indexing fix.}  A second T1 review found that
bridge-failure and terminal-matrix conditions had been made
carrier-intrinsic rather than target-relative.  The carrier fields for
bridge failure, bridge incompatibility, and matrix-family equivalence
were changed to target-indexed predicates, and the \CREBF{} and
\CRECTMT{} predicates were updated to apply the target \(X\).

\item \textbf{Axis-level opacity and consistency review.}  The
section-level mechanization review found several obligations that were
too easily instantiated or too broad.  The fix cycle introduced the
coverage-context guard for T1/T3, narrowed T5 with the framework-scope
guard and opaque internal-move semantics, added the template-applies
guard and opaque shared Part II predicates for T7, made T8's sweep
verdict opaque, made T9's load-bearing content opaque, and made T10's
derivation-route, recognition-landing, and named-recognition predicates
opaque.

\item \textbf{Axis re-review and theorem-scope tightening.}  The
follow-up review required three further structural fixes.  T2 acquired
the bridge-composition guard and was reclassified as partial because it
applies a same-paper bridge obligation.  T15 was strengthened from a
projection-style hole statement to a finite coverage geometry with
covered regions, uncovered regions, typed holes, non-monotone coverage
power, and the bounded-depth overlap lemma.  T16 was reclassified so
the PvNP-specific mismatch is proved directly, while the
framework-general candidate remains a separate guarded obligation not
used by the T16 proof.

\item \textbf{Final polish and prep-arc Phase-Z.}  The last
mechanization polish surfaced the T15 overlap claim explicitly and left
the Lean chain clean.  The later prep-arc Phase-Z review fixed drafting
artifacts rather than Lean content: stale section paths, provisional
bibkeys, appendix inputs in \LeanDecl{main.tex}, and count wording in
the prep documents.  The post-fix state opened the drafting arc with
the validator chain green.
\end{enumerate}

\subsection{Paper-proposal-grade lifting program}\label{app:formalization:lifting-program}

The recognition-mode obligations T7--T10 are intentionally separated
from theorem-grade universal status.  They are paper-proposal-grade
meta-theorems: the PvNP and RH instances supply two cross-track
validations, but two instances do not yet determine the universal
trace-state-only pattern.  The body of the paper records this
once in Part II; this appendix gives the exact lifting program to which
that body caveat refers.

The first future input is a third substrate landing in the
V-Differential trace-state-only column, beyond PvNP and RH.  Hodge is
the leading candidate, but the lifting requirement is structural rather
than nominal: the instance must exhibit a formed closure under the SB
closure assumption, must satisfy the seven-stage recognition-mode
template, and must import its load-bearing source as named recognition
rather than as framework-primitive derivation.  A
third instance that merely resembles the existing cases informally would
not discharge the lifting condition.

The second future input is a structural argument from V-Differential
trace-state-only column-membership to the Stage-6 \SigPattern{} verdict
pattern.  Section~\ref{sec:recognition:verdict-signature} sketches the
shape of that argument: framework-primitive derivation is ruled out by
the closure-content thesis, SAU non-descent elevates to the
V-Differential alternative, and V-Differential elevation is a structural
relabeling rather than a new derivation route.  What remains pending is
the theorem-grade proof that this reasoning follows from column
membership itself, uniformly across all trace-state-only substrates.

Until both inputs are supplied, T7--T10 remain candidate universal
meta-theorems.  Their present Lean status is therefore deliberately
obligation-status, guarded by opaque predicates, while the two-instance
PvNP/RH record remains valid as evidence rather than as a completed
universal proof.

\subsection{Validator runbook}\label{app:formalization:validator-runbook}

The following commands reproduce the formalization and paper-facing
validation chain used for this manuscript.

\begin{enumerate}[label=\arabic*., leftmargin=*]
\item From \LeanDecl{lean/}, run \LeanDecl{lake build}.  The expected
  meta-math close state runs 105 jobs and must complete cleanly, with
  no \LeanDecl{sorry}, \LeanDecl{admit}, or unregistered project axiom.

\item From the repository root, run
  \LeanDecl{python3 scripts/check_manifests.py}.  This checker verifies
  that theorem-tagged declarations have only allowlisted axioms in their
  transitive dependency closure, using
  \LeanDecl{lean/manifests/trust_base.txt} as the allowlist.

\item From the repository root, run
  \LeanDecl{python3 scripts/check_statements_of_record_meta_math.py}.
  This checker cross-checks
  \LeanDecl{paper/meta_math/notes/statements-of-record.yml} against the
  Lean manifest, the 17-row paper inventory, and the nine meta-math
  trust-base obligations.

\item From the repository root, run
  \LeanDecl{make paper-preflight-meta_math}.  This composed gate runs
  the PDF build, \LeanDecl{scripts/paper_lint.sh paper/meta_math}, and
  the statements-of-record validator.  It is the paper-facing close
  gate for drafting dispatches and final review.
\end{enumerate}

The pre-drafting gate state and command order are recorded in the
preflight-signoff note:
\begin{center}
\LeanDecl{paper/meta_math/notes/preflight-signoff.md}
\end{center}
 
\clearpage
\bibliographystyle{plainnat}
\bibliography{references}

\end{document}
